%\usepackage[utf8]{inputenc}
%\usepackage[utf8]{inputenc}
%\usepackage[utf8]{inputenc}
\usepackage[T1]{fontenc}
\usepackage[english]{babel}
\usepackage{amsfonts}
\usepackage{ifthen}
\usepackage{mdframed}
%\usepackage{amssymb}
%\usepackage[margin=0.5cm, paper width=100cm, paper height=100cm]{geometry}
\usepackage{amsmath}
\usepackage{amsthm}
\usepackage{verbatim}
\usepackage{array}
\usepackage{booktabs}
%\usepackage[sc]{mathpazo}
\usepackage{multirow}
\usepackage[table]{xcolor}
\usepackage{datetime}
\usepackage[thinlines]{easytable}

\renewcommand{\familydefault}{\sfdefault}
\renewcommand{\arraystretch}{.3} %% ez állítja a sorok távolságát!!!
\setlength{\heavyrulewidth}{1.5pt}
\setlength{\abovetopsep}{4pt}
\renewcommand{\qedsymbol}{$\blacksquare$}
\newcommand\numberthis{\addtocounter{equation}{1}\tag{\theequation}}
\setlength{\marginparwidth}{3.5cm}
\usepackage{graphicx}
\usepackage[square%round
, colon]{natbib}
\setcitestyle{aysep={},yysep={;}}
\usepackage{capt-of}
\usepackage{tikz}
\usepgflibrary{arrows}
\usetikzlibrary{shapes}
\usetikzlibrary{calc}
\usetikzlibrary{decorations}
\usetikzlibrary{decorations.pathmorphing}
\usetikzlibrary{decorations.shapes}
\usetikzlibrary{decorations.pathreplacing}
\usetikzlibrary{intersections}
\frenchspacing
\title{Órarend}
\author{Attila Molnár}

\newcommand{\oda}[5][]{\begin{tikzpicture}[remember picture,overlay]
\node[#1] at ([xshift=#3 cm, yshift=#4 cm]current page.#2){#5};
\end{tikzpicture}}

\newcommand{\osztalylista}{%
1/3/kny,
2/5/nye,
3/5/nyf,
4/4/9.a,
5/3/9.b,
6/2/9.d,
7/3/9.e,
8/3/9.f,
9/7/10.a,
10/3/10.b,
11/4/10.d,
12/6/10.e,
13/6/10.f,
14/8/11.a,
15/7/11.b,
16/7/11.c,
17/7/11.d,
18/3/11.e,
19/4/11.f,
20/3/12.a,
21/7/12.b,
22/7/12.c,
23/8/12.d,
24/3/12.e,
25/3/12.f%
}


\definecolor{te}{HTML}{FACEBC}  % Pompadour(r�zsasz�n)
\definecolor{gyte}{HTML}{FACEBC}  % Pompadour(r�zsasz�n)
\definecolor{mt}{HTML}{B09CD9}  % ringlo(lila)
\definecolor{msz}{HTML}{CCAEF1}  %  ak�c(lila)
\definecolor{mts}{HTML}{E3A6F8}  % orgona(lila)
\definecolor{of}{HTML}{C70081}  % BAZSAR�ZSAV�R�S (b�bor)
\definecolor{r}{HTML}{FF1900}  % skarl�t(V�R�S)
\definecolor{mte}{HTML}{E80606}  % ribizli(V�R�S)
\definecolor{{\' e}n}{HTML}{FBCB13}  % kukorica(s�rga)
\definecolor{{\' e}k}{HTML}{F5E748}  % kan�ri(s�rga) 
\definecolor{zt}{HTML}{F4CE69}  % barokk(s�rga)
\definecolor{et}{HTML}{EEEAE2}  % Ekr�
\definecolor{i}{HTML}{F0C76A} % MUSKOT�LY(s�rga)
\definecolor{ny}{HTML}{F5C533} % Monarchia(s�rga)
\definecolor{a}{HTML}{FDF74F} % HANZA(s�rga)
\definecolor{asz}{HTML}{FFF566} % k�n(s�rga)
\definecolor{ol}{HTML}{7EC476} % eozin(z�ld)
\definecolor{inf}{HTML}{31BDFF}  % l�zer(k�k)
\definecolor{ao}{HTML}{FFEC80}  % Mim�za
\definecolor{t{\" o}}{HTML}{CC483F}  % S�rk�nyv�r(v�r�s)
\definecolor{{\' e}v}{HTML}{B04A4A}  % Velencei(v�r�s)
\definecolor{civ}{HTML}{73A500}  % KLOROFIL(z�ld)
\definecolor{n{\' e}}{HTML}{FA493C}  % korall(piros)
\definecolor{mt{\" o}}{HTML}{719A0E}    % kadmium(z�ld)
\definecolor{f{\" o}}{HTML}{5FFF30}  % papg�j(z�ld)
\definecolor{k{\' e}}{HTML}{EA71EA}  % cikl�men(lila)
\definecolor{or}{HTML}{EEEAE2}  % selyem(feh�r)
\definecolor{fi}{HTML}{C68A6A}  % v�rhenyes(barna)
\definecolor{fr}{HTML}{99BCCC}  % ac�l(k�k)
\definecolor{bi}{HTML}{F96960}  % eper(piros)
\definecolor{tm}{HTML}{EEC532}  % �szarany
\definecolor{sp}{HTML}{FFA20C}  % s�fr�ny(s�rga)
\definecolor{m{\' e}}{HTML}{A5B5F7}  % kat�ngk�r�(k�k)
\definecolor{{\' a}g}{HTML}{C37ED2}  % perkin(lila)
\definecolor{fil}{HTML}{EDE2CA}  % GYAPJ�(feh�r)

\newcommand{\summag}{125} % ezt már a matlabból, vagy szemre beállítható. Kumulatív összeg, azért.
\newcommand{\tantargynev}[1]{\textbf{#1}}
\newenvironment{blokk}{\tabcolsep=.02cm \begin{tabular}{lccl}}{\end{tabular}}

\newcommand{\keretdraw}[1]{\ifthenelse{\equal{#1}{0}}{none}{black}}
\newcommand{\keretfill}[1]{\ifthenelse{\equal{#1}{0}}{none}{gray!50}}
\newcommand{\keretline}[1]{\ifthenelse{\equal{#1}{0}}{}{thick}}

\newcommand{\ora}[5]{
%1 rajzolja-e le
%2 az osztály sorszáma
%3 Nap
%4 Óra
%5 \oraresz{}{}{}{}, ami parancsban megtalálható a 
	% #1:tanár
	% #2:tantárgy
	% #3:csoport
	% #4:terem
	\ifthenelse{\equal{#1}{#2}}{%csak akkor álljon neki rajzolni, ha az első két argumentum megegyezik. Az első majd egy futó változó lesz egy forciklusban, így tanáronként fog csak plottolni, és nem lesz tex-capacity exceeded.
	% Cella közepe: blokk
	\pgfmathtruncatemacro{\osztsucc}{#2+1}
	\path (BFS#2-#3-#4)--(JFS\osztsucc-#3-#4)
		node[inner sep=1mm, midway, anchor=center, scale=1.1]{#5};
	}{}% end of \ifthenelse{\equal{#1}{#2}}
} % end of \newcommand{ora}

\newcommand{\szel}{5}
\newcommand{\oraresz}[5]%
	% #1:tanár
	% #2:tantárgy
	% #3:csoport
	% #4:terem
	% #5: cellavastagság
{\begin{tikzpicture}
\pgfmathsetmacro{\ymekkora}{#5*1.6	}
\pgfmathsetmacro{\xmekkora}{.8*\szel	}
\node[rectangle, draw = black, thick, fill = #2, rounded corners = 8pt, minimum height = \ymekkora em, minimum width = \xmekkora cm, scale = 1] at (0,0)
	{\ifthenelse{\equal{#3}{-}}{}{(#3) }
	#1 \textbf{#2} 
		\ifthenelse{\equal{#4}{-}}{}{{[}#4{]}}};
\end{tikzpicture}}

\newenvironment{cellatartalom}{\tabcolsep = 0mm\begin{tabular}{c}}{\end{tabular}}

\newcommand{\skeleton}{\begin{tikzpicture}[remember picture,overlay]
		\pgfmathsetmacro\opn{9}
		\pgfmathsetmacro\szelesseg{\szel} 
		\pgfmathsetmacro\evfolyamszelesseg{2}
		\pgfmathsetmacro\magassag{.75}
		\pgfmathsetmacro{\teremletszam}{60}
		\pgfmathsetmacro{\osztalycimkeszelesseg}{.6*\szelesseg}
		\pgfmathsetmacro{\tablaszelesseg}{(\opn+1)*5*\szelesseg}
		\pgfmathsetmacro\hezag{.1} 



% Az egész úgy működik, hogy állítható az elején az egyes osztályok sávjainak vastagsága.
% ez úgy fog megvalósulni, hogy minden vonal "az előző vonalhoz képest" lesz elhelyezve.
% tehát koordinátákat pakolunk le mindig, és a vonalakat ezek közé húzzuk, a koordináták elhelyezésekor pedig egymásra hivatkozunk. 
% a koordinátákat pedig ciklussal pakoljuk le.

%Ez a settings.

% Ez pakolja le a VÍZSZINTES VONALAKat és az osztályok neveit. 
% Mivel relatív megy minden, a nulladik pozíciójával lehet állítani a többit is.
% asszem azért evc, mert elválasztó vonal. A c-t nem tudom.
% Csigolyáknak nevezem a hétfői nulladik órák bal felső sarkait. Minden osztályra jut egy csigolya.
\coordinate (Csigolya0) at (0,0);
%\node[scale=3] at (Csigolya0){X};
\coordinate[xshift=-\osztalycimkeszelesseg cm] (evcStart0) at (Csigolya0);
\pgfmathsetmacro{\teljesszelesseg}{\tablaszelesseg +\osztalycimkeszelesseg-\szelesseg +5*\hezag*\szelesseg }
\coordinate[xshift= \teljesszelesseg cm] (evcEnd0) at (Csigolya0);
\draw[ultra thick] (evcStart0)--(evcEnd0);
	
\foreach \i/\v/\osztaly in \osztalylista % sorszám, vastagság, osztály
{\pgfmathtruncatemacro{\preci}{\i-1}
 \coordinate (Csigolya\i) at ([yshift=-\v*\magassag cm]Csigolya\preci);
 \coordinate (evcStart\i) at ([yshift=-\v*\magassag cm]evcStart\preci);
 \coordinate (evcEnd\i)   at ([yshift=-\v*\magassag cm]evcEnd\preci);
 \path (Csigolya\i)--(Csigolya\preci) 
  node[anchor=base, inner sep=2mm, midway, left]{\scalebox{3}{\sc\MakeUppercase{\osztaly}}};
 \path (evcEnd\i)--(evcEnd\preci) 
	node[anchor=base, inner sep=2mm, midway, right, xshift=-\osztalycimkeszelesseg cm]{\scalebox{3}{\sc\MakeUppercase{\osztaly}}};
}%end of \foreach \i/\v/\osztaly in \osztalylista


\foreach \i/\j/\cimke/\meret/\kijovetel/\isep in {
 0/3/Ny. évf./1/1/0 mm,
 3/9/9. évf./1/1/0 mm,
 9/14/10. évf./1/1/0 mm,
 14/19/11. évf./1/1/0 mm,
 19/25/12. évf./1/1/0 mm
 }
{%balra
\coordinate[xshift=-\kijovetel*\evfolyamszelesseg cm ] (SHIFTevcStart\j) at (evcStart\j);
 \pgfmathsetmacro{\kijovetelcimke}{\kijovetel-1}
 \coordinate[xshift=-\kijovetelcimke*\evfolyamszelesseg cm ] (SHIFT2evcStart\i) at (evcStart\i);
 \draw[draw = black, line width= 2mm, rounded corners=15mm]  (evcEnd\j) -- (SHIFTevcStart\j) |- (evcEnd\i) ;
 \path(SHIFT2evcStart\i) -- (SHIFTevcStart\j) node[midway, rotate=90, scale=3, anchor=center, fill=white, ellipse, inner sep =\isep]{\textsc{\scalebox{\meret}{\cimke}}};
% jobbra
\coordinate[xshift=\kijovetel*\evfolyamszelesseg cm ] (SHIFTevcEnd\j) at (evcEnd\j);
 \coordinate[xshift=\kijovetelcimke*\evfolyamszelesseg cm ] (SHIFT2evcEnd\i) at (evcEnd\i);
 \draw[draw = black, line width= 2mm, rounded corners=15mm]  (evcEnd\j) -- (SHIFTevcEnd\j) |- (evcEnd\i);
 \path(SHIFT2evcEnd\i) -- (SHIFTevcEnd\j) node[midway, rotate=90, scale=3, anchor=center, fill=white, ellipse, inner sep =\isep]{\textsc{\scalebox{\meret}{\cimke}}};
}

%Bal felső sarkok elnevezése
\foreach \osztaly in {1,...,26}
{	\pgfmathtruncatemacro{\predosztaly}{\osztaly-1}
			\foreach\nap in {1,...,5}
	{
	 \foreach\ora in {0,...,6}
		{	\pgfmathsetmacro{\eltolas}{((\nap-1)*(\opn+1+\hezag)+\ora)*\szelesseg}
			\coordinate[xshift=\eltolas cm] (BFS\osztaly-\nap-\ora) at (Csigolya\predosztaly);
			\coordinate[xshift=\szelesseg cm] (JFS\osztaly-\nap-\ora) at (BFS\osztaly-\nap-\ora);
		}%end of \foreach \orak
	 \pgfmathtruncatemacro{\meddig}{\opn}
	 \foreach\ora in {7,...,\meddig}
		{	\pgfmathsetmacro{\eltolashezaggal}{((\nap-1)*(\opn+1+\hezag)+\ora+\hezag)*\szelesseg}
			\coordinate[xshift=\eltolashezaggal cm] (BFS\osztaly-\nap-\ora) at (Csigolya\predosztaly);
			\coordinate[xshift=\szelesseg cm] (JFS\osztaly-\nap-\ora) at (BFS\osztaly-\nap-\ora);
		}%end of \foreach \orak
	}%end of \foreach \napok 
}%end of \foreach \osztaly

% FÜGGŐLEGES VONALAK. Két ciklus egymásban. Az egyik a napokat pakolja, a másik azon belül az órákat.
% plusz utána a lezárása.
\pgfmathtruncatemacro{\oszlopszam}{(\opn+1)*6*\szelesseg}
\foreach \napszam/\napnev in {1/h{\' e}tf{\H o}, 2/kedd, 3/szerda, 4/cs{\" u}t{\" o}rt{\" o}k, 5/p{\' e}ntek}
	{
% 	settings
		\pgfmathsetmacro{\honnan}{(\napszam-1)*((\opn+1+\hezag)*\szelesseg)}		
		\pgfmathsetmacro{\utolsoora}{\honnan+\opn*\szelesseg+\hezag*\szelesseg}
		\pgfmathsetmacro{\szammagassag}{1*\magassag}
		\pgfmathsetmacro{\kisvonalmagassag}{2*\magassag}
		\pgfmathsetmacro{\nagyvonalmagassag}{5*\magassag}
	
%		 napok körüli 3 db vonal a nap nevének feliratával
		 
		\draw[line width= 2mm, rounded corners = 15mm] 
			(\honnan, -\summag) 
			-- (\honnan, \nagyvonalmagassag)
			-- (\utolsoora, \nagyvonalmagassag)node[midway, below, inner sep = 5mm]{\scalebox{5}{\textsc{\napnev}}}
			-- (\utolsoora, -\summag);
			

%		nulladik órák szürke sávja
		\filldraw[fill opacity=.1, draw opacity=1, ultra thick,  rounded corners = 5mm] (\honnan+\szelesseg, \kisvonalmagassag) rectangle (\honnan, -\summag);

%	Függőleges vonalak a duplavonalig.
		\pgfmathtruncatemacro{\eddig}{\opn-2}
		\foreach \oraszam in {2, ..., \eddig} %% a nulladik óra utánit leteszi a nulladik órás szürke téglalap.
			{
				\draw[ultra thick] (\honnan+\oraszam*\szelesseg, \kisvonalmagassag) 
												-- (\honnan+\oraszam*\szelesseg, -\summag);
			}
% Óraszámok számai
		\pgfmathtruncatemacro{\eddig}{\opn-3}
		\foreach \oraszam in {0, ..., \eddig}
			{\node[anchor=center] at (\honnan+\oraszam*\szelesseg+.5*\szelesseg, \szammagassag){\scalebox{3.5}{\oraszam}};}

	 \pgfmathtruncatemacro{\ettol}{\opn-2}
	 \pgfmathtruncatemacro{\eddig}{\opn-1}

%	Függőleges vonalak a duplavonal utántól.
	 \foreach \oraszam in {\ettol, ..., \eddig}
		{
			\draw[ultra thick] (\honnan+\oraszam*\szelesseg+\hezag*\szelesseg, \kisvonalmagassag) 
											-- (\honnan+\oraszam*\szelesseg+\hezag*\szelesseg, -\summag);
		}
% Óraszámok számai a duplavonal utántól
	 \foreach \oraszam in {\ettol, ..., \eddig}
			{
				\node[anchor=center] at 
					(\honnan+\oraszam*\szelesseg+.5*\szelesseg+\hezag*\szelesseg, \szammagassag)
					{\scalebox{3.5}{\oraszam}};
			}
}
	
% Vízszintes vonalak
\foreach \i/\v/\osztaly in \osztalylista {\draw[ultra thick] (evcStart\i)--(evcEnd\i);}

% Napok közti osztálynevek
\foreach \i/\v/\osztaly in \osztalylista 
{
	\foreach \j in {1,2,3,4}
	{
		\pgfmathtruncatemacro{\isucc}{\i+1}
		\pgfmathtruncatemacro{\jsucc}{\j+1}
		\path (BFS\i-\j-\opn)--(BFS\isucc-\jsucc-0) 
			node[inner sep=1mm, anchor=center, midway]{\scalebox{3}{\sc \MakeUppercase{\osztaly}}};
	}
}

\end{tikzpicture}
}%end of skeleton

\newcommand{\orarend}[1]{% ott tartottam, hogy remember overlaybe ki kell rakni a skeletont, és a dorflesht egy 1...25-ös ciklusban kell rem-overlayekkel lepakolni.
% #1: mettől meddig nyomtasson

	\foreach \mitrakjonle in {#1} % Egyedül a DORFLESH.tex-ben van használva ez a változó.
	{
	\begin{tikzpicture}[remember picture, overlay]
		\ora{\mitrakjonle}{1}{1}{0}
  {\begin{cellatartalom}%
    %
  \end{cellatartalom}}
\ora{\mitrakjonle}{1}{1}{1}
  {\begin{cellatartalom}%
     \oraresz{Both}{i}{2}{I.33}{4}%%
  \end{cellatartalom}}
\ora{\mitrakjonle}{1}{1}{2}
  {\begin{cellatartalom}%
     \oraresz{Cserv{\' a}k}{inf}{1}{Szt.1}{2}%
\\ \oraresz{Both}{ny}{2}{I.33}{2}%%
  \end{cellatartalom}}
\ora{\mitrakjonle}{1}{1}{3}
  {\begin{cellatartalom}%
     \oraresz{Barth{\' a}n{\' e}}{{\' e}n}{-}{{\' E}.}{4}%%
  \end{cellatartalom}}
\ora{\mitrakjonle}{1}{1}{4}
  {\begin{cellatartalom}%
     \oraresz{Gismondi}{ol}{ol1}{I.39}{1.33333333333333}%
\\ \oraresz{Krasznay}{ol}{ol3}{Ol.1}{1.33333333333333}%
\\ \oraresz{Bellesi}{ol}{ol2}{Ol.2}{1.33333333333333}%%
  \end{cellatartalom}}
\ora{\mitrakjonle}{1}{1}{5}
  {\begin{cellatartalom}%
     \oraresz{Csuka}{ol}{ol3}{Ol.2}{1.33333333333333}%
\\ \oraresz{Somk{\" o}vi}{ol}{ol2}{F.10}{1.33333333333333}%
\\ \oraresz{Fajcs{\' a}k}{ol}{ol1}{Ol.1}{1.33333333333333}%%
  \end{cellatartalom}}
\ora{\mitrakjonle}{1}{1}{6}
  {\begin{cellatartalom}%
     \oraresz{Csuka}{ol}{ol3}{Ol.2}{1.33333333333333}%
\\ \oraresz{Somk{\" o}vi}{ol}{ol2}{F.10}{1.33333333333333}%
\\ \oraresz{Fajcs{\' a}k}{ol}{ol1}{Ol.1}{1.33333333333333}%%
  \end{cellatartalom}}
\ora{\mitrakjonle}{1}{1}{7}
  {\begin{cellatartalom}%
     \oraresz{Both}{i}{1}{I.33}{2}%
\\ \oraresz{F{\" o}ldesi}{inf}{2}{Szt.1}{2}%%
  \end{cellatartalom}}
\ora{\mitrakjonle}{1}{1}{8}
  {\begin{cellatartalom}%
    %
  \end{cellatartalom}}
\ora{\mitrakjonle}{1}{1}{9}
  {\begin{cellatartalom}%
    %
  \end{cellatartalom}}
\ora{\mitrakjonle}{1}{2}{0}
  {\begin{cellatartalom}%
     \oraresz{{\' A}goston}{mt}{1}{I.34}{4}%%
  \end{cellatartalom}}
\ora{\mitrakjonle}{1}{2}{1}
  {\begin{cellatartalom}%
     \oraresz{F{\" o}ldesi}{mt}{2}{M.}{2}%
\\ \oraresz{Cserv{\' a}k}{inf}{1}{Szt.1}{2}%%
  \end{cellatartalom}}
\ora{\mitrakjonle}{1}{2}{2}
  {\begin{cellatartalom}%
     \oraresz{Csuka}{ol}{ol2}{I.33}{1.33333333333333}%
\\ \oraresz{Fajcs{\' a}k}{ol}{ol1}{II.41}{1.33333333333333}%
\\ \oraresz{Bellesi}{ol}{ol3}{Kt.}{1.33333333333333}%%
  \end{cellatartalom}}
\ora{\mitrakjonle}{1}{2}{3}
  {\begin{cellatartalom}%
     \oraresz{Lad{\' a}nyi}{te}{L}{Tt.}{1.33333333333333}%
\\ \oraresz{Bard{\' o}czky}{te}{L}{Tt.}{1.33333333333333}%
\\ \oraresz{Heged{\H u}s}{te}{F}{Tt.}{1.33333333333333}%%
  \end{cellatartalom}}
\ora{\mitrakjonle}{1}{2}{4}
  {\begin{cellatartalom}%
     \oraresz{Gismondi}{ol}{ol1}{I.39}{1.33333333333333}%
\\ \oraresz{Heged{\" u}s}{ol}{ol3}{Ol.1}{1.33333333333333}%
\\ \oraresz{Somk{\" o}vi}{ol}{ol2}{I.33}{1.33333333333333}%%
  \end{cellatartalom}}
\ora{\mitrakjonle}{1}{2}{5}
  {\begin{cellatartalom}%
     \oraresz{Heged{\" u}s}{ol}{ol3}{Ol.1}{1.33333333333333}%
\\ \oraresz{Somk{\" o}vi}{ol}{ol2}{I.33}{1.33333333333333}%
\\ \oraresz{Gianelli}{ol}{ol1}{Ol.2}{1.33333333333333}%%
  \end{cellatartalom}}
\ora{\mitrakjonle}{1}{2}{6}
  {\begin{cellatartalom}%
     \oraresz{Heged{\" u}s}{ol}{ol3}{Ol.1}{1.33333333333333}%
\\ \oraresz{Somk{\" o}vi}{ol}{ol2}{L.}{1.33333333333333}%
\\ \oraresz{Gianelli}{ol}{ol1}{Ol.2}{1.33333333333333}%%
  \end{cellatartalom}}
\ora{\mitrakjonle}{1}{2}{7}
  {\begin{cellatartalom}%
    %
  \end{cellatartalom}}
\ora{\mitrakjonle}{1}{2}{8}
  {\begin{cellatartalom}%
    %
  \end{cellatartalom}}
\ora{\mitrakjonle}{1}{2}{9}
  {\begin{cellatartalom}%
    %
  \end{cellatartalom}}
\ora{\mitrakjonle}{1}{3}{0}
  {\begin{cellatartalom}%
    %
  \end{cellatartalom}}
\ora{\mitrakjonle}{1}{3}{1}
  {\begin{cellatartalom}%
     \oraresz{Heged{\" u}s}{ol}{ol3}{I.33}{1.33333333333333}%
\\ \oraresz{Csuka}{ol}{ol2}{Ol.1}{1.33333333333333}%
\\ \oraresz{Fajcs{\' a}k}{ol}{ol1}{F.11}{1.33333333333333}%%
  \end{cellatartalom}}
\ora{\mitrakjonle}{1}{3}{2}
  {\begin{cellatartalom}%
     \oraresz{Heged{\" u}s}{ol}{ol3}{I.33}{1.33333333333333}%
\\ \oraresz{Csuka}{ol}{ol2}{Ol.1}{1.33333333333333}%
\\ \oraresz{Fajcs{\' a}k}{ol}{ol1}{Ol.2}{1.33333333333333}%%
  \end{cellatartalom}}
\ora{\mitrakjonle}{1}{3}{3}
  {\begin{cellatartalom}%
     \oraresz{Heged{\" u}s}{ol}{ol3}{I.33}{1}%
\\ \oraresz{Somk{\" o}vi}{ol}{ol2}{Ol.1}{1}%
\\ \oraresz{Gianelli}{ol}{ol1}{R.}{1}%
\\ \oraresz{Sum}{mid}{mid}{FS.t}{1}%%
  \end{cellatartalom}}
\ora{\mitrakjonle}{1}{3}{4}
  {\begin{cellatartalom}%
     \oraresz{Csuka}{ol}{ol3}{I.33}{1.33333333333333}%
\\ \oraresz{Somk{\" o}vi}{ol}{ol2}{Ol.1}{1.33333333333333}%
\\ \oraresz{Gianelli}{ol}{ol1}{R.}{1.33333333333333}%%
  \end{cellatartalom}}
\ora{\mitrakjonle}{1}{3}{5}
  {\begin{cellatartalom}%
     \oraresz{Csuka}{ol}{ol1}{Ol.2}{1.33333333333333}%
\\ \oraresz{Somk{\" o}vi}{ol}{ol3}{I.30}{1.33333333333333}%
\\ \oraresz{Gianelli}{ol}{ol2}{II.40}{1.33333333333333}%%
  \end{cellatartalom}}
\ora{\mitrakjonle}{1}{3}{6}
  {\begin{cellatartalom}%
    %
  \end{cellatartalom}}
\ora{\mitrakjonle}{1}{3}{7}
  {\begin{cellatartalom}%
    %
  \end{cellatartalom}}
\ora{\mitrakjonle}{1}{3}{8}
  {\begin{cellatartalom}%
    %
  \end{cellatartalom}}
\ora{\mitrakjonle}{1}{3}{9}
  {\begin{cellatartalom}%
    %
  \end{cellatartalom}}
\ora{\mitrakjonle}{1}{4}{0}
  {\begin{cellatartalom}%
    %
  \end{cellatartalom}}
\ora{\mitrakjonle}{1}{4}{1}
  {\begin{cellatartalom}%
     \oraresz{Gianelli}{ol}{ol1}{L.}{1.33333333333333}%
\\ \oraresz{Csuka}{ol}{ol3}{I.33}{1.33333333333333}%
\\ \oraresz{Somk{\" o}vi}{ol}{ol2}{Ol.2}{1.33333333333333}%%
  \end{cellatartalom}}
\ora{\mitrakjonle}{1}{4}{2}
  {\begin{cellatartalom}%
     \oraresz{Somk{\" o}vi}{ol}{ol2}{Ol.2}{1.33333333333333}%
\\ \oraresz{Csuka}{ol}{ol1}{I.33}{1.33333333333333}%
\\ \oraresz{Bellesi}{ol}{ol3}{F.18}{1.33333333333333}%%
  \end{cellatartalom}}
\ora{\mitrakjonle}{1}{4}{3}
  {\begin{cellatartalom}%
     \oraresz{Somk{\" o}vi}{ol}{ol1}{Ol.2}{1.33333333333333}%
\\ \oraresz{Heged{\" u}s}{ol}{ol3}{I.33}{1.33333333333333}%
\\ \oraresz{Csuka}{ol}{ol2}{L.}{1.33333333333333}%%
  \end{cellatartalom}}
\ora{\mitrakjonle}{1}{4}{4}
  {\begin{cellatartalom}%
     \oraresz{Heged{\" u}s}{ol}{ol3}{I.33}{1}%
\\ \oraresz{Csuka}{ol}{ol2}{L.}{1}%
\\ \oraresz{Fajcs{\' a}k}{ol}{ol1}{Ol.1}{1}%
\\ \oraresz{Sum}{mid}{mid}{FS.t}{1}%%
  \end{cellatartalom}}
\ora{\mitrakjonle}{1}{4}{5}
  {\begin{cellatartalom}%
     \oraresz{Both}{i}{1}{I.33}{2}%
\\ \oraresz{F{\" o}ldesi}{inf}{2}{Szt.2}{2}%%
  \end{cellatartalom}}
\ora{\mitrakjonle}{1}{4}{6}
  {\begin{cellatartalom}%
     \oraresz{Both}{ny}{1}{I.33}{2}%
\\ \oraresz{F{\" o}ldesi}{mt}{2}{M.}{2}%%
  \end{cellatartalom}}
\ora{\mitrakjonle}{1}{4}{7}
  {\begin{cellatartalom}%
     \oraresz{Both}{i}{2}{I.33}{2}%
\\ \oraresz{{\' A}goston}{mt}{1}{M.}{2}%%
  \end{cellatartalom}}
\ora{\mitrakjonle}{1}{4}{8}
  {\begin{cellatartalom}%
    %
  \end{cellatartalom}}
\ora{\mitrakjonle}{1}{4}{9}
  {\begin{cellatartalom}%
    %
  \end{cellatartalom}}
\ora{\mitrakjonle}{1}{5}{0}
  {\begin{cellatartalom}%
    %
  \end{cellatartalom}}
\ora{\mitrakjonle}{1}{5}{1}
  {\begin{cellatartalom}%
     \oraresz{P{\' o}k}{r}{-}{R.}{4}%%
  \end{cellatartalom}}
\ora{\mitrakjonle}{1}{5}{2}
  {\begin{cellatartalom}%
     \oraresz{Csuka}{of}{-}{I.33}{4}%%
  \end{cellatartalom}}
\ora{\mitrakjonle}{1}{5}{3}
  {\begin{cellatartalom}%
     \oraresz{Csuka}{ol}{ol3}{Ol.1}{1.33333333333333}%
\\ \oraresz{Fajcs{\' a}k}{ol}{ol1}{I.33}{1.33333333333333}%
\\ \oraresz{Gianelli}{ol}{ol2}{Ol.2}{1.33333333333333}%%
  \end{cellatartalom}}
\ora{\mitrakjonle}{1}{5}{4}
  {\begin{cellatartalom}%
     \oraresz{Bellesi}{ol}{ol2}{Ol.2}{1.33333333333333}%
\\ \oraresz{Fajcs{\' a}k}{ol}{ol1}{Ol.1}{1.33333333333333}%
\\ \oraresz{Krasznay}{ol}{ol3}{I.33}{1.33333333333333}%%
  \end{cellatartalom}}
\ora{\mitrakjonle}{1}{5}{5}
  {\begin{cellatartalom}%
     \oraresz{F{\" o}ldesi}{inf}{2}{Szt.2}{2}%
\\ \oraresz{Cserv{\' a}k}{inf}{1}{Szt.1}{2}%%
  \end{cellatartalom}}
\ora{\mitrakjonle}{1}{5}{6}
  {\begin{cellatartalom}%
     \oraresz{Lad{\' a}nyi}{te}{L}{Tt.}{1.33333333333333}%
\\ \oraresz{Bard{\' o}czky}{te}{L}{Tt.}{1.33333333333333}%
\\ \oraresz{Heged{\H u}s}{te}{F}{Tt.}{1.33333333333333}%%
  \end{cellatartalom}}
\ora{\mitrakjonle}{1}{5}{7}
  {\begin{cellatartalom}%
     \oraresz{Lad{\' a}nyi}{te}{L}{Tt.}{1.33333333333333}%
\\ \oraresz{Bard{\' o}czky}{te}{L}{Tt.}{1.33333333333333}%
\\ \oraresz{Heged{\H u}s}{te}{F}{Tt.}{1.33333333333333}%%
  \end{cellatartalom}}
\ora{\mitrakjonle}{1}{5}{8}
  {\begin{cellatartalom}%
    %
  \end{cellatartalom}}
\ora{\mitrakjonle}{1}{5}{9}
  {\begin{cellatartalom}%
    %
  \end{cellatartalom}}
\ora{\mitrakjonle}{2}{1}{0}
  {\begin{cellatartalom}%
    %
  \end{cellatartalom}}
\ora{\mitrakjonle}{2}{1}{1}
  {\begin{cellatartalom}%
     \oraresz{T{\' o}th}{a}{2}{II.41}{5}%%
  \end{cellatartalom}}
\ora{\mitrakjonle}{2}{1}{2}
  {\begin{cellatartalom}%
     \oraresz{Horv{\' a}th G.}{of}{-}{II.41}{5}%%
  \end{cellatartalom}}
\ora{\mitrakjonle}{2}{1}{3}
  {\begin{cellatartalom}%
     \oraresz{S{\' o}tin{\' e}}{n{\' e}}{n{\' e}2}{N.}{1}%
\\ \oraresz{Hal{\' a}sz}{n{\' e}}{n{\' e}1}{R.}{1}%
\\ \oraresz{Szabolcsi}{fr}{fr}{Fr.}{1}%
\\ \oraresz{Gianelli}{ol}{ol}{Ol.2}{1}%
\\ \oraresz{Sum}{sp}{sp}{Sp.}{1}%%
  \end{cellatartalom}}
\ora{\mitrakjonle}{2}{1}{4}
  {\begin{cellatartalom}%
     \oraresz{Lad{\' a}nyi}{te}{L}{Tt.}{2.5}%
\\ \oraresz{De{\' a}k}{te}{F}{Tt.}{2.5}%%
  \end{cellatartalom}}
\ora{\mitrakjonle}{2}{1}{5}
  {\begin{cellatartalom}%
     \oraresz{T{\' o}th}{a}{an4}{A.3}{1.25}%
\\ \oraresz{Varga Bo.}{a}{an1}{I.30}{1.25}%
\\ \oraresz{Seprenyi}{a}{an3}{II.41}{1.25}%
\\ \oraresz{Hajba}{a}{an2}{Szt.3}{1.25}%%
  \end{cellatartalom}}
\ora{\mitrakjonle}{2}{1}{6}
  {\begin{cellatartalom}%
     \oraresz{T{\' o}th}{a}{an4}{A.3}{1.25}%
\\ \oraresz{Varga Bo.}{a}{an1}{I.30}{1.25}%
\\ \oraresz{Bert{\' o}k}{a}{an2}{Szt.3}{1.25}%
\\ \oraresz{Seprenyi}{a}{an3}{II.41}{1.25}%%
  \end{cellatartalom}}
\ora{\mitrakjonle}{2}{1}{7}
  {\begin{cellatartalom}%
    %
  \end{cellatartalom}}
\ora{\mitrakjonle}{2}{1}{8}
  {\begin{cellatartalom}%
    %
  \end{cellatartalom}}
\ora{\mitrakjonle}{2}{1}{9}
  {\begin{cellatartalom}%
    %
  \end{cellatartalom}}
\ora{\mitrakjonle}{2}{2}{0}
  {\begin{cellatartalom}%
     \oraresz{Fekete G.}{mt}{2}{II.41}{2.5}%
\\ \oraresz{Horv{\' a}th G.}{bi}{1}{B.ea}{2.5}%%
  \end{cellatartalom}}
\ora{\mitrakjonle}{2}{2}{1}
  {\begin{cellatartalom}%
     \oraresz{Sum}{sp}{sp}{Sp.}{1}%
\\ \oraresz{Hal{\' a}sz}{n{\' e}}{n{\' e}1}{N.}{1}%
\\ \oraresz{Szabolcsi}{fr}{fr}{Fr.}{1}%
\\ \oraresz{Gianelli}{ol}{ol}{Ol.1}{1}%
\\ \oraresz{S{\' o}tin{\' e}}{n{\' e}}{n{\' e}2}{A.1}{1}%%
  \end{cellatartalom}}
\ora{\mitrakjonle}{2}{2}{2}
  {\begin{cellatartalom}%
     \oraresz{Sum}{sp}{sp}{Sp.}{1}%
\\ \oraresz{Hal{\' a}sz}{n{\' e}}{n{\' e}1}{N.}{1}%
\\ \oraresz{Szabolcsi}{fr}{fr}{Fr.}{1}%
\\ \oraresz{Gianelli}{ol}{ol}{Ol.1}{1}%
\\ \oraresz{S{\' o}tin{\' e}}{n{\' e}}{n{\' e}2}{A.1}{1}%%
  \end{cellatartalom}}
\ora{\mitrakjonle}{2}{2}{3}
  {\begin{cellatartalom}%
     \oraresz{S{\" u}le}{i}{2}{II.41}{2.5}%
\\ \oraresz{Cserv{\' a}k}{inf}{1}{Szt.1}{2.5}%%
  \end{cellatartalom}}
\ora{\mitrakjonle}{2}{2}{4}
  {\begin{cellatartalom}%
     \oraresz{Sebesty{\' e}n}{mt}{1}{M.}{2.5}%
\\ \oraresz{Cserv{\' a}k}{inf}{2}{Szt.1}{2.5}%%
  \end{cellatartalom}}
\ora{\mitrakjonle}{2}{2}{5}
  {\begin{cellatartalom}%
    %
  \end{cellatartalom}}
\ora{\mitrakjonle}{2}{2}{6}
  {\begin{cellatartalom}%
     \oraresz{Juh{\' a}sz A.}{a}{an4}{A.3}{1.25}%
\\ \oraresz{Varga Bo.}{a}{an1}{N.}{1.25}%
\\ \oraresz{Bert{\' o}k}{a}{an2}{A.2}{1.25}%
\\ \oraresz{Seprenyi}{a}{an3}{II.41}{1.25}%%
  \end{cellatartalom}}
\ora{\mitrakjonle}{2}{2}{7}
  {\begin{cellatartalom}%
     \oraresz{Juh{\' a}sz A.}{a}{an4}{A.3}{1.25}%
\\ \oraresz{Varga Bo.}{a}{an1}{N.}{1.25}%
\\ \oraresz{Bert{\' o}k}{a}{an2}{A.2}{1.25}%
\\ \oraresz{Seprenyi}{a}{an3}{II.41}{1.25}%%
  \end{cellatartalom}}
\ora{\mitrakjonle}{2}{2}{8}
  {\begin{cellatartalom}%
     \oraresz{Krajny{\' a}k}{inf}{szak}{Szt.1}{5}%%
  \end{cellatartalom}}
\ora{\mitrakjonle}{2}{2}{9}
  {\begin{cellatartalom}%
    %
  \end{cellatartalom}}
\ora{\mitrakjonle}{2}{3}{0}
  {\begin{cellatartalom}%
     \oraresz{S{\' o}tin{\' e}}{n{\' e}}{n{\' e}2}{N.}{1}%
\\ \oraresz{Hal{\' a}sz}{n{\' e}}{n{\' e}1}{{\' E}.}{1}%
\\ \oraresz{Szabolcsi}{fr}{fr}{Fr.}{1}%
\\ \oraresz{Gianelli}{ol}{ol}{Ol.1}{1}%
\\ \oraresz{Sum}{sp}{sp}{Sp.}{1}%%
  \end{cellatartalom}}
\ora{\mitrakjonle}{2}{3}{1}
  {\begin{cellatartalom}%
     \oraresz{S{\" u}le}{i}{1}{II.41}{2.5}%
\\ \oraresz{Horv{\' a}th G.}{bi}{2}{B.ea}{2.5}%%
  \end{cellatartalom}}
\ora{\mitrakjonle}{2}{3}{2}
  {\begin{cellatartalom}%
     \oraresz{Hajba}{a}{an4}{Szt.3}{1.25}%
\\ \oraresz{Bert{\' o}k}{a}{an2}{A.1}{1.25}%
\\ \oraresz{Varga Bo.}{a}{an1}{I.30}{1.25}%
\\ \oraresz{Seprenyi}{a}{an3}{A.2}{1.25}%%
  \end{cellatartalom}}
\ora{\mitrakjonle}{2}{3}{3}
  {\begin{cellatartalom}%
     \oraresz{T{\' o}th}{a}{an4}{F.18}{1.25}%
\\ \oraresz{Varga Bo.}{a}{an1}{A.1}{1.25}%
\\ \oraresz{Husz{\' a}r}{a}{an2}{I.30}{1.25}%
\\ \oraresz{Seprenyi}{a}{an3}{A.2}{1.25}%%
  \end{cellatartalom}}
\ora{\mitrakjonle}{2}{3}{4}
  {\begin{cellatartalom}%
     \oraresz{T{\' o}th}{a}{an4}{F.18}{1.25}%
\\ \oraresz{Varga Bo.}{a}{an1}{A.1}{1.25}%
\\ \oraresz{Husz{\' a}r}{a}{an2}{I.30}{1.25}%
\\ \oraresz{Seprenyi}{a}{an3}{A.2}{1.25}%%
  \end{cellatartalom}}
\ora{\mitrakjonle}{2}{3}{5}
  {\begin{cellatartalom}%
    %
  \end{cellatartalom}}
\ora{\mitrakjonle}{2}{3}{6}
  {\begin{cellatartalom}%
     \oraresz{Lad{\' a}nyi}{te}{L}{Tt.}{2.5}%
\\ \oraresz{De{\' a}k}{te}{F}{Tt.}{2.5}%%
  \end{cellatartalom}}
\ora{\mitrakjonle}{2}{3}{7}
  {\begin{cellatartalom}%
     \oraresz{Lad{\' a}nyi}{te}{L}{Tt.}{2.5}%
\\ \oraresz{De{\' a}k}{te}{F}{Tt.}{2.5}%%
  \end{cellatartalom}}
\ora{\mitrakjonle}{2}{3}{8}
  {\begin{cellatartalom}%
    %
  \end{cellatartalom}}
\ora{\mitrakjonle}{2}{3}{9}
  {\begin{cellatartalom}%
    %
  \end{cellatartalom}}
\ora{\mitrakjonle}{2}{4}{0}
  {\begin{cellatartalom}%
     \oraresz{T{\' o}th}{a}{an4}{I.30}{1.25}%
\\ \oraresz{Varga Bo.}{a}{an1}{F.8}{1.25}%
\\ \oraresz{Husz{\' a}r}{a}{an2}{A.2}{1.25}%
\\ \oraresz{Seprenyi}{a}{an3}{II.41}{1.25}%%
  \end{cellatartalom}}
\ora{\mitrakjonle}{2}{4}{1}
  {\begin{cellatartalom}%
     \oraresz{T{\' o}th}{a}{an4}{I.30}{1.25}%
\\ \oraresz{Varga Bo.}{a}{an1}{F.8}{1.25}%
\\ \oraresz{Husz{\' a}r}{a}{an2}{A.2}{1.25}%
\\ \oraresz{Seprenyi}{a}{an3}{II.41}{1.25}%%
  \end{cellatartalom}}
\ora{\mitrakjonle}{2}{4}{2}
  {\begin{cellatartalom}%
     \oraresz{S{\' o}tin{\' e}}{n{\' e}}{n{\' e}2}{II.51}{1}%
\\ \oraresz{Hal{\' a}sz}{n{\' e}}{n{\' e}1}{Szt.3}{1}%
\\ \oraresz{Szabolcsi}{fr}{fr}{II.42}{1}%
\\ \oraresz{Gianelli}{ol}{ol}{F.9}{1}%
\\ \oraresz{Sum}{sp}{sp}{Kt.}{1}%%
  \end{cellatartalom}}
\ora{\mitrakjonle}{2}{4}{3}
  {\begin{cellatartalom}%
     \oraresz{T{\' o}th}{a}{1}{II.41}{2.5}%
\\ \oraresz{Cserv{\' a}k}{inf}{2}{Szt.3}{2.5}%%
  \end{cellatartalom}}
\ora{\mitrakjonle}{2}{4}{4}
  {\begin{cellatartalom}%
     \oraresz{Fekete G.}{mt}{2}{II.41}{2.5}%
\\ \oraresz{Cserv{\' a}k}{inf}{1}{Szt.1}{2.5}%%
  \end{cellatartalom}}
\ora{\mitrakjonle}{2}{4}{5}
  {\begin{cellatartalom}%
     \oraresz{S{\" u}le}{ny}{2}{II.41}{2.5}%
\\ \oraresz{S{\' o}skuthyn{\' e}}{k{\' e}}{1}{K.gy}{2.5}%%
  \end{cellatartalom}}
\ora{\mitrakjonle}{2}{4}{6}
  {\begin{cellatartalom}%
     \oraresz{Horv{\' a}th G.}{bi}{2}{B.ea}{2.5}%
\\ \oraresz{S{\' o}skuthyn{\' e}}{k{\' e}}{1}{K.gy}{2.5}%%
  \end{cellatartalom}}
\ora{\mitrakjonle}{2}{4}{7}
  {\begin{cellatartalom}%
    %
  \end{cellatartalom}}
\ora{\mitrakjonle}{2}{4}{8}
  {\begin{cellatartalom}%
     \oraresz{B{\" o}g{\" o}s}{a}{fz}{A.1}{5}%%
  \end{cellatartalom}}
\ora{\mitrakjonle}{2}{4}{9}
  {\begin{cellatartalom}%
    %
  \end{cellatartalom}}
\ora{\mitrakjonle}{2}{5}{0}
  {\begin{cellatartalom}%
     \oraresz{Cserv{\' a}k}{inf}{1}{Szt.1}{5}%%
  \end{cellatartalom}}
\ora{\mitrakjonle}{2}{5}{1}
  {\begin{cellatartalom}%
     \oraresz{Horv{\' a}th G.}{bi}{1}{B.gy}{2.5}%
\\ \oraresz{Cserv{\' a}k}{inf}{2}{Szt.1}{2.5}%%
  \end{cellatartalom}}
\ora{\mitrakjonle}{2}{5}{2}
  {\begin{cellatartalom}%
     \oraresz{T{\' o}th}{a}{an4}{R.}{1.25}%
\\ \oraresz{Varga Bo.}{a}{an1}{I.30}{1.25}%
\\ \oraresz{Husz{\' a}r}{a}{an2}{I.32}{1.25}%
\\ \oraresz{Seprenyi}{a}{an3}{Szt.3}{1.25}%%
  \end{cellatartalom}}
\ora{\mitrakjonle}{2}{5}{3}
  {\begin{cellatartalom}%
     \oraresz{T{\' o}th}{a}{an4}{R.}{1.25}%
\\ \oraresz{Varga Bo.}{a}{an1}{I.30}{1.25}%
\\ \oraresz{Husz{\' a}r}{a}{an2}{I.32}{1.25}%
\\ \oraresz{Seprenyi}{a}{an3}{Szt.3}{1.25}%%
  \end{cellatartalom}}
\ora{\mitrakjonle}{2}{5}{4}
  {\begin{cellatartalom}%
     \oraresz{S{\" u}le}{ny}{1}{II.41}{2.5}%
\\ \oraresz{S{\' o}skuthyn{\' e}}{k{\' e}}{2}{K.gy}{2.5}%%
  \end{cellatartalom}}
\ora{\mitrakjonle}{2}{5}{5}
  {\begin{cellatartalom}%
     \oraresz{S{\' o}skuthyn{\' e}}{k{\' e}}{2}{K.gy}{2.5}%
\\ \oraresz{Sebesty{\' e}n}{mt}{1}{II.41}{2.5}%%
  \end{cellatartalom}}
\ora{\mitrakjonle}{2}{5}{6}
  {\begin{cellatartalom}%
     \oraresz{S{\' o}tin{\' e}}{n{\' e}}{n{\' e}2}{N.}{1}%
\\ \oraresz{Hal{\' a}sz}{n{\' e}}{n{\' e}1}{{\' E}.}{1}%
\\ \oraresz{Szabolcsi}{fr}{fr}{Fr.}{1}%
\\ \oraresz{Gianelli}{ol}{ol}{L.}{1}%
\\ \oraresz{Sum}{sp}{sp}{Sp.}{1}%%
  \end{cellatartalom}}
\ora{\mitrakjonle}{2}{5}{7}
  {\begin{cellatartalom}%
    %
  \end{cellatartalom}}
\ora{\mitrakjonle}{2}{5}{8}
  {\begin{cellatartalom}%
    %
  \end{cellatartalom}}
\ora{\mitrakjonle}{2}{5}{9}
  {\begin{cellatartalom}%
    %
  \end{cellatartalom}}
\ora{\mitrakjonle}{3}{1}{0}
  {\begin{cellatartalom}%
     \oraresz{Darab{\' a}nt}{mt}{szt}{I.30}{5}%%
  \end{cellatartalom}}
\ora{\mitrakjonle}{3}{1}{1}
  {\begin{cellatartalom}%
     \oraresz{Szabados}{inf}{tk}{Tk.}{2.5}%
\\ \oraresz{Krajny{\' a}k}{inf}{szt}{Szt.2}{2.5}%%
  \end{cellatartalom}}
\ora{\mitrakjonle}{3}{1}{2}
  {\begin{cellatartalom}%
     \oraresz{Krajny{\' a}k}{inf}{szt}{Szt.2}{2.5}%
\\ \oraresz{P{\' o}k}{inf}{tk}{Tk.}{2.5}%%
  \end{cellatartalom}}
\ora{\mitrakjonle}{3}{1}{3}
  {\begin{cellatartalom}%
     \oraresz{S{\' o}tin{\' e}}{n{\' e}}{n{\' e}2}{N.}{1}%
\\ \oraresz{Hal{\' a}sz}{n{\' e}}{n{\' e}1}{R.}{1}%
\\ \oraresz{Szabolcsi}{fr}{fr}{Fr.}{1}%
\\ \oraresz{Gianelli}{ol}{ol}{Ol.2}{1}%
\\ \oraresz{Sum}{sp}{sp}{Sp.}{1}%%
  \end{cellatartalom}}
\ora{\mitrakjonle}{3}{1}{4}
  {\begin{cellatartalom}%
     \oraresz{Bard{\' o}czky}{te}{L}{Tt.}{2.5}%
\\ \oraresz{Furi{\' a}k}{te}{F}{Tt.}{2.5}%%
  \end{cellatartalom}}
\ora{\mitrakjonle}{3}{1}{5}
  {\begin{cellatartalom}%
     \oraresz{T{\' o}th}{a}{an4}{A.3}{1.25}%
\\ \oraresz{Varga Bo.}{a}{an1}{I.30}{1.25}%
\\ \oraresz{Seprenyi}{a}{an3}{II.41}{1.25}%
\\ \oraresz{Hajba}{a}{an2}{Szt.3}{1.25}%%
  \end{cellatartalom}}
\ora{\mitrakjonle}{3}{1}{6}
  {\begin{cellatartalom}%
     \oraresz{T{\' o}th}{a}{an4}{A.3}{1.25}%
\\ \oraresz{Varga Bo.}{a}{an1}{I.30}{1.25}%
\\ \oraresz{Bert{\' o}k}{a}{an2}{Szt.3}{1.25}%
\\ \oraresz{Seprenyi}{a}{an3}{II.41}{1.25}%%
  \end{cellatartalom}}
\ora{\mitrakjonle}{3}{1}{7}
  {\begin{cellatartalom}%
    %
  \end{cellatartalom}}
\ora{\mitrakjonle}{3}{1}{8}
  {\begin{cellatartalom}%
    %
  \end{cellatartalom}}
\ora{\mitrakjonle}{3}{1}{9}
  {\begin{cellatartalom}%
    %
  \end{cellatartalom}}
\ora{\mitrakjonle}{3}{2}{0}
  {\begin{cellatartalom}%
    %
  \end{cellatartalom}}
\ora{\mitrakjonle}{3}{2}{1}
  {\begin{cellatartalom}%
     \oraresz{Sum}{sp}{sp}{Sp.}{1}%
\\ \oraresz{Hal{\' a}sz}{n{\' e}}{n{\' e}1}{N.}{1}%
\\ \oraresz{Szabolcsi}{fr}{fr}{Fr.}{1}%
\\ \oraresz{Gianelli}{ol}{ol}{Ol.1}{1}%
\\ \oraresz{S{\' o}tin{\' e}}{n{\' e}}{n{\' e}2}{A.1}{1}%%
  \end{cellatartalom}}
\ora{\mitrakjonle}{3}{2}{2}
  {\begin{cellatartalom}%
     \oraresz{Sum}{sp}{sp}{Sp.}{1}%
\\ \oraresz{Hal{\' a}sz}{n{\' e}}{n{\' e}1}{N.}{1}%
\\ \oraresz{Szabolcsi}{fr}{fr}{Fr.}{1}%
\\ \oraresz{Gianelli}{ol}{ol}{Ol.1}{1}%
\\ \oraresz{S{\' o}tin{\' e}}{n{\' e}}{n{\' e}2}{A.1}{1}%%
  \end{cellatartalom}}
\ora{\mitrakjonle}{3}{2}{3}
  {\begin{cellatartalom}%
     \oraresz{Krajny{\' a}k}{mt}{tk}{I.30}{2.5}%
\\ \oraresz{T{\' o}th}{a}{szt}{F.10}{2.5}%%
  \end{cellatartalom}}
\ora{\mitrakjonle}{3}{2}{4}
  {\begin{cellatartalom}%
     \oraresz{Darab{\' a}nt}{mt}{szt}{I.30}{2.5}%
\\ \oraresz{Papp L.}{inf}{tk}{Tk.}{2.5}%%
  \end{cellatartalom}}
\ora{\mitrakjonle}{3}{2}{5}
  {\begin{cellatartalom}%
     \oraresz{Krajny{\' a}k}{inf}{szt}{Szt.2}{2.5}%
\\ \oraresz{P.-P{\' a}sztor}{i}{tk}{I.30}{2.5}%%
  \end{cellatartalom}}
\ora{\mitrakjonle}{3}{2}{6}
  {\begin{cellatartalom}%
     \oraresz{Juh{\' a}sz A.}{a}{an4}{A.3}{1.25}%
\\ \oraresz{Varga Bo.}{a}{an1}{N.}{1.25}%
\\ \oraresz{Bert{\' o}k}{a}{an2}{A.2}{1.25}%
\\ \oraresz{Seprenyi}{a}{an3}{II.41}{1.25}%%
  \end{cellatartalom}}
\ora{\mitrakjonle}{3}{2}{7}
  {\begin{cellatartalom}%
     \oraresz{Juh{\' a}sz A.}{a}{an4}{A.3}{1.25}%
\\ \oraresz{Varga Bo.}{a}{an1}{N.}{1.25}%
\\ \oraresz{Bert{\' o}k}{a}{an2}{A.2}{1.25}%
\\ \oraresz{Seprenyi}{a}{an3}{II.41}{1.25}%%
  \end{cellatartalom}}
\ora{\mitrakjonle}{3}{2}{8}
  {\begin{cellatartalom}%
     \oraresz{Krajny{\' a}k}{inf}{szak}{Szt.1}{5}%%
  \end{cellatartalom}}
\ora{\mitrakjonle}{3}{2}{9}
  {\begin{cellatartalom}%
    %
  \end{cellatartalom}}
\ora{\mitrakjonle}{3}{3}{0}
  {\begin{cellatartalom}%
     \oraresz{S{\' o}tin{\' e}}{n{\' e}}{n{\' e}2}{N.}{1}%
\\ \oraresz{Hal{\' a}sz}{n{\' e}}{n{\' e}1}{{\' E}.}{1}%
\\ \oraresz{Szabolcsi}{fr}{fr}{Fr.}{1}%
\\ \oraresz{Gianelli}{ol}{ol}{Ol.1}{1}%
\\ \oraresz{Sum}{sp}{sp}{Sp.}{1}%%
  \end{cellatartalom}}
\ora{\mitrakjonle}{3}{3}{1}
  {\begin{cellatartalom}%
     \oraresz{Krajny{\' a}k}{of}{-}{I.30}{5}%%
  \end{cellatartalom}}
\ora{\mitrakjonle}{3}{3}{2}
  {\begin{cellatartalom}%
     \oraresz{Hajba}{a}{an4}{Szt.3}{1.25}%
\\ \oraresz{Bert{\' o}k}{a}{an2}{A.1}{1.25}%
\\ \oraresz{Varga Bo.}{a}{an1}{I.30}{1.25}%
\\ \oraresz{Seprenyi}{a}{an3}{A.2}{1.25}%%
  \end{cellatartalom}}
\ora{\mitrakjonle}{3}{3}{3}
  {\begin{cellatartalom}%
     \oraresz{T{\' o}th}{a}{an4}{F.18}{1.25}%
\\ \oraresz{Varga Bo.}{a}{an1}{A.1}{1.25}%
\\ \oraresz{Husz{\' a}r}{a}{an2}{I.30}{1.25}%
\\ \oraresz{Seprenyi}{a}{an3}{A.2}{1.25}%%
  \end{cellatartalom}}
\ora{\mitrakjonle}{3}{3}{4}
  {\begin{cellatartalom}%
     \oraresz{T{\' o}th}{a}{an4}{F.18}{1.25}%
\\ \oraresz{Varga Bo.}{a}{an1}{A.1}{1.25}%
\\ \oraresz{Husz{\' a}r}{a}{an2}{I.30}{1.25}%
\\ \oraresz{Seprenyi}{a}{an3}{A.2}{1.25}%%
  \end{cellatartalom}}
\ora{\mitrakjonle}{3}{3}{5}
  {\begin{cellatartalom}%
     \oraresz{Krajny{\' a}k}{mt}{tk}{T.2}{2.5}%
\\ \oraresz{P.-P{\' a}sztor}{i}{szt}{I.39}{2.5}%%
  \end{cellatartalom}}
\ora{\mitrakjonle}{3}{3}{6}
  {\begin{cellatartalom}%
     \oraresz{Krajny{\' a}k}{inf}{szt}{Szt.1}{5}%%
  \end{cellatartalom}}
\ora{\mitrakjonle}{3}{3}{7}
  {\begin{cellatartalom}%
    %
  \end{cellatartalom}}
\ora{\mitrakjonle}{3}{3}{8}
  {\begin{cellatartalom}%
    %
  \end{cellatartalom}}
\ora{\mitrakjonle}{3}{3}{9}
  {\begin{cellatartalom}%
    %
  \end{cellatartalom}}
\ora{\mitrakjonle}{3}{4}{0}
  {\begin{cellatartalom}%
     \oraresz{T{\' o}th}{a}{an4}{I.30}{1.25}%
\\ \oraresz{Varga Bo.}{a}{an1}{F.8}{1.25}%
\\ \oraresz{Husz{\' a}r}{a}{an2}{A.2}{1.25}%
\\ \oraresz{Seprenyi}{a}{an3}{II.41}{1.25}%%
  \end{cellatartalom}}
\ora{\mitrakjonle}{3}{4}{1}
  {\begin{cellatartalom}%
     \oraresz{T{\' o}th}{a}{an4}{I.30}{1.25}%
\\ \oraresz{Varga Bo.}{a}{an1}{F.8}{1.25}%
\\ \oraresz{Husz{\' a}r}{a}{an2}{A.2}{1.25}%
\\ \oraresz{Seprenyi}{a}{an3}{II.41}{1.25}%%
  \end{cellatartalom}}
\ora{\mitrakjonle}{3}{4}{2}
  {\begin{cellatartalom}%
     \oraresz{S{\' o}tin{\' e}}{n{\' e}}{n{\' e}2}{II.51}{1}%
\\ \oraresz{Hal{\' a}sz}{n{\' e}}{n{\' e}1}{Szt.3}{1}%
\\ \oraresz{Szabolcsi}{fr}{fr}{II.42}{1}%
\\ \oraresz{Gianelli}{ol}{ol}{F.9}{1}%
\\ \oraresz{Sum}{sp}{sp}{Kt.}{1}%%
  \end{cellatartalom}}
\ora{\mitrakjonle}{3}{4}{3}
  {\begin{cellatartalom}%
     \oraresz{Krajny{\' a}k}{inf}{szt}{Szt.1}{2.5}%
\\ \oraresz{P.-P{\' a}sztor}{ny}{tk}{I.30}{2.5}%%
  \end{cellatartalom}}
\ora{\mitrakjonle}{3}{4}{4}
  {\begin{cellatartalom}%
     \oraresz{Darab{\' a}nt}{mt}{szt}{M.}{2.5}%
\\ \oraresz{Krajny{\' a}k}{mt}{tk}{T.2}{2.5}%%
  \end{cellatartalom}}
\ora{\mitrakjonle}{3}{4}{5}
  {\begin{cellatartalom}%
     \oraresz{Papp L.}{inf}{tk}{Tk.}{5}%%
  \end{cellatartalom}}
\ora{\mitrakjonle}{3}{4}{6}
  {\begin{cellatartalom}%
     \oraresz{P{\' o}k}{inf}{tk}{Tk.}{5}%%
  \end{cellatartalom}}
\ora{\mitrakjonle}{3}{4}{7}
  {\begin{cellatartalom}%
    %
  \end{cellatartalom}}
\ora{\mitrakjonle}{3}{4}{8}
  {\begin{cellatartalom}%
     \oraresz{B{\" o}g{\" o}s}{a}{fz}{A.1}{5}%%
  \end{cellatartalom}}
\ora{\mitrakjonle}{3}{4}{9}
  {\begin{cellatartalom}%
    %
  \end{cellatartalom}}
\ora{\mitrakjonle}{3}{5}{0}
  {\begin{cellatartalom}%
     \oraresz{Bard{\' o}czky}{te}{L}{Tt.}{2.5}%
\\ \oraresz{Furi{\' a}k}{te}{F}{Tt.}{2.5}%%
  \end{cellatartalom}}
\ora{\mitrakjonle}{3}{5}{1}
  {\begin{cellatartalom}%
     \oraresz{Bard{\' o}czky}{te}{L}{Tt.}{2.5}%
\\ \oraresz{Furi{\' a}k}{te}{F}{Tt.}{2.5}%%
  \end{cellatartalom}}
\ora{\mitrakjonle}{3}{5}{2}
  {\begin{cellatartalom}%
     \oraresz{T{\' o}th}{a}{an4}{R.}{1.25}%
\\ \oraresz{Varga Bo.}{a}{an1}{I.30}{1.25}%
\\ \oraresz{Husz{\' a}r}{a}{an2}{I.32}{1.25}%
\\ \oraresz{Seprenyi}{a}{an3}{Szt.3}{1.25}%%
  \end{cellatartalom}}
\ora{\mitrakjonle}{3}{5}{3}
  {\begin{cellatartalom}%
     \oraresz{T{\' o}th}{a}{an4}{R.}{1.25}%
\\ \oraresz{Varga Bo.}{a}{an1}{I.30}{1.25}%
\\ \oraresz{Husz{\' a}r}{a}{an2}{I.32}{1.25}%
\\ \oraresz{Seprenyi}{a}{an3}{Szt.3}{1.25}%%
  \end{cellatartalom}}
\ora{\mitrakjonle}{3}{5}{4}
  {\begin{cellatartalom}%
     \oraresz{P.-P{\' a}sztor}{ny}{szt}{I.30}{2.5}%
\\ \oraresz{T{\' o}th}{a}{tk}{F.18}{2.5}%%
  \end{cellatartalom}}
\ora{\mitrakjonle}{3}{5}{5}
  {\begin{cellatartalom}%
    %
  \end{cellatartalom}}
\ora{\mitrakjonle}{3}{5}{6}
  {\begin{cellatartalom}%
     \oraresz{S{\' o}tin{\' e}}{n{\' e}}{n{\' e}2}{N.}{1}%
\\ \oraresz{Hal{\' a}sz}{n{\' e}}{n{\' e}1}{{\' E}.}{1}%
\\ \oraresz{Szabolcsi}{fr}{fr}{Fr.}{1}%
\\ \oraresz{Gianelli}{ol}{ol}{L.}{1}%
\\ \oraresz{Sum}{sp}{sp}{Sp.}{1}%%
  \end{cellatartalom}}
\ora{\mitrakjonle}{3}{5}{7}
  {\begin{cellatartalom}%
    %
  \end{cellatartalom}}
\ora{\mitrakjonle}{3}{5}{8}
  {\begin{cellatartalom}%
    %
  \end{cellatartalom}}
\ora{\mitrakjonle}{3}{5}{9}
  {\begin{cellatartalom}%
    %
  \end{cellatartalom}}
\ora{\mitrakjonle}{4}{1}{0}
  {\begin{cellatartalom}%
     \oraresz{V. Bartha}{i}{-}{F.11}{6}%%
  \end{cellatartalom}}
\ora{\mitrakjonle}{4}{1}{1}
  {\begin{cellatartalom}%
     \oraresz{Horv{\' a}th G.}{k{\' e}}{1}{K.ea}{3}%
\\ \oraresz{M{\" u}ller}{f{\" o}}{2}{L.}{3}%%
  \end{cellatartalom}}
\ora{\mitrakjonle}{4}{1}{2}
  {\begin{cellatartalom}%
     \oraresz{Fajcs{\' a}k}{ol}{ol1}{F.11}{2}%
\\ \oraresz{Somk{\" o}vi}{ol}{ol3}{F.12}{2}%
\\ \oraresz{Bellesi}{ol}{ol2}{L.}{2}%%
  \end{cellatartalom}}
\ora{\mitrakjonle}{4}{1}{3}
  {\begin{cellatartalom}%
     \oraresz{Schlachter}{mt}{2}{F.11}{3}%
\\ \oraresz{Cserv{\' a}k}{inf}{1}{Szt.1}{3}%%
  \end{cellatartalom}}
\ora{\mitrakjonle}{4}{1}{4}
  {\begin{cellatartalom}%
     \oraresz{Kudella}{fr}{fr}{Fr.}{1}%
\\ \oraresz{Pap}{sp}{sp2}{Sp.}{1}%
\\ \oraresz{Tandory}{a}{an3}{I.31}{1}%
\\ \oraresz{T{\' o}th}{a}{an1}{F.12}{1}%
\\ \oraresz{Szabados}{a}{an2}{A.1}{1}%
\\ \oraresz{Szab{\' o} A.}{or}{or}{Kt.}{1}%%
  \end{cellatartalom}}
\ora{\mitrakjonle}{4}{1}{5}
  {\begin{cellatartalom}%
     \oraresz{Barth{\' a}n{\' e}}{{\' e}n}{-}{{\' E}.}{6}%%
  \end{cellatartalom}}
\ora{\mitrakjonle}{4}{1}{6}
  {\begin{cellatartalom}%
     \oraresz{Lovas}{t{\" o}}{2}{II.40}{3}%
\\ \oraresz{M{\" u}ller}{f{\" o}}{1}{L.}{3}%%
  \end{cellatartalom}}
\ora{\mitrakjonle}{4}{1}{7}
  {\begin{cellatartalom}%
     \oraresz{Fajcs{\' a}k}{ol}{ol2}{Ol.1}{6}%%
  \end{cellatartalom}}
\ora{\mitrakjonle}{4}{1}{8}
  {\begin{cellatartalom}%
    %
  \end{cellatartalom}}
\ora{\mitrakjonle}{4}{1}{9}
  {\begin{cellatartalom}%
    %
  \end{cellatartalom}}
\ora{\mitrakjonle}{4}{2}{0}
  {\begin{cellatartalom}%
     \oraresz{Fajcs{\' a}k}{ol}{ol3}{Ol.1}{6}%%
  \end{cellatartalom}}
\ora{\mitrakjonle}{4}{2}{1}
  {\begin{cellatartalom}%
     \oraresz{Horv{\' a}th}{fi}{-}{Fiz.}{6}%%
  \end{cellatartalom}}
\ora{\mitrakjonle}{4}{2}{2}
  {\begin{cellatartalom}%
     \oraresz{Krajny{\' a}k}{mt}{1}{F.11}{3}%
\\ \oraresz{Lovas}{t{\" o}}{2}{II.40}{3}%%
  \end{cellatartalom}}
\ora{\mitrakjonle}{4}{2}{3}
  {\begin{cellatartalom}%
     \oraresz{Fajcs{\' a}k}{ol}{ol1}{F.11}{1.5}%
\\ \oraresz{Somk{\" o}vi}{ol}{ol3}{Ol.2}{1.5}%
\\ \oraresz{Bellesi}{ol}{ol2}{L.}{1.5}%
\\ \oraresz{Sum}{mid}{mid}{FS.t}{1.5}%%
  \end{cellatartalom}}
\ora{\mitrakjonle}{4}{2}{4}
  {\begin{cellatartalom}%
     \oraresz{Vas E.}{te}{L}{Tt.}{2}%
\\ \oraresz{Lad{\' a}nyi}{te}{L}{Tt.}{2}%
\\ \oraresz{De{\' a}k}{te}{F}{Tt.}{2}%%
  \end{cellatartalom}}
\ora{\mitrakjonle}{4}{2}{5}
  {\begin{cellatartalom}%
     \oraresz{Vas E.}{te}{L}{Tt.}{2}%
\\ \oraresz{Lad{\' a}nyi}{te}{L}{Tt.}{2}%
\\ \oraresz{De{\' a}k}{te}{F}{Tt.}{2}%%
  \end{cellatartalom}}
\ora{\mitrakjonle}{4}{2}{6}
  {\begin{cellatartalom}%
     \oraresz{Balanyi}{r}{-}{R.}{6}%%
  \end{cellatartalom}}
\ora{\mitrakjonle}{4}{2}{7}
  {\begin{cellatartalom}%
     \oraresz{Horv{\' a}th G.}{k{\' e}}{2}{II.48}{3}%
\\ \oraresz{Cserv{\' a}k}{inf}{1}{Szt.1}{3}%%
  \end{cellatartalom}}
\ora{\mitrakjonle}{4}{2}{8}
  {\begin{cellatartalom}%
    %
  \end{cellatartalom}}
\ora{\mitrakjonle}{4}{2}{9}
  {\begin{cellatartalom}%
    %
  \end{cellatartalom}}
\ora{\mitrakjonle}{4}{3}{0}
  {\begin{cellatartalom}%
    %
  \end{cellatartalom}}
\ora{\mitrakjonle}{4}{3}{1}
  {\begin{cellatartalom}%
     \oraresz{Tandory}{a}{an3}{A.3}{1}%
\\ \oraresz{Kudella}{fr}{fr}{Fr.}{1}%
\\ \oraresz{Pap}{sp}{sp2}{Sp.}{1}%
\\ \oraresz{Szabados}{a}{an2}{I.31}{1}%
\\ \oraresz{T{\' o}th}{a}{an1}{II.43}{1}%
\\ \oraresz{Szab{\' o} A.}{or}{or}{F.16}{1}%%
  \end{cellatartalom}}
\ora{\mitrakjonle}{4}{3}{2}
  {\begin{cellatartalom}%
     \oraresz{Tandory}{a}{an3}{A.3}{1}%
\\ \oraresz{Kudella}{fr}{fr}{Fr.}{1}%
\\ \oraresz{Pap}{sp}{sp2}{Sp.}{1}%
\\ \oraresz{Szabados}{a}{an2}{I.31}{1}%
\\ \oraresz{T{\' o}th}{a}{an1}{II.43}{1}%
\\ \oraresz{Szab{\' o} A.}{or}{or}{F.16}{1}%%
  \end{cellatartalom}}
\ora{\mitrakjonle}{4}{3}{3}
  {\begin{cellatartalom}%
     \oraresz{Schlachter}{mt}{2}{M.}{3}%
\\ \oraresz{Lovas}{t{\" o}}{1}{II.40}{3}%%
  \end{cellatartalom}}
\ora{\mitrakjonle}{4}{3}{4}
  {\begin{cellatartalom}%
     \oraresz{Krajny{\' a}k}{mt}{1}{F.11}{3}%
\\ \oraresz{M{\" u}ller}{f{\" o}}{2}{L.}{3}%%
  \end{cellatartalom}}
\ora{\mitrakjonle}{4}{3}{5}
  {\begin{cellatartalom}%
     \oraresz{Horv{\' a}th G.}{bi}{-}{B.ea}{6}%%
  \end{cellatartalom}}
\ora{\mitrakjonle}{4}{3}{6}
  {\begin{cellatartalom}%
     \oraresz{Horv{\' a}th}{fi}{-}{F.11}{6}%%
  \end{cellatartalom}}
\ora{\mitrakjonle}{4}{3}{7}
  {\begin{cellatartalom}%
     \oraresz{Fajcs{\' a}k}{ol}{ol3}{Ol.1}{6}%%
  \end{cellatartalom}}
\ora{\mitrakjonle}{4}{3}{8}
  {\begin{cellatartalom}%
    %
  \end{cellatartalom}}
\ora{\mitrakjonle}{4}{3}{9}
  {\begin{cellatartalom}%
    %
  \end{cellatartalom}}
\ora{\mitrakjonle}{4}{4}{0}
  {\begin{cellatartalom}%
     \oraresz{Bellesi}{ol}{ol1}{Ol.1}{6}%%
  \end{cellatartalom}}
\ora{\mitrakjonle}{4}{4}{1}
  {\begin{cellatartalom}%
     \oraresz{V. Bartha}{ny}{-}{F.11}{6}%%
  \end{cellatartalom}}
\ora{\mitrakjonle}{4}{4}{2}
  {\begin{cellatartalom}%
     \oraresz{Horv{\' a}th G.}{k{\' e}}{1}{K.gy}{3}%
\\ \oraresz{Cserv{\' a}k}{inf}{2}{Szt.1}{3}%%
  \end{cellatartalom}}
\ora{\mitrakjonle}{4}{4}{3}
  {\begin{cellatartalom}%
     \oraresz{Vas E.}{te}{L}{Tt.}{2}%
\\ \oraresz{Lad{\' a}nyi}{te}{L}{Tt.}{2}%
\\ \oraresz{De{\' a}k}{te}{F}{Tt.}{2}%%
  \end{cellatartalom}}
\ora{\mitrakjonle}{4}{4}{4}
  {\begin{cellatartalom}%
     \oraresz{Kudella}{fr}{fr}{Fr.}{1}%
\\ \oraresz{Pap}{sp}{sp2}{Sp.}{1}%
\\ \oraresz{Tandory}{a}{an3}{I.34}{1}%
\\ \oraresz{T{\' o}th}{a}{an1}{A.2}{1}%
\\ \oraresz{Szabados}{a}{an2}{{\' E}.}{1}%
\\ \oraresz{Szab{\' o} A.}{or}{or}{I.30}{1}%%
  \end{cellatartalom}}
\ora{\mitrakjonle}{4}{4}{5}
  {\begin{cellatartalom}%
     \oraresz{Schlachter}{mt}{2}{F.11}{3}%
\\ \oraresz{M{\" u}ller}{f{\" o}}{1}{L.}{3}%%
  \end{cellatartalom}}
\ora{\mitrakjonle}{4}{4}{6}
  {\begin{cellatartalom}%
     \oraresz{Fajcs{\' a}k}{ol}{ol1}{Ol.2}{2}%
\\ \oraresz{Somk{\" o}vi}{ol}{ol3}{Ol.1}{2}%
\\ \oraresz{Bellesi}{ol}{ol2}{F.11}{2}%%
  \end{cellatartalom}}
\ora{\mitrakjonle}{4}{4}{7}
  {\begin{cellatartalom}%
     \oraresz{Fajcs{\' a}k}{ol}{ol3}{Ol.1}{6}%%
  \end{cellatartalom}}
\ora{\mitrakjonle}{4}{4}{8}
  {\begin{cellatartalom}%
    %
  \end{cellatartalom}}
\ora{\mitrakjonle}{4}{4}{9}
  {\begin{cellatartalom}%
    %
  \end{cellatartalom}}
\ora{\mitrakjonle}{4}{5}{0}
  {\begin{cellatartalom}%
     \oraresz{Fajcs{\' a}k}{ol}{ol2}{Ol.2}{3}%
\\ \oraresz{Bellesi}{ol}{ol1}{F.11}{3}%%
  \end{cellatartalom}}
\ora{\mitrakjonle}{4}{5}{1}
  {\begin{cellatartalom}%
     \oraresz{Fajcs{\' a}k}{ol}{ol2}{Ol.2}{3}%
\\ \oraresz{Bellesi}{ol}{ol1}{F.11}{3}%%
  \end{cellatartalom}}
\ora{\mitrakjonle}{4}{5}{2}
  {\begin{cellatartalom}%
     \oraresz{Horv{\' a}th G.}{bi}{-}{F.11}{6}%%
  \end{cellatartalom}}
\ora{\mitrakjonle}{4}{5}{3}
  {\begin{cellatartalom}%
     \oraresz{Krajny{\' a}k}{mt}{1}{F.11}{3}%
\\ \oraresz{Cserv{\' a}k}{inf}{2}{Szt.1}{3}%%
  \end{cellatartalom}}
\ora{\mitrakjonle}{4}{5}{4}
  {\begin{cellatartalom}%
     \oraresz{V. Bartha}{i}{-}{F.11}{6}%%
  \end{cellatartalom}}
\ora{\mitrakjonle}{4}{5}{5}
  {\begin{cellatartalom}%
     \oraresz{Fajcs{\' a}k}{of}{-}{F.11}{6}%%
  \end{cellatartalom}}
\ora{\mitrakjonle}{4}{5}{6}
  {\begin{cellatartalom}%
     \oraresz{Horv{\' a}th G.}{k{\' e}}{2}{K.ea}{3}%
\\ \oraresz{Lovas}{t{\" o}}{1}{II.40}{3}%%
  \end{cellatartalom}}
\ora{\mitrakjonle}{4}{5}{7}
  {\begin{cellatartalom}%
    %
  \end{cellatartalom}}
\ora{\mitrakjonle}{4}{5}{8}
  {\begin{cellatartalom}%
    %
  \end{cellatartalom}}
\ora{\mitrakjonle}{4}{5}{9}
  {\begin{cellatartalom}%
    %
  \end{cellatartalom}}
\ora{\mitrakjonle}{5}{1}{0}
  {\begin{cellatartalom}%
    %
  \end{cellatartalom}}
\ora{\mitrakjonle}{5}{1}{1}
  {\begin{cellatartalom}%
     \oraresz{Tandory}{f{\" o}}{-}{F.12}{2}%%
  \end{cellatartalom}}
\ora{\mitrakjonle}{5}{1}{2}
  {\begin{cellatartalom}%
     \oraresz{Barth{\' a}n{\' e}}{{\' e}n}{-}{{\' E}.}{2}%%
  \end{cellatartalom}}
\ora{\mitrakjonle}{5}{1}{3}
  {\begin{cellatartalom}%
     \oraresz{F{\" o}ldesi}{inf}{1}{Szt.2}{1}%
\\ \oraresz{Sebesty{\' e}n}{mt}{2}{F.12}{1}%%
  \end{cellatartalom}}
\ora{\mitrakjonle}{5}{1}{4}
  {\begin{cellatartalom}%
     \oraresz{Balanyi}{r}{-}{R.}{2}%%
  \end{cellatartalom}}
\ora{\mitrakjonle}{5}{1}{5}
  {\begin{cellatartalom}%
     \oraresz{Sz{\' e}kely}{a}{an2}{I.31}{1}%
\\ \oraresz{Kassai}{a}{an1}{F.12}{1}%%
  \end{cellatartalom}}
\ora{\mitrakjonle}{5}{1}{6}
  {\begin{cellatartalom}%
     \oraresz{Schlachter}{mt}{1}{F.12}{1}%
\\ \oraresz{Kudella}{i}{2}{I.34}{1}%%
  \end{cellatartalom}}
\ora{\mitrakjonle}{5}{1}{7}
  {\begin{cellatartalom}%
     \oraresz{Schlachter}{mt}{1}{F.12}{1}%
\\ \oraresz{Kudella}{ny}{2}{I.34}{1}%%
  \end{cellatartalom}}
\ora{\mitrakjonle}{5}{1}{8}
  {\begin{cellatartalom}%
    %
  \end{cellatartalom}}
\ora{\mitrakjonle}{5}{1}{9}
  {\begin{cellatartalom}%
    %
  \end{cellatartalom}}
\ora{\mitrakjonle}{5}{2}{0}
  {\begin{cellatartalom}%
     \oraresz{Bard{\' o}czky}{te}{L}{Tt.}{2}%%
  \end{cellatartalom}}
\ora{\mitrakjonle}{5}{2}{1}
  {\begin{cellatartalom}%
     \oraresz{Furi{\' a}k}{te}{F}{Tt.}{1}%
\\ \oraresz{Bard{\' o}czky}{te}{L}{Tt.}{1}%%
  \end{cellatartalom}}
\ora{\mitrakjonle}{5}{2}{2}
  {\begin{cellatartalom}%
     \oraresz{Cserv{\' a}k}{inf}{2}{Szt.1}{1}%
\\ \oraresz{Kudella}{i}{1}{F.12}{1}%%
  \end{cellatartalom}}
\ora{\mitrakjonle}{5}{2}{3}
  {\begin{cellatartalom}%
     \oraresz{Menyh{\' a}rt}{bi}{-}{F.12}{2}%%
  \end{cellatartalom}}
\ora{\mitrakjonle}{5}{2}{4}
  {\begin{cellatartalom}%
     \oraresz{Pat{\' o}}{t{\" o}}{1}{F.16}{1}%
\\ \oraresz{Kerling}{t{\" o}}{2}{F.12}{1}%%
  \end{cellatartalom}}
\ora{\mitrakjonle}{5}{2}{5}
  {\begin{cellatartalom}%
     \oraresz{Malatinszky}{k{\' e}}{1}{K.gy}{1}%
\\ \oraresz{Sebesty{\' e}n}{mt}{2}{M.}{1}%%
  \end{cellatartalom}}
\ora{\mitrakjonle}{5}{2}{6}
  {\begin{cellatartalom}%
     \oraresz{Sz{\' e}kely}{a}{an1}{F.17}{1}%
\\ \oraresz{Kassai}{a}{an2}{A.1}{1}%%
  \end{cellatartalom}}
\ora{\mitrakjonle}{5}{2}{7}
  {\begin{cellatartalom}%
     \oraresz{Sz{\' e}kely}{a}{an1}{F.17}{1}%
\\ \oraresz{Kassai}{a}{an2}{A.1}{1}%%
  \end{cellatartalom}}
\ora{\mitrakjonle}{5}{2}{8}
  {\begin{cellatartalom}%
    %
  \end{cellatartalom}}
\ora{\mitrakjonle}{5}{2}{9}
  {\begin{cellatartalom}%
    %
  \end{cellatartalom}}
\ora{\mitrakjonle}{5}{3}{0}
  {\begin{cellatartalom}%
    %
  \end{cellatartalom}}
\ora{\mitrakjonle}{5}{3}{1}
  {\begin{cellatartalom}%
     \oraresz{P{\' e}teri}{fi}{-}{Fiz.}{2}%%
  \end{cellatartalom}}
\ora{\mitrakjonle}{5}{3}{2}
  {\begin{cellatartalom}%
     \oraresz{Dickerson}{a}{an2}{F.11}{1}%
\\ \oraresz{Sz{\' e}kely}{a}{an1}{F.12}{1}%%
  \end{cellatartalom}}
\ora{\mitrakjonle}{5}{3}{3}
  {\begin{cellatartalom}%
     \oraresz{Kudella}{i}{1}{F.12}{1}%
\\ \oraresz{Malatinszky}{k{\' e}}{2}{K.gy}{1}%%
  \end{cellatartalom}}
\ora{\mitrakjonle}{5}{3}{4}
  {\begin{cellatartalom}%
     \oraresz{Schlachter}{mt}{1}{M.}{1}%
\\ \oraresz{Sebesty{\' e}n}{mt}{2}{F.12}{1}%%
  \end{cellatartalom}}
\ora{\mitrakjonle}{5}{3}{5}
  {\begin{cellatartalom}%
     \oraresz{Kudella}{i}{2}{F.12}{1}%
\\ \oraresz{F{\" o}ldesi}{inf}{1}{Szt.1}{1}%%
  \end{cellatartalom}}
\ora{\mitrakjonle}{5}{3}{6}
  {\begin{cellatartalom}%
     \oraresz{Furi{\' a}k}{te}{F}{Tt.}{1}%
\\ \oraresz{Bard{\' o}czky}{te}{L}{Tt.}{1}%%
  \end{cellatartalom}}
\ora{\mitrakjonle}{5}{3}{7}
  {\begin{cellatartalom}%
     \oraresz{Furi{\' a}k}{te}{F}{Tt.}{2}%%
  \end{cellatartalom}}
\ora{\mitrakjonle}{5}{3}{8}
  {\begin{cellatartalom}%
    %
  \end{cellatartalom}}
\ora{\mitrakjonle}{5}{3}{9}
  {\begin{cellatartalom}%
    %
  \end{cellatartalom}}
\ora{\mitrakjonle}{5}{4}{0}
  {\begin{cellatartalom}%
     \oraresz{Cserv{\' a}k}{inf}{2}{Szt.1}{2}%%
  \end{cellatartalom}}
\ora{\mitrakjonle}{5}{4}{1}
  {\begin{cellatartalom}%
     \oraresz{Schlachter}{mt}{1}{F.12}{1}%
\\ \oraresz{Sebesty{\' e}n}{mt}{2}{I.39}{1}%%
  \end{cellatartalom}}
\ora{\mitrakjonle}{5}{4}{2}
  {\begin{cellatartalom}%
     \oraresz{Kassai}{a}{an1}{F.11}{1}%
\\ \oraresz{Sz{\' e}kely}{a}{an2}{A.2}{1}%%
  \end{cellatartalom}}
\ora{\mitrakjonle}{5}{4}{3}
  {\begin{cellatartalom}%
     \oraresz{Kassai}{a}{an1}{F.11}{1}%
\\ \oraresz{Sz{\' e}kely}{a}{an2}{A.2}{1}%%
  \end{cellatartalom}}
\ora{\mitrakjonle}{5}{4}{4}
  {\begin{cellatartalom}%
     \oraresz{Pat{\' o}}{t{\" o}}{1}{F.12}{1}%
\\ \oraresz{Kerling}{t{\" o}}{2}{II.50}{1}%%
  \end{cellatartalom}}
\ora{\mitrakjonle}{5}{4}{5}
  {\begin{cellatartalom}%
     \oraresz{Tandory}{f{\" o}}{-}{F.10}{2}%%
  \end{cellatartalom}}
\ora{\mitrakjonle}{5}{4}{6}
  {\begin{cellatartalom}%
     \oraresz{Sz{\' e}kely}{of}{-}{F.12}{2}%%
  \end{cellatartalom}}
\ora{\mitrakjonle}{5}{4}{7}
  {\begin{cellatartalom}%
     \oraresz{Kudella}{ny}{1}{I.31}{2}%%
  \end{cellatartalom}}
\ora{\mitrakjonle}{5}{4}{8}
  {\begin{cellatartalom}%
    %
  \end{cellatartalom}}
\ora{\mitrakjonle}{5}{4}{9}
  {\begin{cellatartalom}%
    %
  \end{cellatartalom}}
\ora{\mitrakjonle}{5}{5}{0}
  {\begin{cellatartalom}%
    %
  \end{cellatartalom}}
\ora{\mitrakjonle}{5}{5}{1}
  {\begin{cellatartalom}%
     \oraresz{P{\' e}teri}{fi}{-}{F.12}{2}%%
  \end{cellatartalom}}
\ora{\mitrakjonle}{5}{5}{2}
  {\begin{cellatartalom}%
     \oraresz{Kudella}{i}{2}{F.12}{1}%
\\ \oraresz{Malatinszky}{k{\' e}}{1}{K.gy}{1}%%
  \end{cellatartalom}}
\ora{\mitrakjonle}{5}{5}{3}
  {\begin{cellatartalom}%
     \oraresz{Menyh{\' a}rt}{bi}{-}{B.ea}{2}%%
  \end{cellatartalom}}
\ora{\mitrakjonle}{5}{5}{4}
  {\begin{cellatartalom}%
     \oraresz{Kudella}{i}{1}{F.12}{1}%
\\ \oraresz{Malatinszky}{k{\' e}}{2}{K.ea}{1}%%
  \end{cellatartalom}}
\ora{\mitrakjonle}{5}{5}{5}
  {\begin{cellatartalom}%
     \oraresz{Sz{\' e}kely}{a}{an1}{A.1}{1}%
\\ \oraresz{Kassai}{a}{an2}{F.12}{1}%%
  \end{cellatartalom}}
\ora{\mitrakjonle}{5}{5}{6}
  {\begin{cellatartalom}%
     \oraresz{Dickerson}{a}{an1}{A.1}{1}%
\\ \oraresz{Kassai}{a}{an2}{F.12}{1}%%
  \end{cellatartalom}}
\ora{\mitrakjonle}{5}{5}{7}
  {\begin{cellatartalom}%
    %
  \end{cellatartalom}}
\ora{\mitrakjonle}{5}{5}{8}
  {\begin{cellatartalom}%
    %
  \end{cellatartalom}}
\ora{\mitrakjonle}{5}{5}{9}
  {\begin{cellatartalom}%
    %
  \end{cellatartalom}}
\ora{\mitrakjonle}{6}{1}{0}
  {\begin{cellatartalom}%
    %
  \end{cellatartalom}}
\ora{\mitrakjonle}{6}{1}{1}
  {\begin{cellatartalom}%
     \oraresz{Szendrei}{t{\" o}}{hu}{F.8}{1.5}%
\\ \oraresz{Dezs{\H o}}{mt}{mt}{T.2}{1.5}%%
  \end{cellatartalom}}
\ora{\mitrakjonle}{6}{1}{2}
  {\begin{cellatartalom}%
     \oraresz{De{\' a}k}{te}{F}{Tt.}{1.5}%
\\ \oraresz{Bard{\' o}czky}{te}{L}{Tt.}{1.5}%%
  \end{cellatartalom}}
\ora{\mitrakjonle}{6}{1}{3}
  {\begin{cellatartalom}%
     \oraresz{B{\" o}g{\" o}s}{a}{an3}{A.3}{1}%
\\ \oraresz{Husz{\' a}r}{a}{an1}{I.30}{1}%
\\ \oraresz{Varga Bo.}{a}{an2}{II.41}{1}%%
  \end{cellatartalom}}
\ora{\mitrakjonle}{6}{1}{4}
  {\begin{cellatartalom}%
     \oraresz{B{\" o}g{\" o}s}{a}{an3}{A.3}{1}%
\\ \oraresz{Husz{\' a}r}{a}{an1}{I.30}{1}%
\\ \oraresz{Varga Bo.}{a}{an2}{II.41}{1}%%
  \end{cellatartalom}}
\ora{\mitrakjonle}{6}{1}{5}
  {\begin{cellatartalom}%
     \oraresz{S{\' o}skuthyn{\' e}}{bi}{-}{K.ea}{3}%%
  \end{cellatartalom}}
\ora{\mitrakjonle}{6}{1}{6}
  {\begin{cellatartalom}%
     \oraresz{Jantner}{mt}{hu}{M.}{1.5}%
\\ \oraresz{L{\' a}z{\' a}r}{i}{mt}{I.39}{1.5}%%
  \end{cellatartalom}}
\ora{\mitrakjonle}{6}{1}{7}
  {\begin{cellatartalom}%
     \oraresz{Horv{\' a}th}{fi}{-}{F.8}{3}%%
  \end{cellatartalom}}
\ora{\mitrakjonle}{6}{1}{8}
  {\begin{cellatartalom}%
    %
  \end{cellatartalom}}
\ora{\mitrakjonle}{6}{1}{9}
  {\begin{cellatartalom}%
    %
  \end{cellatartalom}}
\ora{\mitrakjonle}{6}{2}{0}
  {\begin{cellatartalom}%
     \oraresz{Adamis}{f{\" o}}{-}{F.8}{3}%%
  \end{cellatartalom}}
\ora{\mitrakjonle}{6}{2}{1}
  {\begin{cellatartalom}%
     \oraresz{Jantner}{mt}{hu}{F.8}{1.5}%
\\ \oraresz{Dezs{\H o}}{mt}{mt}{T.2}{1.5}%%
  \end{cellatartalom}}
\ora{\mitrakjonle}{6}{2}{2}
  {\begin{cellatartalom}%
     \oraresz{S{\' o}skuthyn{\' e}}{k{\' e}}{hu}{K.gy}{1.5}%
\\ \oraresz{Dezs{\H o}}{mt}{mt}{T.2}{1.5}%%
  \end{cellatartalom}}
\ora{\mitrakjonle}{6}{2}{3}
  {\begin{cellatartalom}%
     \oraresz{F{\" o}ldesi}{inf}{mt}{Szt.3}{1.5}%
\\ \oraresz{Teremy}{ny}{hu}{F.8}{1.5}%%
  \end{cellatartalom}}
\ora{\mitrakjonle}{6}{2}{4}
  {\begin{cellatartalom}%
     \oraresz{Adamis}{t{\" o}}{mt}{F.8}{1.5}%
\\ \oraresz{M{\" u}ller}{la}{hu}{L.}{1.5}%%
  \end{cellatartalom}}
\ora{\mitrakjonle}{6}{2}{5}
  {\begin{cellatartalom}%
     \oraresz{Szendrei}{t{\" o}}{hu}{II.50}{1.5}%
\\ \oraresz{L{\' a}z{\' a}r}{ny}{mt}{F.8}{1.5}%%
  \end{cellatartalom}}
\ora{\mitrakjonle}{6}{2}{6}
  {\begin{cellatartalom}%
     \oraresz{F{\" o}ldesi}{inf}{hu}{Szt.3}{1.5}%
\\ \oraresz{S{\' o}skuthyn{\' e}}{k{\' e}}{mt}{F.8}{1.5}%%
  \end{cellatartalom}}
\ora{\mitrakjonle}{6}{2}{7}
  {\begin{cellatartalom}%
    %
  \end{cellatartalom}}
\ora{\mitrakjonle}{6}{2}{8}
  {\begin{cellatartalom}%
    %
  \end{cellatartalom}}
\ora{\mitrakjonle}{6}{2}{9}
  {\begin{cellatartalom}%
    %
  \end{cellatartalom}}
\ora{\mitrakjonle}{6}{3}{0}
  {\begin{cellatartalom}%
     \oraresz{B{\" o}g{\" o}s}{a}{an3}{F.17}{1}%
\\ \oraresz{Husz{\' a}r}{a}{an1}{A.2}{1}%
\\ \oraresz{Varga Bo.}{a}{an2}{A.1}{1}%%
  \end{cellatartalom}}
\ora{\mitrakjonle}{6}{3}{1}
  {\begin{cellatartalom}%
     \oraresz{B{\" o}g{\" o}s}{a}{an3}{F.17}{1}%
\\ \oraresz{Husz{\' a}r}{a}{an1}{A.2}{1}%
\\ \oraresz{Varga Bo.}{a}{an2}{A.1}{1}%%
  \end{cellatartalom}}
\ora{\mitrakjonle}{6}{3}{2}
  {\begin{cellatartalom}%
     \oraresz{F{\" o}ldesi}{inf}{hu}{Szt.1}{1.5}%
\\ \oraresz{S{\' o}skuthyn{\' e}}{k{\' e}}{mt}{K.gy}{1.5}%%
  \end{cellatartalom}}
\ora{\mitrakjonle}{6}{3}{3}
  {\begin{cellatartalom}%
     \oraresz{S{\' o}skuthyn{\' e}}{bi}{-}{F.8}{3}%%
  \end{cellatartalom}}
\ora{\mitrakjonle}{6}{3}{4}
  {\begin{cellatartalom}%
     \oraresz{L{\' a}z{\' a}r}{i}{mt}{F.8}{1.5}%
\\ \oraresz{Teremy}{i}{hu}{F.10}{1.5}%%
  \end{cellatartalom}}
\ora{\mitrakjonle}{6}{3}{5}
  {\begin{cellatartalom}%
     \oraresz{Adamis}{of}{-}{F.8}{3}%%
  \end{cellatartalom}}
\ora{\mitrakjonle}{6}{3}{6}
  {\begin{cellatartalom}%
     \oraresz{Dezs{\H o}}{mt}{mt}{T.2}{1.5}%
\\ \oraresz{P{\' o}k}{r}{hu}{II.51}{1.5}%%
  \end{cellatartalom}}
\ora{\mitrakjonle}{6}{3}{7}
  {\begin{cellatartalom}%
    %
  \end{cellatartalom}}
\ora{\mitrakjonle}{6}{3}{8}
  {\begin{cellatartalom}%
    %
  \end{cellatartalom}}
\ora{\mitrakjonle}{6}{3}{9}
  {\begin{cellatartalom}%
    %
  \end{cellatartalom}}
\ora{\mitrakjonle}{6}{4}{0}
  {\begin{cellatartalom}%
     \oraresz{De{\' a}k}{te}{F}{Tt.}{1.5}%
\\ \oraresz{Bard{\' o}czky}{te}{L}{Tt.}{1.5}%%
  \end{cellatartalom}}
\ora{\mitrakjonle}{6}{4}{1}
  {\begin{cellatartalom}%
     \oraresz{De{\' a}k}{te}{F}{Tt.}{1.5}%
\\ \oraresz{Bard{\' o}czky}{te}{L}{Tt.}{1.5}%%
  \end{cellatartalom}}
\ora{\mitrakjonle}{6}{4}{2}
  {\begin{cellatartalom}%
     \oraresz{B{\" o}g{\" o}s}{a}{an3}{F.10}{1}%
\\ \oraresz{Varga Bo.}{a}{an2}{F.8}{1}%
\\ \oraresz{Husz{\' a}r}{a}{an1}{II.41}{1}%%
  \end{cellatartalom}}
\ora{\mitrakjonle}{6}{4}{3}
  {\begin{cellatartalom}%
     \oraresz{Barth{\' a}n{\' e}}{{\' e}n}{-}{{\' E}.}{3}%%
  \end{cellatartalom}}
\ora{\mitrakjonle}{6}{4}{4}
  {\begin{cellatartalom}%
     \oraresz{Horv{\' a}th}{fi}{-}{Fiz.}{3}%%
  \end{cellatartalom}}
\ora{\mitrakjonle}{6}{4}{5}
  {\begin{cellatartalom}%
     \oraresz{Jantner}{mt}{hu}{F.8}{1.5}%
\\ \oraresz{Dezs{\H o}}{mt}{mt}{T.2}{1.5}%%
  \end{cellatartalom}}
\ora{\mitrakjonle}{6}{4}{6}
  {\begin{cellatartalom}%
     \oraresz{Teremy}{i}{hu}{F.9}{1.5}%
\\ \oraresz{Adamis}{t{\" o}}{mt}{F.8}{1.5}%%
  \end{cellatartalom}}
\ora{\mitrakjonle}{6}{4}{7}
  {\begin{cellatartalom}%
     \oraresz{S{\' o}skuthyn{\' e}}{k{\' e}}{hu}{K.gy}{3}%%
  \end{cellatartalom}}
\ora{\mitrakjonle}{6}{4}{8}
  {\begin{cellatartalom}%
     \oraresz{B{\" o}g{\" o}s}{a}{fz}{A.1}{3}%%
  \end{cellatartalom}}
\ora{\mitrakjonle}{6}{4}{9}
  {\begin{cellatartalom}%
    %
  \end{cellatartalom}}
\ora{\mitrakjonle}{6}{5}{0}
  {\begin{cellatartalom}%
     \oraresz{P{\' o}k}{r}{mt}{R.}{3}%%
  \end{cellatartalom}}
\ora{\mitrakjonle}{6}{5}{1}
  {\begin{cellatartalom}%
     \oraresz{L{\' a}z{\' a}r}{i}{mt}{F.8}{1.5}%
\\ \oraresz{Jantner}{mt}{hu}{M.}{1.5}%%
  \end{cellatartalom}}
\ora{\mitrakjonle}{6}{5}{2}
  {\begin{cellatartalom}%
     \oraresz{Szendrei}{t{\" o}}{hu}{F.8}{1.5}%
\\ \oraresz{F{\" o}ldesi}{inf}{mt}{Szt.1}{1.5}%%
  \end{cellatartalom}}
\ora{\mitrakjonle}{6}{5}{3}
  {\begin{cellatartalom}%
     \oraresz{Teremy}{i}{hu}{F.8}{1.5}%
\\ \oraresz{Dezs{\H o}}{mt}{mt}{T.2}{1.5}%%
  \end{cellatartalom}}
\ora{\mitrakjonle}{6}{5}{4}
  {\begin{cellatartalom}%
     \oraresz{Teremy}{i}{hu}{F.8}{1.5}%
\\ \oraresz{Dezs{\H o}}{mt}{mt}{T.2}{1.5}%%
  \end{cellatartalom}}
\ora{\mitrakjonle}{6}{5}{5}
  {\begin{cellatartalom}%
     \oraresz{Adamis}{f{\" o}}{-}{F.8}{3}%%
  \end{cellatartalom}}
\ora{\mitrakjonle}{6}{5}{6}
  {\begin{cellatartalom}%
    %
  \end{cellatartalom}}
\ora{\mitrakjonle}{6}{5}{7}
  {\begin{cellatartalom}%
    %
  \end{cellatartalom}}
\ora{\mitrakjonle}{6}{5}{8}
  {\begin{cellatartalom}%
    %
  \end{cellatartalom}}
\ora{\mitrakjonle}{6}{5}{9}
  {\begin{cellatartalom}%
    %
  \end{cellatartalom}}
\ora{\mitrakjonle}{7}{1}{0}
  {\begin{cellatartalom}%
     \oraresz{G{\' a}l}{n{\' e}}{n{\' e}1}{N.}{2}%
\\ \oraresz{L{\' a}z{\' a}r}{i}{r}{F.16}{2}%%
  \end{cellatartalom}}
\ora{\mitrakjonle}{7}{1}{1}
  {\begin{cellatartalom}%
     \oraresz{G{\' a}l}{n{\' e}}{n{\' e}1}{N.}{1.33333333333333}%
\\ \oraresz{P{\' o}k}{r}{r}{R.}{1.33333333333333}%
\\ \oraresz{Balanyi}{r}{r}{R.}{1.33333333333333}%%
  \end{cellatartalom}}
\ora{\mitrakjonle}{7}{1}{2}
  {\begin{cellatartalom}%
     \oraresz{Jantner}{k{\' e}}{2}{K.gy}{2}%
\\ \oraresz{Szab{\' o} A.}{t{\" o}}{n{\' e}}{F.16}{2}%%
  \end{cellatartalom}}
\ora{\mitrakjonle}{7}{1}{3}
  {\begin{cellatartalom}%
     \oraresz{Fekete G.}{mt}{n{\' e}}{F.16}{1}%
\\ \oraresz{B{\" o}g{\" o}s}{a}{an3}{A.3}{1}%
\\ \oraresz{Husz{\' a}r}{a}{an1}{I.30}{1}%
\\ \oraresz{Varga Bo.}{a}{an2}{II.41}{1}%%
  \end{cellatartalom}}
\ora{\mitrakjonle}{7}{1}{4}
  {\begin{cellatartalom}%
     \oraresz{B{\" o}g{\" o}s}{a}{an3}{A.3}{1}%
\\ \oraresz{Husz{\' a}r}{a}{an1}{I.30}{1}%
\\ \oraresz{Varga Bo.}{a}{an2}{II.41}{1}%
\\ \oraresz{Teremy}{ny}{n{\' e}}{F.16}{1}%%
  \end{cellatartalom}}
\ora{\mitrakjonle}{7}{1}{5}
  {\begin{cellatartalom}%
     \oraresz{Darab{\' a}nt}{mt}{r}{F.16}{2}%
\\ \oraresz{F{\" o}ldesi}{inf}{n{\' e}}{Szt.1}{2}%%
  \end{cellatartalom}}
\ora{\mitrakjonle}{7}{1}{6}
  {\begin{cellatartalom}%
     \oraresz{Menyh{\' a}rt}{bi}{-}{B.ea}{4}%%
  \end{cellatartalom}}
\ora{\mitrakjonle}{7}{1}{7}
  {\begin{cellatartalom}%
    %
  \end{cellatartalom}}
\ora{\mitrakjonle}{7}{1}{8}
  {\begin{cellatartalom}%
    %
  \end{cellatartalom}}
\ora{\mitrakjonle}{7}{1}{9}
  {\begin{cellatartalom}%
    %
  \end{cellatartalom}}
\ora{\mitrakjonle}{7}{2}{0}
  {\begin{cellatartalom}%
    %
  \end{cellatartalom}}
\ora{\mitrakjonle}{7}{2}{1}
  {\begin{cellatartalom}%
     \oraresz{Szab{\' o} A.}{t{\" o}}{r}{F.18}{2}%
\\ \oraresz{Teremy}{i}{n{\' e}}{F.16}{2}%%
  \end{cellatartalom}}
\ora{\mitrakjonle}{7}{2}{2}
  {\begin{cellatartalom}%
     \oraresz{F{\" o}ldesi}{inf}{r}{Szt.2}{2}%
\\ \oraresz{Teremy}{i}{n{\' e}}{F.16}{2}%%
  \end{cellatartalom}}
\ora{\mitrakjonle}{7}{2}{3}
  {\begin{cellatartalom}%
     \oraresz{Bard{\' o}czky}{te}{L}{Tt.}{1.33333333333333}%
\\ \oraresz{Heged{\H u}s}{te}{F}{Tt.}{1.33333333333333}%
\\ \oraresz{Nagy Sz.}{te}{L}{Tt.}{1.33333333333333}%%
  \end{cellatartalom}}
\ora{\mitrakjonle}{7}{2}{4}
  {\begin{cellatartalom}%
     \oraresz{Balanyi}{r}{r}{R.}{2}%
\\ \oraresz{Jantner}{k{\' e}}{1}{K.gy}{2}%%
  \end{cellatartalom}}
\ora{\mitrakjonle}{7}{2}{5}
  {\begin{cellatartalom}%
     \oraresz{Darab{\' a}nt}{mt}{r}{F.12}{2}%
\\ \oraresz{J.-Varga}{n{\' e}}{n{\' e}1}{F.16}{2}%%
  \end{cellatartalom}}
\ora{\mitrakjonle}{7}{2}{6}
  {\begin{cellatartalom}%
     \oraresz{J.-Varga}{n{\' e}}{n{\' e}1}{F.16}{2}%
\\ \oraresz{L{\' a}z{\' a}r}{ny}{r}{F.12}{2}%%
  \end{cellatartalom}}
\ora{\mitrakjonle}{7}{2}{7}
  {\begin{cellatartalom}%
    %
  \end{cellatartalom}}
\ora{\mitrakjonle}{7}{2}{8}
  {\begin{cellatartalom}%
    %
  \end{cellatartalom}}
\ora{\mitrakjonle}{7}{2}{9}
  {\begin{cellatartalom}%
    %
  \end{cellatartalom}}
\ora{\mitrakjonle}{7}{3}{0}
  {\begin{cellatartalom}%
     \oraresz{B{\" o}g{\" o}s}{a}{an3}{F.17}{1.33333333333333}%
\\ \oraresz{Husz{\' a}r}{a}{an1}{A.2}{1.33333333333333}%
\\ \oraresz{Varga Bo.}{a}{an2}{A.1}{1.33333333333333}%%
  \end{cellatartalom}}
\ora{\mitrakjonle}{7}{3}{1}
  {\begin{cellatartalom}%
     \oraresz{F{\" o}ldesi}{inf}{n{\' e}}{Szt.1}{1}%
\\ \oraresz{B{\" o}g{\" o}s}{a}{an3}{F.17}{1}%
\\ \oraresz{Husz{\' a}r}{a}{an1}{A.2}{1}%
\\ \oraresz{Varga Bo.}{a}{an2}{A.1}{1}%%
  \end{cellatartalom}}
\ora{\mitrakjonle}{7}{3}{2}
  {\begin{cellatartalom}%
     \oraresz{P{\' e}teri}{fi}{-}{Fiz.}{4}%%
  \end{cellatartalom}}
\ora{\mitrakjonle}{7}{3}{3}
  {\begin{cellatartalom}%
     \oraresz{Fekete G.}{mt}{n{\' e}}{F.16}{2}%
\\ \oraresz{Jantner}{k{\' e}}{2}{K.ea}{2}%%
  \end{cellatartalom}}
\ora{\mitrakjonle}{7}{3}{4}
  {\begin{cellatartalom}%
     \oraresz{Menyh{\' a}rt}{bi}{-}{B.ea}{4}%%
  \end{cellatartalom}}
\ora{\mitrakjonle}{7}{3}{5}
  {\begin{cellatartalom}%
     \oraresz{Schlachter}{f{\" o}}{-}{F.16}{4}%%
  \end{cellatartalom}}
\ora{\mitrakjonle}{7}{3}{6}
  {\begin{cellatartalom}%
     \oraresz{Barth{\' a}n{\' e}}{{\' e}n}{-}{{\' E}.}{4}%%
  \end{cellatartalom}}
\ora{\mitrakjonle}{7}{3}{7}
  {\begin{cellatartalom}%
     \oraresz{Heged{\H u}s}{of}{-}{I.33}{4}%%
  \end{cellatartalom}}
\ora{\mitrakjonle}{7}{3}{8}
  {\begin{cellatartalom}%
    %
  \end{cellatartalom}}
\ora{\mitrakjonle}{7}{3}{9}
  {\begin{cellatartalom}%
    %
  \end{cellatartalom}}
\ora{\mitrakjonle}{7}{4}{0}
  {\begin{cellatartalom}%
    %
  \end{cellatartalom}}
\ora{\mitrakjonle}{7}{4}{1}
  {\begin{cellatartalom}%
     \oraresz{Fekete G.}{mt}{n{\' e}}{F.16}{2}%
\\ \oraresz{Darab{\' a}nt}{mt}{r}{M.}{2}%%
  \end{cellatartalom}}
\ora{\mitrakjonle}{7}{4}{2}
  {\begin{cellatartalom}%
     \oraresz{Balanyi}{r}{n{\' e}}{R.}{1}%
\\ \oraresz{B{\" o}g{\" o}s}{a}{an3}{F.10}{1}%
\\ \oraresz{Varga Bo.}{a}{an2}{F.8}{1}%
\\ \oraresz{Husz{\' a}r}{a}{an1}{II.41}{1}%%
  \end{cellatartalom}}
\ora{\mitrakjonle}{7}{4}{3}
  {\begin{cellatartalom}%
     \oraresz{Balanyi}{r}{r}{R.}{2}%
\\ \oraresz{G{\' a}l}{n{\' e}}{n{\' e}1}{N.}{2}%%
  \end{cellatartalom}}
\ora{\mitrakjonle}{7}{4}{4}
  {\begin{cellatartalom}%
     \oraresz{G{\' a}l}{n{\' e}}{n{\' e}1}{N.}{1.33333333333333}%
\\ \oraresz{P{\' o}k}{r}{r}{R.}{1.33333333333333}%
\\ \oraresz{Balanyi}{r}{r}{R.}{1.33333333333333}%%
  \end{cellatartalom}}
\ora{\mitrakjonle}{7}{4}{5}
  {\begin{cellatartalom}%
     \oraresz{L{\' a}z{\' a}r}{i}{r}{F.16}{2}%
\\ \oraresz{Teremy}{i}{n{\' e}}{F.9}{2}%%
  \end{cellatartalom}}
\ora{\mitrakjonle}{7}{4}{6}
  {\begin{cellatartalom}%
     \oraresz{Schlachter}{f{\" o}}{-}{F.16}{4}%%
  \end{cellatartalom}}
\ora{\mitrakjonle}{7}{4}{7}
  {\begin{cellatartalom}%
    %
  \end{cellatartalom}}
\ora{\mitrakjonle}{7}{4}{8}
  {\begin{cellatartalom}%
     \oraresz{B{\" o}g{\" o}s}{a}{fz}{A.1}{4}%%
  \end{cellatartalom}}
\ora{\mitrakjonle}{7}{4}{9}
  {\begin{cellatartalom}%
    %
  \end{cellatartalom}}
\ora{\mitrakjonle}{7}{5}{0}
  {\begin{cellatartalom}%
     \oraresz{J.-Varga}{n{\' e}}{n{\' e}1}{I.39}{4}%%
  \end{cellatartalom}}
\ora{\mitrakjonle}{7}{5}{1}
  {\begin{cellatartalom}%
     \oraresz{F{\" o}ldesi}{inf}{r}{Szt.2}{2}%
\\ \oraresz{J.-Varga}{n{\' e}}{n{\' e}1}{I.39}{2}%%
  \end{cellatartalom}}
\ora{\mitrakjonle}{7}{5}{2}
  {\begin{cellatartalom}%
     \oraresz{Darab{\' a}nt}{mt}{r}{F.16}{2}%
\\ \oraresz{Szab{\' o} A.}{t{\" o}}{n{\' e}}{F.18}{2}%%
  \end{cellatartalom}}
\ora{\mitrakjonle}{7}{5}{3}
  {\begin{cellatartalom}%
     \oraresz{Jantner}{k{\' e}}{1}{F.12}{2}%
\\ \oraresz{Szab{\' o} A.}{t{\" o}}{r}{F.16}{2}%%
  \end{cellatartalom}}
\ora{\mitrakjonle}{7}{5}{4}
  {\begin{cellatartalom}%
     \oraresz{P{\' e}teri}{fi}{-}{Fiz.}{4}%%
  \end{cellatartalom}}
\ora{\mitrakjonle}{7}{5}{5}
  {\begin{cellatartalom}%
     \oraresz{L{\' a}z{\' a}r}{i}{r}{I.30}{2}%
\\ \oraresz{Fekete G.}{mt}{n{\' e}}{F.16}{2}%%
  \end{cellatartalom}}
\ora{\mitrakjonle}{7}{5}{6}
  {\begin{cellatartalom}%
     \oraresz{Bard{\' o}czky}{te}{L}{Tt.}{1.33333333333333}%
\\ \oraresz{Heged{\H u}s}{te}{F}{Tt.}{1.33333333333333}%
\\ \oraresz{Nagy Sz.}{te}{L}{Tt.}{1.33333333333333}%%
  \end{cellatartalom}}
\ora{\mitrakjonle}{7}{5}{7}
  {\begin{cellatartalom}%
     \oraresz{Bard{\' o}czky}{te}{L}{Tt.}{1.33333333333333}%
\\ \oraresz{Heged{\H u}s}{te}{F}{Tt.}{1.33333333333333}%
\\ \oraresz{Nagy Sz.}{te}{L}{Tt.}{1.33333333333333}%%
  \end{cellatartalom}}
\ora{\mitrakjonle}{7}{5}{8}
  {\begin{cellatartalom}%
    %
  \end{cellatartalom}}
\ora{\mitrakjonle}{7}{5}{9}
  {\begin{cellatartalom}%
    %
  \end{cellatartalom}}
\ora{\mitrakjonle}{8}{1}{0}
  {\begin{cellatartalom}%
     \oraresz{Jantner}{k{\' e}}{1}{K.ea}{6}%%
  \end{cellatartalom}}
\ora{\mitrakjonle}{8}{1}{1}
  {\begin{cellatartalom}%
     \oraresz{Kudella}{fr}{fr}{Fr.}{1}%
\\ \oraresz{Pap}{sp}{sp1}{Sp.}{1}%
\\ \oraresz{Krasznay}{ol}{ol4}{Ol.1}{1}%
\\ \oraresz{M{\' e}sz{\' a}ros}{n{\' e}}{n{\' e}1}{F.16}{1}%
\\ \oraresz{S{\' o}tin{\' e}}{n{\' e}}{n{\' e}2}{F.9}{1}%
\\ \oraresz{Hal{\' a}sz}{n{\' e}}{n{\' e}3}{F.18}{1}%%
  \end{cellatartalom}}
\ora{\mitrakjonle}{8}{1}{2}
  {\begin{cellatartalom}%
     \oraresz{Balanyi}{r}{-}{R.}{6}%%
  \end{cellatartalom}}
\ora{\mitrakjonle}{8}{1}{3}
  {\begin{cellatartalom}%
     \oraresz{Krajny{\' a}k}{mt}{1}{II.42}{3}%
\\ \oraresz{Hanga}{k{\' e}}{2}{K.gy}{3}%%
  \end{cellatartalom}}
\ora{\mitrakjonle}{8}{1}{4}
  {\begin{cellatartalom}%
     \oraresz{V. Bartha}{i}{-}{II.42}{6}%%
  \end{cellatartalom}}
\ora{\mitrakjonle}{8}{1}{5}
  {\begin{cellatartalom}%
     \oraresz{Juh{\' a}sz A.}{a}{an2}{F.8}{1.5}%
\\ \oraresz{Luk{\' a}csi}{a}{an1}{Fr.}{1.5}%
\\ \oraresz{Szabados}{a}{an4}{Sp.}{1.5}%
\\ \oraresz{Csan{\' a}di}{a}{an3}{A.2}{1.5}%%
  \end{cellatartalom}}
\ora{\mitrakjonle}{8}{1}{6}
  {\begin{cellatartalom}%
     \oraresz{Juh{\' a}sz A.}{a}{an2}{F.8}{1.5}%
\\ \oraresz{Luk{\' a}csi}{a}{an1}{Fr.}{1.5}%
\\ \oraresz{Szabados}{a}{an4}{Sp.}{1.5}%
\\ \oraresz{Csan{\' a}di}{a}{an3}{A.2}{1.5}%%
  \end{cellatartalom}}
\ora{\mitrakjonle}{8}{1}{7}
  {\begin{cellatartalom}%
     \oraresz{Lovas}{t{\" o}}{2}{II.40}{6}%%
  \end{cellatartalom}}
\ora{\mitrakjonle}{8}{1}{8}
  {\begin{cellatartalom}%
    %
  \end{cellatartalom}}
\ora{\mitrakjonle}{8}{1}{9}
  {\begin{cellatartalom}%
    %
  \end{cellatartalom}}
\ora{\mitrakjonle}{8}{2}{0}
  {\begin{cellatartalom}%
     \oraresz{Darab{\' a}nt}{mt}{2}{II.42}{6}%%
  \end{cellatartalom}}
\ora{\mitrakjonle}{8}{2}{1}
  {\begin{cellatartalom}%
     \oraresz{Heged{\H u}s}{te}{F}{Tt.}{3}%
\\ \oraresz{Nagy Sz.}{te}{L}{Tt.}{3}%%
  \end{cellatartalom}}
\ora{\mitrakjonle}{8}{2}{2}
  {\begin{cellatartalom}%
     \oraresz{Heged{\H u}s}{te}{F}{Tt.}{3}%
\\ \oraresz{Nagy Sz.}{te}{L}{Tt.}{3}%%
  \end{cellatartalom}}
\ora{\mitrakjonle}{8}{2}{3}
  {\begin{cellatartalom}%
     \oraresz{Horv{\' a}th}{fi}{-}{Fiz.}{6}%%
  \end{cellatartalom}}
\ora{\mitrakjonle}{8}{2}{4}
  {\begin{cellatartalom}%
     \oraresz{Kudella}{fr}{fr}{Fr.}{1}%
\\ \oraresz{S{\' o}tin{\' e}}{n{\' e}}{n{\' e}2}{{\' E}.}{1}%
\\ \oraresz{Hal{\' a}sz}{n{\' e}}{n{\' e}3}{N.}{1}%
\\ \oraresz{Pap}{sp}{sp1}{Sp.}{1}%
\\ \oraresz{M{\' e}sz{\' a}ros}{n{\' e}}{n{\' e}1}{II.48}{1}%
\\ \oraresz{Krasznay}{ol}{ol4}{F.11}{1}%%
  \end{cellatartalom}}
\ora{\mitrakjonle}{8}{2}{5}
  {\begin{cellatartalom}%
     \oraresz{Kudella}{fr}{fr}{Fr.}{1}%
\\ \oraresz{S{\' o}tin{\' e}}{n{\' e}}{n{\' e}2}{{\' E}.}{1}%
\\ \oraresz{Hal{\' a}sz}{n{\' e}}{n{\' e}3}{N.}{1}%
\\ \oraresz{Pap}{sp}{sp1}{Sp.}{1}%
\\ \oraresz{M{\' e}sz{\' a}ros}{n{\' e}}{n{\' e}1}{II.48}{1}%
\\ \oraresz{Krasznay}{ol}{ol4}{F.11}{1}%%
  \end{cellatartalom}}
\ora{\mitrakjonle}{8}{2}{6}
  {\begin{cellatartalom}%
     \oraresz{Cserv{\' a}k}{inf}{1}{Szt.1}{3}%
\\ \oraresz{Hanga}{k{\' e}}{2}{K.gy}{3}%%
  \end{cellatartalom}}
\ora{\mitrakjonle}{8}{2}{7}
  {\begin{cellatartalom}%
     \oraresz{S{\' o}skuthyn{\' e}}{bi}{1}{B.gy}{3}%
\\ \oraresz{Hanga}{k{\' e}}{2}{K.gy}{3}%%
  \end{cellatartalom}}
\ora{\mitrakjonle}{8}{2}{8}
  {\begin{cellatartalom}%
    %
  \end{cellatartalom}}
\ora{\mitrakjonle}{8}{2}{9}
  {\begin{cellatartalom}%
    %
  \end{cellatartalom}}
\ora{\mitrakjonle}{8}{3}{0}
  {\begin{cellatartalom}%
     \oraresz{Darab{\' a}nt}{mt}{2}{II.42}{6}%%
  \end{cellatartalom}}
\ora{\mitrakjonle}{8}{3}{1}
  {\begin{cellatartalom}%
     \oraresz{V. Bartha}{ny}{-}{II.42}{6}%%
  \end{cellatartalom}}
\ora{\mitrakjonle}{8}{3}{2}
  {\begin{cellatartalom}%
     \oraresz{Barth{\' a}n{\' e}}{{\' e}n}{-}{{\' E}.}{6}%%
  \end{cellatartalom}}
\ora{\mitrakjonle}{8}{3}{3}
  {\begin{cellatartalom}%
     \oraresz{Tandory}{f{\" o}}{-}{II.42}{6}%%
  \end{cellatartalom}}
\ora{\mitrakjonle}{8}{3}{4}
  {\begin{cellatartalom}%
     \oraresz{Kudella}{fr}{fr}{Fr.}{1}%
\\ \oraresz{Pap}{sp}{sp1}{Sp.}{1}%
\\ \oraresz{Krasznay}{ol}{ol4}{Ol.2}{1}%
\\ \oraresz{M{\' e}sz{\' a}ros}{n{\' e}}{n{\' e}1}{N.}{1}%
\\ \oraresz{S{\' o}tin{\' e}}{n{\' e}}{n{\' e}2}{{\' E}.}{1}%
\\ \oraresz{Hal{\' a}sz}{n{\' e}}{n{\' e}3}{Szt.3}{1}%%
  \end{cellatartalom}}
\ora{\mitrakjonle}{8}{3}{5}
  {\begin{cellatartalom}%
     \oraresz{S{\' o}skuthyn{\' e}}{bi}{2}{K.ea}{3}%
\\ \oraresz{Jantner}{k{\' e}}{1}{K.gy}{3}%%
  \end{cellatartalom}}
\ora{\mitrakjonle}{8}{3}{6}
  {\begin{cellatartalom}%
     \oraresz{Jantner}{k{\' e}}{1}{K.gy}{3}%
\\ \oraresz{Lovas}{t{\" o}}{2}{II.40}{3}%%
  \end{cellatartalom}}
\ora{\mitrakjonle}{8}{3}{7}
  {\begin{cellatartalom}%
     \oraresz{Krajny{\' a}k}{mt}{1}{II.42}{6}%%
  \end{cellatartalom}}
\ora{\mitrakjonle}{8}{3}{8}
  {\begin{cellatartalom}%
    %
  \end{cellatartalom}}
\ora{\mitrakjonle}{8}{3}{9}
  {\begin{cellatartalom}%
    %
  \end{cellatartalom}}
\ora{\mitrakjonle}{8}{4}{0}
  {\begin{cellatartalom}%
     \oraresz{V. Bartha}{i}{-}{II.42}{6}%%
  \end{cellatartalom}}
\ora{\mitrakjonle}{8}{4}{1}
  {\begin{cellatartalom}%
     \oraresz{Juh{\' a}sz A.}{a}{an2}{A.1}{1.5}%
\\ \oraresz{Luk{\' a}csi}{a}{an1}{I.34}{1.5}%
\\ \oraresz{Szabados}{a}{an4}{A.3}{1.5}%
\\ \oraresz{Csan{\' a}di}{a}{an3}{I.32}{1.5}%%
  \end{cellatartalom}}
\ora{\mitrakjonle}{8}{4}{2}
  {\begin{cellatartalom}%
     \oraresz{Juh{\' a}sz A.}{a}{an2}{A.1}{1.5}%
\\ \oraresz{Luk{\' a}csi}{a}{an1}{I.34}{1.5}%
\\ \oraresz{Szabados}{a}{an4}{A.3}{1.5}%
\\ \oraresz{Csan{\' a}di}{a}{an3}{I.32}{1.5}%%
  \end{cellatartalom}}
\ora{\mitrakjonle}{8}{4}{3}
  {\begin{cellatartalom}%
     \oraresz{Horv{\' a}th}{fi}{-}{Fiz.}{6}%%
  \end{cellatartalom}}
\ora{\mitrakjonle}{8}{4}{4}
  {\begin{cellatartalom}%
     \oraresz{Heged{\H u}s}{te}{F}{Tt.}{3}%
\\ \oraresz{Nagy Sz.}{te}{L}{Tt.}{3}%%
  \end{cellatartalom}}
\ora{\mitrakjonle}{8}{4}{5}
  {\begin{cellatartalom}%
     \oraresz{Lovas}{t{\" o}}{1}{II.40}{3}%
\\ \oraresz{Cserv{\' a}k}{inf}{2}{Szt.1}{3}%%
  \end{cellatartalom}}
\ora{\mitrakjonle}{8}{4}{6}
  {\begin{cellatartalom}%
    %
  \end{cellatartalom}}
\ora{\mitrakjonle}{8}{4}{7}
  {\begin{cellatartalom}%
    %
  \end{cellatartalom}}
\ora{\mitrakjonle}{8}{4}{8}
  {\begin{cellatartalom}%
    %
  \end{cellatartalom}}
\ora{\mitrakjonle}{8}{4}{9}
  {\begin{cellatartalom}%
    %
  \end{cellatartalom}}
\ora{\mitrakjonle}{8}{5}{0}
  {\begin{cellatartalom}%
    %
  \end{cellatartalom}}
\ora{\mitrakjonle}{8}{5}{1}
  {\begin{cellatartalom}%
     \oraresz{S{\' o}skuthyn{\' e}}{bi}{1}{K.gy}{6}%%
  \end{cellatartalom}}
\ora{\mitrakjonle}{8}{5}{2}
  {\begin{cellatartalom}%
     \oraresz{S{\' o}skuthyn{\' e}}{bi}{2}{B.ea}{3}%
\\ \oraresz{Jantner}{k{\' e}}{1}{II.41}{3}%%
  \end{cellatartalom}}
\ora{\mitrakjonle}{8}{5}{3}
  {\begin{cellatartalom}%
     \oraresz{Tandory}{f{\" o}}{-}{II.42}{6}%%
  \end{cellatartalom}}
\ora{\mitrakjonle}{8}{5}{4}
  {\begin{cellatartalom}%
     \oraresz{Krajny{\' a}k}{mt}{1}{II.42}{3}%
\\ \oraresz{Cserv{\' a}k}{inf}{2}{Szt.1}{3}%%
  \end{cellatartalom}}
\ora{\mitrakjonle}{8}{5}{5}
  {\begin{cellatartalom}%
     \oraresz{Lovas}{t{\" o}}{1}{II.40}{3}%
\\ \oraresz{Darab{\' a}nt}{mt}{2}{II.42}{3}%%
  \end{cellatartalom}}
\ora{\mitrakjonle}{8}{5}{6}
  {\begin{cellatartalom}%
     \oraresz{Cserv{\' a}k}{inf}{1}{Szt.1}{3}%
\\ \oraresz{Hanga}{k{\' e}}{2}{K.gy}{3}%%
  \end{cellatartalom}}
\ora{\mitrakjonle}{8}{5}{7}
  {\begin{cellatartalom}%
     \oraresz{V. Bartha}{of}{-}{II.42}{6}%%
  \end{cellatartalom}}
\ora{\mitrakjonle}{8}{5}{8}
  {\begin{cellatartalom}%
    %
  \end{cellatartalom}}
\ora{\mitrakjonle}{8}{5}{9}
  {\begin{cellatartalom}%
    %
  \end{cellatartalom}}
\ora{\mitrakjonle}{9}{1}{0}
  {\begin{cellatartalom}%
    %
  \end{cellatartalom}}
\ora{\mitrakjonle}{9}{1}{1}
  {\begin{cellatartalom}%
     \oraresz{Kudella}{fr}{fr}{Fr.}{1}%
\\ \oraresz{Pap}{sp}{sp1}{Sp.}{1}%
\\ \oraresz{Krasznay}{ol}{ol4}{Ol.1}{1}%
\\ \oraresz{M{\' e}sz{\' a}ros}{n{\' e}}{n{\' e}1}{F.16}{1}%
\\ \oraresz{S{\' o}tin{\' e}}{n{\' e}}{n{\' e}2}{F.9}{1}%
\\ \oraresz{Hal{\' a}sz}{n{\' e}}{n{\' e}3}{F.18}{1}%%
  \end{cellatartalom}}
\ora{\mitrakjonle}{9}{1}{2}
  {\begin{cellatartalom}%
     \oraresz{Tandory}{f{\" o}}{-}{I.34}{6}%%
  \end{cellatartalom}}
\ora{\mitrakjonle}{9}{1}{3}
  {\begin{cellatartalom}%
     \oraresz{Pat{\' o}}{t{\" o}}{szt}{II.51}{3}%
\\ \oraresz{Sch{\H o}n}{mt}{tk}{I.34}{3}%%
  \end{cellatartalom}}
\ora{\mitrakjonle}{9}{1}{4}
  {\begin{cellatartalom}%
     \oraresz{P{\' o}k}{r}{-}{II.51}{6}%%
  \end{cellatartalom}}
\ora{\mitrakjonle}{9}{1}{5}
  {\begin{cellatartalom}%
     \oraresz{Juh{\' a}sz A.}{a}{an2}{F.8}{1.5}%
\\ \oraresz{Luk{\' a}csi}{a}{an1}{Fr.}{1.5}%
\\ \oraresz{Szabados}{a}{an4}{Sp.}{1.5}%
\\ \oraresz{Csan{\' a}di}{a}{an3}{A.2}{1.5}%%
  \end{cellatartalom}}
\ora{\mitrakjonle}{9}{1}{6}
  {\begin{cellatartalom}%
     \oraresz{Juh{\' a}sz A.}{a}{an2}{F.8}{1.5}%
\\ \oraresz{Luk{\' a}csi}{a}{an1}{Fr.}{1.5}%
\\ \oraresz{Szabados}{a}{an4}{Sp.}{1.5}%
\\ \oraresz{Csan{\' a}di}{a}{an3}{A.2}{1.5}%%
  \end{cellatartalom}}
\ora{\mitrakjonle}{9}{1}{7}
  {\begin{cellatartalom}%
     \oraresz{Sch{\H o}n}{mt}{szt}{M.}{6}%%
  \end{cellatartalom}}
\ora{\mitrakjonle}{9}{1}{8}
  {\begin{cellatartalom}%
    %
  \end{cellatartalom}}
\ora{\mitrakjonle}{9}{1}{9}
  {\begin{cellatartalom}%
    %
  \end{cellatartalom}}
\ora{\mitrakjonle}{9}{2}{0}
  {\begin{cellatartalom}%
     \oraresz{Fekete R.}{k{\' e}}{tk}{K.gy}{6}%%
  \end{cellatartalom}}
\ora{\mitrakjonle}{9}{2}{1}
  {\begin{cellatartalom}%
     \oraresz{Oravecz}{i}{-}{I.34}{6}%%
  \end{cellatartalom}}
\ora{\mitrakjonle}{9}{2}{2}
  {\begin{cellatartalom}%
     \oraresz{Furi{\' a}k}{te}{F}{Tt.}{3}%
\\ \oraresz{Vas E.}{te}{L}{Tt.}{3}%%
  \end{cellatartalom}}
\ora{\mitrakjonle}{9}{2}{3}
  {\begin{cellatartalom}%
     \oraresz{Fekete R.}{k{\' e}}{szt}{K.gy}{3}%
\\ \oraresz{Sch{\H o}n}{mt}{tk}{I.34}{3}%%
  \end{cellatartalom}}
\ora{\mitrakjonle}{9}{2}{4}
  {\begin{cellatartalom}%
     \oraresz{Kudella}{fr}{fr}{Fr.}{1}%
\\ \oraresz{S{\' o}tin{\' e}}{n{\' e}}{n{\' e}2}{{\' E}.}{1}%
\\ \oraresz{Hal{\' a}sz}{n{\' e}}{n{\' e}3}{N.}{1}%
\\ \oraresz{Pap}{sp}{sp1}{Sp.}{1}%
\\ \oraresz{M{\' e}sz{\' a}ros}{n{\' e}}{n{\' e}1}{II.48}{1}%
\\ \oraresz{Krasznay}{ol}{ol4}{F.11}{1}%%
  \end{cellatartalom}}
\ora{\mitrakjonle}{9}{2}{5}
  {\begin{cellatartalom}%
     \oraresz{Kudella}{fr}{fr}{Fr.}{1}%
\\ \oraresz{S{\' o}tin{\' e}}{n{\' e}}{n{\' e}2}{{\' E}.}{1}%
\\ \oraresz{Hal{\' a}sz}{n{\' e}}{n{\' e}3}{N.}{1}%
\\ \oraresz{Pap}{sp}{sp1}{Sp.}{1}%
\\ \oraresz{M{\' e}sz{\' a}ros}{n{\' e}}{n{\' e}1}{II.48}{1}%
\\ \oraresz{Krasznay}{ol}{ol4}{F.11}{1}%%
  \end{cellatartalom}}
\ora{\mitrakjonle}{9}{2}{6}
  {\begin{cellatartalom}%
     \oraresz{Moln{\' a}r}{inf}{szt}{Szt.2}{3}%
\\ \oraresz{Papp L.}{m{\' e}}{tk}{Tk.}{3}%%
  \end{cellatartalom}}
\ora{\mitrakjonle}{9}{2}{7}
  {\begin{cellatartalom}%
     \oraresz{Moln{\' a}r}{inf}{szt}{Szt.2}{3}%
\\ \oraresz{Papp L.}{m{\' e}}{tk}{Tk.}{3}%%
  \end{cellatartalom}}
\ora{\mitrakjonle}{9}{2}{8}
  {\begin{cellatartalom}%
     \oraresz{Krajny{\' a}k}{inf}{szak}{Szt.1}{6}%%
  \end{cellatartalom}}
\ora{\mitrakjonle}{9}{2}{9}
  {\begin{cellatartalom}%
    %
  \end{cellatartalom}}
\ora{\mitrakjonle}{9}{3}{0}
  {\begin{cellatartalom}%
    %
  \end{cellatartalom}}
\ora{\mitrakjonle}{9}{3}{1}
  {\begin{cellatartalom}%
     \oraresz{Moln{\' a}r}{inf}{szt}{Szt.2}{3}%
\\ \oraresz{Papp L.}{m{\' e}}{tk}{Tk.}{3}%%
  \end{cellatartalom}}
\ora{\mitrakjonle}{9}{3}{2}
  {\begin{cellatartalom}%
     \oraresz{Sch{\H o}n}{mt}{szt}{M.}{3}%
\\ \oraresz{P{\' o}k}{m{\' e}}{tk}{Tk.}{3}%%
  \end{cellatartalom}}
\ora{\mitrakjonle}{9}{3}{3}
  {\begin{cellatartalom}%
     \oraresz{Barth{\' a}n{\' e}}{{\' e}n}{-}{{\' E}.}{6}%%
  \end{cellatartalom}}
\ora{\mitrakjonle}{9}{3}{4}
  {\begin{cellatartalom}%
     \oraresz{Kudella}{fr}{fr}{Fr.}{1}%
\\ \oraresz{Pap}{sp}{sp1}{Sp.}{1}%
\\ \oraresz{Krasznay}{ol}{ol4}{Ol.2}{1}%
\\ \oraresz{M{\' e}sz{\' a}ros}{n{\' e}}{n{\' e}1}{N.}{1}%
\\ \oraresz{S{\' o}tin{\' e}}{n{\' e}}{n{\' e}2}{{\' E}.}{1}%
\\ \oraresz{Hal{\' a}sz}{n{\' e}}{n{\' e}3}{Szt.3}{1}%%
  \end{cellatartalom}}
\ora{\mitrakjonle}{9}{3}{5}
  {\begin{cellatartalom}%
     \oraresz{Papp L.}{fi}{-}{Fiz.}{6}%%
  \end{cellatartalom}}
\ora{\mitrakjonle}{9}{3}{6}
  {\begin{cellatartalom}%
     \oraresz{Menyh{\' a}rt}{bi}{-}{B.ea}{6}%%
  \end{cellatartalom}}
\ora{\mitrakjonle}{9}{3}{7}
  {\begin{cellatartalom}%
     \oraresz{Oravecz}{ny}{-}{I.34}{6}%%
  \end{cellatartalom}}
\ora{\mitrakjonle}{9}{3}{8}
  {\begin{cellatartalom}%
    %
  \end{cellatartalom}}
\ora{\mitrakjonle}{9}{3}{9}
  {\begin{cellatartalom}%
    %
  \end{cellatartalom}}
\ora{\mitrakjonle}{9}{4}{0}
  {\begin{cellatartalom}%
     \oraresz{Sch{\H o}n}{mt}{tk}{M.}{6}%%
  \end{cellatartalom}}
\ora{\mitrakjonle}{9}{4}{1}
  {\begin{cellatartalom}%
     \oraresz{Juh{\' a}sz A.}{a}{an2}{A.1}{1.5}%
\\ \oraresz{Luk{\' a}csi}{a}{an1}{I.34}{1.5}%
\\ \oraresz{Szabados}{a}{an4}{A.3}{1.5}%
\\ \oraresz{Csan{\' a}di}{a}{an3}{I.32}{1.5}%%
  \end{cellatartalom}}
\ora{\mitrakjonle}{9}{4}{2}
  {\begin{cellatartalom}%
     \oraresz{Juh{\' a}sz A.}{a}{an2}{A.1}{1.5}%
\\ \oraresz{Luk{\' a}csi}{a}{an1}{I.34}{1.5}%
\\ \oraresz{Szabados}{a}{an4}{A.3}{1.5}%
\\ \oraresz{Csan{\' a}di}{a}{an3}{I.32}{1.5}%%
  \end{cellatartalom}}
\ora{\mitrakjonle}{9}{4}{3}
  {\begin{cellatartalom}%
     \oraresz{Pat{\' o}}{t{\" o}}{tk}{II.50}{3}%
\\ \oraresz{Sch{\H o}n}{mt}{szt}{M.}{3}%%
  \end{cellatartalom}}
\ora{\mitrakjonle}{9}{4}{4}
  {\begin{cellatartalom}%
     \oraresz{Furi{\' a}k}{te}{F}{Tt.}{3}%
\\ \oraresz{Vas E.}{te}{L}{Tt.}{3}%%
  \end{cellatartalom}}
\ora{\mitrakjonle}{9}{4}{5}
  {\begin{cellatartalom}%
     \oraresz{Furi{\' a}k}{te}{F}{Tt.}{3}%
\\ \oraresz{Vas E.}{te}{L}{Tt.}{3}%%
  \end{cellatartalom}}
\ora{\mitrakjonle}{9}{4}{6}
  {\begin{cellatartalom}%
     \oraresz{Papp L.}{fi}{-}{I.34}{6}%%
  \end{cellatartalom}}
\ora{\mitrakjonle}{9}{4}{7}
  {\begin{cellatartalom}%
    %
  \end{cellatartalom}}
\ora{\mitrakjonle}{9}{4}{8}
  {\begin{cellatartalom}%
    %
  \end{cellatartalom}}
\ora{\mitrakjonle}{9}{4}{9}
  {\begin{cellatartalom}%
    %
  \end{cellatartalom}}
\ora{\mitrakjonle}{9}{5}{0}
  {\begin{cellatartalom}%
     \oraresz{Fekete R.}{k{\' e}}{szt}{B.ea}{6}%%
  \end{cellatartalom}}
\ora{\mitrakjonle}{9}{5}{1}
  {\begin{cellatartalom}%
     \oraresz{Tandory}{f{\" o}}{-}{I.34}{6}%%
  \end{cellatartalom}}
\ora{\mitrakjonle}{9}{5}{2}
  {\begin{cellatartalom}%
     \oraresz{Pat{\' o}}{t{\" o}}{tk}{I.34}{3}%
\\ \oraresz{Sch{\H o}n}{mt}{szt}{M.}{3}%%
  \end{cellatartalom}}
\ora{\mitrakjonle}{9}{5}{3}
  {\begin{cellatartalom}%
     \oraresz{Pat{\' o}}{t{\" o}}{szt}{I.34}{3}%
\\ \oraresz{Fekete R.}{k{\' e}}{tk}{K.ea}{3}%%
  \end{cellatartalom}}
\ora{\mitrakjonle}{9}{5}{4}
  {\begin{cellatartalom}%
     \oraresz{Oravecz}{i}{-}{I.34}{6}%%
  \end{cellatartalom}}
\ora{\mitrakjonle}{9}{5}{5}
  {\begin{cellatartalom}%
     \oraresz{Menyh{\' a}rt}{bi}{-}{I.34}{6}%%
  \end{cellatartalom}}
\ora{\mitrakjonle}{9}{5}{6}
  {\begin{cellatartalom}%
     \oraresz{Sch{\H o}n}{of}{-}{I.34}{6}%%
  \end{cellatartalom}}
\ora{\mitrakjonle}{9}{5}{7}
  {\begin{cellatartalom}%
    %
  \end{cellatartalom}}
\ora{\mitrakjonle}{9}{5}{8}
  {\begin{cellatartalom}%
    %
  \end{cellatartalom}}
\ora{\mitrakjonle}{9}{5}{9}
  {\begin{cellatartalom}%
    %
  \end{cellatartalom}}
\ora{\mitrakjonle}{10}{1}{0}
  {\begin{cellatartalom}%
     \oraresz{Dezs{\H o}}{mt}{1}{T.2}{6}%%
  \end{cellatartalom}}
\ora{\mitrakjonle}{10}{1}{1}
  {\begin{cellatartalom}%
     \oraresz{Gismondi}{civ}{1}{I.39}{3}%
\\ \oraresz{Bellesi}{civ}{2}{F.10}{3}%%
  \end{cellatartalom}}
\ora{\mitrakjonle}{10}{1}{2}
  {\begin{cellatartalom}%
     \oraresz{Tomor}{ol}{ol1}{F.10}{2}%
\\ \oraresz{Krasznay}{ol}{ol2}{Ol.1}{2}%
\\ \oraresz{Eletto}{ol}{ol3}{Ol.2}{2}%%
  \end{cellatartalom}}
\ora{\mitrakjonle}{10}{1}{3}
  {\begin{cellatartalom}%
     \oraresz{Dezs{\H o}}{mt}{2}{T.2}{3}%
\\ \oraresz{Tomor}{t{\" o}}{1}{L.}{3}%%
  \end{cellatartalom}}
\ora{\mitrakjonle}{10}{1}{4}
  {\begin{cellatartalom}%
     \oraresz{Malatinszky}{k{\' e}}{-}{K.ea}{6}%%
  \end{cellatartalom}}
\ora{\mitrakjonle}{10}{1}{5}
  {\begin{cellatartalom}%
     \oraresz{Gismondi}{mt{\" o}}{1}{I.39}{3}%
\\ \oraresz{M{\" u}ller}{f{\" o}}{1}{L.}{3}%%
  \end{cellatartalom}}
\ora{\mitrakjonle}{10}{1}{6}
  {\begin{cellatartalom}%
     \oraresz{S{\' a}rk{\' a}ny}{fi}{-}{Fiz.}{6}%%
  \end{cellatartalom}}
\ora{\mitrakjonle}{10}{1}{7}
  {\begin{cellatartalom}%
     \oraresz{P.-P{\' a}sztor}{i}{-}{II.50}{6}%%
  \end{cellatartalom}}
\ora{\mitrakjonle}{10}{1}{8}
  {\begin{cellatartalom}%
    %
  \end{cellatartalom}}
\ora{\mitrakjonle}{10}{1}{9}
  {\begin{cellatartalom}%
    %
  \end{cellatartalom}}
\ora{\mitrakjonle}{10}{2}{0}
  {\begin{cellatartalom}%
    %
  \end{cellatartalom}}
\ora{\mitrakjonle}{10}{2}{1}
  {\begin{cellatartalom}%
     \oraresz{Tomor}{t{\" o}}{2}{F.9}{3}%
\\ \oraresz{Sebesty{\' e}n}{mt}{1}{F.10}{3}%%
  \end{cellatartalom}}
\ora{\mitrakjonle}{10}{2}{2}
  {\begin{cellatartalom}%
     \oraresz{Tomor}{ol}{ol1}{L.}{2}%
\\ \oraresz{Krasznay}{ol}{ol2}{Ol.2}{2}%
\\ \oraresz{Eletto}{ol}{ol3}{F.10}{2}%%
  \end{cellatartalom}}
\ora{\mitrakjonle}{10}{2}{3}
  {\begin{cellatartalom}%
     \oraresz{Szabolcsi}{fr}{fr}{Fr.}{1}%
\\ \oraresz{Pap}{sp}{sp2}{Sp.}{1}%
\\ \oraresz{Csuka}{a}{an2}{A.2}{1}%
\\ \oraresz{Nagy Kata}{a}{an3}{F.18}{1}%
\\ \oraresz{Varga Bo.}{a}{an1}{A.3}{1}%
\\ \oraresz{J.-Varga}{n{\' e}}{n{\' e}5}{I.33}{1}%%
  \end{cellatartalom}}
\ora{\mitrakjonle}{10}{2}{4}
  {\begin{cellatartalom}%
     \oraresz{Malatinszky}{bi}{-}{B.ea}{6}%%
  \end{cellatartalom}}
\ora{\mitrakjonle}{10}{2}{5}
  {\begin{cellatartalom}%
     \oraresz{Gismondi}{civ}{1}{I.39}{3}%
\\ \oraresz{M{\" u}ller}{f{\" o}}{2}{L.}{3}%%
  \end{cellatartalom}}
\ora{\mitrakjonle}{10}{2}{6}
  {\begin{cellatartalom}%
     \oraresz{Eletto}{t{\" o}}{1}{F.10}{3}%
\\ \oraresz{Dezs{\H o}}{mt}{2}{T.2}{3}%%
  \end{cellatartalom}}
\ora{\mitrakjonle}{10}{2}{7}
  {\begin{cellatartalom}%
     \oraresz{M{\" u}ller}{f{\" o}}{1}{L.}{6}%%
  \end{cellatartalom}}
\ora{\mitrakjonle}{10}{2}{8}
  {\begin{cellatartalom}%
    %
  \end{cellatartalom}}
\ora{\mitrakjonle}{10}{2}{9}
  {\begin{cellatartalom}%
    %
  \end{cellatartalom}}
\ora{\mitrakjonle}{10}{3}{0}
  {\begin{cellatartalom}%
     \oraresz{Bellesi}{civ}{2}{F.10}{6}%%
  \end{cellatartalom}}
\ora{\mitrakjonle}{10}{3}{1}
  {\begin{cellatartalom}%
     \oraresz{P.-P{\' a}sztor}{i}{-}{F.10}{6}%%
  \end{cellatartalom}}
\ora{\mitrakjonle}{10}{3}{2}
  {\begin{cellatartalom}%
     \oraresz{Sebesty{\' e}n}{mt}{1}{F.10}{3}%
\\ \oraresz{Dezs{\H o}}{mt}{2}{T.2}{3}%%
  \end{cellatartalom}}
\ora{\mitrakjonle}{10}{3}{3}
  {\begin{cellatartalom}%
     \oraresz{Lad{\' a}nyi}{te}{L}{Tt.}{2}%
\\ \oraresz{De{\' a}k}{te}{F}{Tt.}{2}%
\\ \oraresz{Bard{\' o}czky}{te}{L}{Tt.}{2}%%
  \end{cellatartalom}}
\ora{\mitrakjonle}{10}{3}{4}
  {\begin{cellatartalom}%
     \oraresz{Lad{\' a}nyi}{te}{L}{Tt.}{2}%
\\ \oraresz{De{\' a}k}{te}{F}{Tt.}{2}%
\\ \oraresz{Bard{\' o}czky}{te}{L}{Tt.}{2}%%
  \end{cellatartalom}}
\ora{\mitrakjonle}{10}{3}{5}
  {\begin{cellatartalom}%
     \oraresz{M{\" u}ller}{ol}{ol3}{L.}{2}%
\\ \oraresz{Tomor}{ol}{ol1}{F.10}{2}%
\\ \oraresz{Krasznay}{ol}{ol2}{Ol.1}{2}%%
  \end{cellatartalom}}
\ora{\mitrakjonle}{10}{3}{6}
  {\begin{cellatartalom}%
     \oraresz{Gismondi}{mt{\" o}}{1}{I.39}{3}%
\\ \oraresz{M{\" u}ller}{f{\" o}}{2}{L.}{3}%%
  \end{cellatartalom}}
\ora{\mitrakjonle}{10}{3}{7}
  {\begin{cellatartalom}%
    %
  \end{cellatartalom}}
\ora{\mitrakjonle}{10}{3}{8}
  {\begin{cellatartalom}%
    %
  \end{cellatartalom}}
\ora{\mitrakjonle}{10}{3}{9}
  {\begin{cellatartalom}%
    %
  \end{cellatartalom}}
\ora{\mitrakjonle}{10}{4}{0}
  {\begin{cellatartalom}%
     \oraresz{Sebesty{\' e}n}{mt}{1}{I.39}{6}%%
  \end{cellatartalom}}
\ora{\mitrakjonle}{10}{4}{1}
  {\begin{cellatartalom}%
     \oraresz{P.-P{\' a}sztor}{ny}{-}{F.10}{6}%%
  \end{cellatartalom}}
\ora{\mitrakjonle}{10}{4}{2}
  {\begin{cellatartalom}%
     \oraresz{Malatinszky}{bi}{-}{F.12}{6}%%
  \end{cellatartalom}}
\ora{\mitrakjonle}{10}{4}{3}
  {\begin{cellatartalom}%
     \oraresz{Eletto}{t{\" o}}{2}{F.10}{3}%
\\ \oraresz{Tomor}{t{\" o}}{1}{F.8}{3}%%
  \end{cellatartalom}}
\ora{\mitrakjonle}{10}{4}{4}
  {\begin{cellatartalom}%
     \oraresz{M{\" u}ller}{of}{-}{F.10}{6}%%
  \end{cellatartalom}}
\ora{\mitrakjonle}{10}{4}{5}
  {\begin{cellatartalom}%
     \oraresz{Szabolcsi}{fr}{fr}{Fr.}{1}%
\\ \oraresz{Pap}{sp}{sp2}{Sp.}{1}%
\\ \oraresz{Csuka}{a}{an2}{II.51}{1}%
\\ \oraresz{Nagy Kata}{a}{an3}{II.42}{1}%
\\ \oraresz{Varga Bo.}{a}{an1}{Szt.3}{1}%
\\ \oraresz{J.-Varga}{n{\' e}}{n{\' e}5}{F.12}{1}%%
  \end{cellatartalom}}
\ora{\mitrakjonle}{10}{4}{6}
  {\begin{cellatartalom}%
     \oraresz{Malatinszky}{k{\' e}}{-}{K.ea}{6}%%
  \end{cellatartalom}}
\ora{\mitrakjonle}{10}{4}{7}
  {\begin{cellatartalom}%
     \oraresz{Dezs{\H o}}{mt}{2}{T.2}{6}%%
  \end{cellatartalom}}
\ora{\mitrakjonle}{10}{4}{8}
  {\begin{cellatartalom}%
    %
  \end{cellatartalom}}
\ora{\mitrakjonle}{10}{4}{9}
  {\begin{cellatartalom}%
    %
  \end{cellatartalom}}
\ora{\mitrakjonle}{10}{5}{0}
  {\begin{cellatartalom}%
     \oraresz{Tomor}{t{\" o}}{2}{F.10}{6}%%
  \end{cellatartalom}}
\ora{\mitrakjonle}{10}{5}{1}
  {\begin{cellatartalom}%
     \oraresz{Tomor}{ol}{ol1}{Ol.1}{2}%
\\ \oraresz{Krasznay}{ol}{ol2}{L.}{2}%
\\ \oraresz{Eletto}{ol}{ol3}{F.10}{2}%%
  \end{cellatartalom}}
\ora{\mitrakjonle}{10}{5}{2}
  {\begin{cellatartalom}%
     \oraresz{Tomor}{ol}{ol1}{Ol.1}{2}%
\\ \oraresz{Krasznay}{ol}{ol2}{L.}{2}%
\\ \oraresz{Eletto}{ol}{ol3}{F.10}{2}%%
  \end{cellatartalom}}
\ora{\mitrakjonle}{10}{5}{3}
  {\begin{cellatartalom}%
     \oraresz{Lad{\' a}nyi}{te}{L}{Tt.}{2}%
\\ \oraresz{De{\' a}k}{te}{F}{Tt.}{2}%
\\ \oraresz{Bard{\' o}czky}{te}{L}{Tt.}{2}%%
  \end{cellatartalom}}
\ora{\mitrakjonle}{10}{5}{4}
  {\begin{cellatartalom}%
     \oraresz{Nagy Kata}{a}{an3}{A.2}{1}%
\\ \oraresz{Szabolcsi}{fr}{fr}{Fr.}{1}%
\\ \oraresz{Pap}{sp}{sp2}{Sp.}{1}%
\\ \oraresz{Csuka}{a}{an2}{F.17}{1}%
\\ \oraresz{Varga Bo.}{a}{an1}{{\' E}.}{1}%
\\ \oraresz{J.-Varga}{n{\' e}}{n{\' e}5}{F.10}{1}%%
  \end{cellatartalom}}
\ora{\mitrakjonle}{10}{5}{5}
  {\begin{cellatartalom}%
     \oraresz{Nagy Kata}{a}{an3}{A.2}{1}%
\\ \oraresz{Szabolcsi}{fr}{fr}{Fr.}{1}%
\\ \oraresz{Pap}{sp}{sp2}{Sp.}{1}%
\\ \oraresz{Csuka}{a}{an2}{F.17}{1}%
\\ \oraresz{Varga Bo.}{a}{an1}{{\' E}.}{1}%
\\ \oraresz{J.-Varga}{n{\' e}}{n{\' e}5}{F.10}{1}%%
  \end{cellatartalom}}
\ora{\mitrakjonle}{10}{5}{6}
  {\begin{cellatartalom}%
     \oraresz{S{\' a}rk{\' a}ny}{fi}{-}{F.10}{6}%%
  \end{cellatartalom}}
\ora{\mitrakjonle}{10}{5}{7}
  {\begin{cellatartalom}%
    %
  \end{cellatartalom}}
\ora{\mitrakjonle}{10}{5}{8}
  {\begin{cellatartalom}%
    %
  \end{cellatartalom}}
\ora{\mitrakjonle}{10}{5}{9}
  {\begin{cellatartalom}%
    %
  \end{cellatartalom}}
\ora{\mitrakjonle}{11}{1}{0}
  {\begin{cellatartalom}%
    %
  \end{cellatartalom}}
\ora{\mitrakjonle}{11}{1}{1}
  {\begin{cellatartalom}%
     \oraresz{Sz{\' e}kely}{a}{an1}{A.3}{1.25}%
\\ \oraresz{B{\" o}g{\" o}s}{a}{an1}{A.1}{1.25}%
\\ \oraresz{Hajba}{a}{an2}{A.2}{1.25}%
\\ \oraresz{Juh{\' a}sz A.}{a}{an3}{I.31}{1.25}%%
  \end{cellatartalom}}
\ora{\mitrakjonle}{11}{1}{2}
  {\begin{cellatartalom}%
     \oraresz{Sz{\' e}kely}{a}{an1}{A.3}{1.25}%
\\ \oraresz{B{\" o}g{\" o}s}{a}{an1}{A.1}{1.25}%
\\ \oraresz{Hajba}{a}{an2}{A.2}{1.25}%
\\ \oraresz{Juh{\' a}sz A.}{a}{an3}{I.31}{1.25}%%
  \end{cellatartalom}}
\ora{\mitrakjonle}{11}{1}{3}
  {\begin{cellatartalom}%
     \oraresz{Varga Be.}{mt}{mt}{B.ea}{2.5}%
\\ \oraresz{Jantner}{mt}{an}{B.gy}{2.5}%%
  \end{cellatartalom}}
\ora{\mitrakjonle}{11}{1}{4}
  {\begin{cellatartalom}%
     \oraresz{Kudella}{fr}{fr}{Fr.}{1}%
\\ \oraresz{Gianelli}{ol}{ol5}{L.}{1}%
\\ \oraresz{Pap}{sp}{sp2}{Sp.}{1}%
\\ \oraresz{J.-Varga}{n{\' e}}{n{\' e}4}{Szt.1}{1}%
\\ \oraresz{Szab{\' o} A.}{or}{or}{Kt.}{1}%%
  \end{cellatartalom}}
\ora{\mitrakjonle}{11}{1}{5}
  {\begin{cellatartalom}%
     \oraresz{S{\' a}rk{\' a}ny}{fi}{-}{Fiz.}{5}%%
  \end{cellatartalom}}
\ora{\mitrakjonle}{11}{1}{6}
  {\begin{cellatartalom}%
     \oraresz{Tandory}{f{\" o}}{-}{I.31}{5}%%
  \end{cellatartalom}}
\ora{\mitrakjonle}{11}{1}{7}
  {\begin{cellatartalom}%
    %
  \end{cellatartalom}}
\ora{\mitrakjonle}{11}{1}{8}
  {\begin{cellatartalom}%
    %
  \end{cellatartalom}}
\ora{\mitrakjonle}{11}{1}{9}
  {\begin{cellatartalom}%
    %
  \end{cellatartalom}}
\ora{\mitrakjonle}{11}{2}{0}
  {\begin{cellatartalom}%
     \oraresz{Varga Be.}{mt}{mt}{I.31}{5}%%
  \end{cellatartalom}}
\ora{\mitrakjonle}{11}{2}{1}
  {\begin{cellatartalom}%
     \oraresz{Dickerson}{a}{an1}{I.31}{1.25}%
\\ \oraresz{Hajba}{a}{an2}{A.2}{1.25}%
\\ \oraresz{Juh{\' a}sz A.}{a}{an3}{A.3}{1.25}%
\\ \oraresz{Sz{\' e}kely}{a}{an1}{F.11}{1.25}%%
  \end{cellatartalom}}
\ora{\mitrakjonle}{11}{2}{2}
  {\begin{cellatartalom}%
     \oraresz{B{\" u}kk{\" o}si}{ny}{an}{I.31}{2.5}%
\\ \oraresz{P.-P{\' a}sztor}{t{\" o}}{mt}{II.50}{2.5}%%
  \end{cellatartalom}}
\ora{\mitrakjonle}{11}{2}{3}
  {\begin{cellatartalom}%
     \oraresz{B{\" u}kk{\" o}si}{ny}{mt}{I.31}{2.5}%
\\ \oraresz{Jantner}{mt}{an}{M.}{2.5}%%
  \end{cellatartalom}}
\ora{\mitrakjonle}{11}{2}{4}
  {\begin{cellatartalom}%
     \oraresz{S{\' a}rk{\' a}ny}{fi}{-}{I.31}{5}%%
  \end{cellatartalom}}
\ora{\mitrakjonle}{11}{2}{5}
  {\begin{cellatartalom}%
     \oraresz{Balanyi}{r}{-}{R.}{5}%%
  \end{cellatartalom}}
\ora{\mitrakjonle}{11}{2}{6}
  {\begin{cellatartalom}%
     \oraresz{Barth{\' a}n{\' e}}{{\' e}n}{-}{{\' E}.}{5}%%
  \end{cellatartalom}}
\ora{\mitrakjonle}{11}{2}{7}
  {\begin{cellatartalom}%
     \oraresz{Jantner}{k{\' e}}{-}{I.31}{5}%%
  \end{cellatartalom}}
\ora{\mitrakjonle}{11}{2}{8}
  {\begin{cellatartalom}%
    %
  \end{cellatartalom}}
\ora{\mitrakjonle}{11}{2}{9}
  {\begin{cellatartalom}%
    %
  \end{cellatartalom}}
\ora{\mitrakjonle}{11}{3}{0}
  {\begin{cellatartalom}%
    %
  \end{cellatartalom}}
\ora{\mitrakjonle}{11}{3}{1}
  {\begin{cellatartalom}%
     \oraresz{Kudella}{fr}{fr}{Fr.}{1}%
\\ \oraresz{Pap}{sp}{sp2}{Sp.}{1}%
\\ \oraresz{J.-Varga}{n{\' e}}{n{\' e}4}{N.}{1}%
\\ \oraresz{Gianelli}{ol}{ol5}{I.34}{1}%
\\ \oraresz{Szab{\' o} A.}{or}{or}{F.16}{1}%%
  \end{cellatartalom}}
\ora{\mitrakjonle}{11}{3}{2}
  {\begin{cellatartalom}%
     \oraresz{Kudella}{fr}{fr}{Fr.}{1}%
\\ \oraresz{Pap}{sp}{sp2}{Sp.}{1}%
\\ \oraresz{J.-Varga}{n{\' e}}{n{\' e}4}{N.}{1}%
\\ \oraresz{Gianelli}{ol}{ol5}{I.34}{1}%
\\ \oraresz{Szab{\' o} A.}{or}{or}{F.16}{1}%%
  \end{cellatartalom}}
\ora{\mitrakjonle}{11}{3}{3}
  {\begin{cellatartalom}%
     \oraresz{B{\" o}g{\" o}s}{a}{an1}{I.31}{1.25}%
\\ \oraresz{Hajba}{a}{an2}{Szt.3}{1.25}%
\\ \oraresz{Juh{\' a}sz A.}{a}{an3}{A.3}{1.25}%
\\ \oraresz{Sz{\' e}kely}{a}{an1}{II.44}{1.25}%%
  \end{cellatartalom}}
\ora{\mitrakjonle}{11}{3}{4}
  {\begin{cellatartalom}%
     \oraresz{B{\" u}kk{\" o}si}{i}{mt}{I.31}{2.5}%
\\ \oraresz{Jantner}{mt}{an}{F.16}{2.5}%%
  \end{cellatartalom}}
\ora{\mitrakjonle}{11}{3}{5}
  {\begin{cellatartalom}%
     \oraresz{B{\" u}kk{\" o}si}{i}{mt}{I.31}{2.5}%
\\ \oraresz{Teremy}{t{\" o}}{an}{I.32}{2.5}%%
  \end{cellatartalom}}
\ora{\mitrakjonle}{11}{3}{6}
  {\begin{cellatartalom}%
     \oraresz{B{\" u}kk{\" o}si}{i}{an}{I.31}{2.5}%
\\ \oraresz{Varga Be.}{mt}{mt}{I.32}{2.5}%%
  \end{cellatartalom}}
\ora{\mitrakjonle}{11}{3}{7}
  {\begin{cellatartalom}%
     \oraresz{Varga Be.}{mt}{mt}{I.32}{5}%%
  \end{cellatartalom}}
\ora{\mitrakjonle}{11}{3}{8}
  {\begin{cellatartalom}%
    %
  \end{cellatartalom}}
\ora{\mitrakjonle}{11}{3}{9}
  {\begin{cellatartalom}%
    %
  \end{cellatartalom}}
\ora{\mitrakjonle}{11}{4}{0}
  {\begin{cellatartalom}%
     \oraresz{Heged{\H u}s}{te}{F}{Tt.}{2.5}%
\\ \oraresz{Nagy Sz.}{te}{L}{Tt.}{2.5}%%
  \end{cellatartalom}}
\ora{\mitrakjonle}{11}{4}{1}
  {\begin{cellatartalom}%
     \oraresz{Heged{\H u}s}{te}{F}{Tt.}{2.5}%
\\ \oraresz{Nagy Sz.}{te}{L}{Tt.}{2.5}%%
  \end{cellatartalom}}
\ora{\mitrakjonle}{11}{4}{2}
  {\begin{cellatartalom}%
     \oraresz{Tandory}{f{\" o}}{-}{I.31}{5}%%
  \end{cellatartalom}}
\ora{\mitrakjonle}{11}{4}{3}
  {\begin{cellatartalom}%
     \oraresz{Jantner}{k{\' e}}{-}{K.ea}{5}%%
  \end{cellatartalom}}
\ora{\mitrakjonle}{11}{4}{4}
  {\begin{cellatartalom}%
     \oraresz{Kudella}{fr}{fr}{Fr.}{1}%
\\ \oraresz{Gianelli}{ol}{ol5}{Ol.2}{1}%
\\ \oraresz{Pap}{sp}{sp2}{Sp.}{1}%
\\ \oraresz{J.-Varga}{n{\' e}}{n{\' e}4}{Szt.3}{1}%
\\ \oraresz{Szab{\' o} A.}{or}{or}{I.30}{1}%%
  \end{cellatartalom}}
\ora{\mitrakjonle}{11}{4}{5}
  {\begin{cellatartalom}%
     \oraresz{Krasny{\' a}nszki}{bi}{-}{K.ea}{5}%%
  \end{cellatartalom}}
\ora{\mitrakjonle}{11}{4}{6}
  {\begin{cellatartalom}%
     \oraresz{B{\" o}g{\" o}s}{a}{an1}{A.3}{1.25}%
\\ \oraresz{Dickerson}{a}{an1}{A.2}{1.25}%
\\ \oraresz{Hajba}{a}{an2}{F.10}{1.25}%
\\ \oraresz{Juh{\' a}sz A.}{a}{an3}{I.32}{1.25}%%
  \end{cellatartalom}}
\ora{\mitrakjonle}{11}{4}{7}
  {\begin{cellatartalom}%
     \oraresz{B{\" o}g{\" o}s}{a}{an1}{A.3}{2.5}%
\\ \oraresz{Dickerson}{a}{an1}{A.2}{2.5}%%
  \end{cellatartalom}}
\ora{\mitrakjonle}{11}{4}{8}
  {\begin{cellatartalom}%
    %
  \end{cellatartalom}}
\ora{\mitrakjonle}{11}{4}{9}
  {\begin{cellatartalom}%
    %
  \end{cellatartalom}}
\ora{\mitrakjonle}{11}{5}{0}
  {\begin{cellatartalom}%
     \oraresz{B{\" u}kk{\" o}si}{i}{mt}{I.31}{2.5}%
\\ \oraresz{Jantner}{mt}{an}{I.32}{2.5}%%
  \end{cellatartalom}}
\ora{\mitrakjonle}{11}{5}{1}
  {\begin{cellatartalom}%
     \oraresz{B{\" u}kk{\" o}si}{i}{an}{I.31}{2.5}%
\\ \oraresz{Varga Be.}{mt}{mt}{II.42}{2.5}%%
  \end{cellatartalom}}
\ora{\mitrakjonle}{11}{5}{2}
  {\begin{cellatartalom}%
     \oraresz{B{\" u}kk{\" o}si}{i}{an}{I.31}{2.5}%
\\ \oraresz{Varga Be.}{mt}{mt}{II.42}{2.5}%%
  \end{cellatartalom}}
\ora{\mitrakjonle}{11}{5}{3}
  {\begin{cellatartalom}%
     \oraresz{Heged{\H u}s}{te}{F}{Tt.}{2.5}%
\\ \oraresz{Nagy Sz.}{te}{L}{Tt.}{2.5}%%
  \end{cellatartalom}}
\ora{\mitrakjonle}{11}{5}{4}
  {\begin{cellatartalom}%
     \oraresz{B{\" u}kk{\" o}si}{of}{-}{I.31}{5}%%
  \end{cellatartalom}}
\ora{\mitrakjonle}{11}{5}{5}
  {\begin{cellatartalom}%
     \oraresz{P.-P{\' a}sztor}{t{\" o}}{mt}{R.}{2.5}%
\\ \oraresz{Teremy}{t{\" o}}{an}{I.31}{2.5}%%
  \end{cellatartalom}}
\ora{\mitrakjonle}{11}{5}{6}
  {\begin{cellatartalom}%
     \oraresz{Krasny{\' a}nszki}{bi}{-}{B.ea}{5}%%
  \end{cellatartalom}}
\ora{\mitrakjonle}{11}{5}{7}
  {\begin{cellatartalom}%
    %
  \end{cellatartalom}}
\ora{\mitrakjonle}{11}{5}{8}
  {\begin{cellatartalom}%
    %
  \end{cellatartalom}}
\ora{\mitrakjonle}{11}{5}{9}
  {\begin{cellatartalom}%
    %
  \end{cellatartalom}}
\ora{\mitrakjonle}{12}{1}{0}
  {\begin{cellatartalom}%
     \oraresz{P{\' o}k}{r}{r}{R.}{8}%%
  \end{cellatartalom}}
\ora{\mitrakjonle}{12}{1}{1}
  {\begin{cellatartalom}%
     \oraresz{Sz{\' e}kely}{a}{an1}{A.3}{2}%
\\ \oraresz{Hajba}{a}{an2}{A.2}{2}%
\\ \oraresz{Juh{\' a}sz A.}{a}{an3}{I.31}{2}%
\\ \oraresz{J.-Varga}{n{\' e}}{n{\' e}1}{I.30}{2}%%
  \end{cellatartalom}}
\ora{\mitrakjonle}{12}{1}{2}
  {\begin{cellatartalom}%
     \oraresz{Sz{\' e}kely}{a}{an1}{A.3}{2}%
\\ \oraresz{Hajba}{a}{an2}{A.2}{2}%
\\ \oraresz{Juh{\' a}sz A.}{a}{an3}{I.31}{2}%
\\ \oraresz{J.-Varga}{n{\' e}}{n{\' e}1}{I.30}{2}%%
  \end{cellatartalom}}
\ora{\mitrakjonle}{12}{1}{3}
  {\begin{cellatartalom}%
     \oraresz{Darab{\' a}nt}{mt}{n{\' e}}{M.}{4}%
\\ \oraresz{Adamis}{t{\" o}}{r}{II.50}{4}%%
  \end{cellatartalom}}
\ora{\mitrakjonle}{12}{1}{4}
  {\begin{cellatartalom}%
     \oraresz{Kudella}{fr}{fr}{Fr.}{1}%
\\ \oraresz{Gianelli}{ol}{ol5}{L.}{1}%
\\ \oraresz{Pap}{sp}{sp2}{Sp.}{1}%
\\ \oraresz{Tandory}{a}{an3}{I.31}{1}%
\\ \oraresz{T{\' o}th}{a}{an1}{F.12}{1}%
\\ \oraresz{Szabados}{a}{an2}{A.1}{1}%
\\ \oraresz{J.-Varga}{n{\' e}}{n{\' e}4}{Szt.1}{1}%
\\ \oraresz{Szab{\' o} A.}{or}{or}{Kt.}{1}%%
  \end{cellatartalom}}
\ora{\mitrakjonle}{12}{1}{5}
  {\begin{cellatartalom}%
     \oraresz{Balanyi}{r}{n{\' e}1}{R.}{4}%
\\ \oraresz{Schlachter}{mt}{r}{F.18}{4}%%
  \end{cellatartalom}}
\ora{\mitrakjonle}{12}{1}{6}
  {\begin{cellatartalom}%
     \oraresz{P{\' e}teri}{fi}{-}{F.18}{8}%%
  \end{cellatartalom}}
\ora{\mitrakjonle}{12}{1}{7}
  {\begin{cellatartalom}%
     \oraresz{S{\' o}skuthyn{\' e}}{k{\' e}}{-}{F.16}{8}%%
  \end{cellatartalom}}
\ora{\mitrakjonle}{12}{1}{8}
  {\begin{cellatartalom}%
    %
  \end{cellatartalom}}
\ora{\mitrakjonle}{12}{1}{9}
  {\begin{cellatartalom}%
    %
  \end{cellatartalom}}
\ora{\mitrakjonle}{12}{2}{0}
  {\begin{cellatartalom}%
     \oraresz{J.-Varga}{n{\' e}}{n{\' e}1}{II.44}{8}%%
  \end{cellatartalom}}
\ora{\mitrakjonle}{12}{2}{1}
  {\begin{cellatartalom}%
     \oraresz{J.-Varga}{n{\' e}}{n{\' e}1}{II.44}{2}%
\\ \oraresz{Hajba}{a}{an2}{A.2}{2}%
\\ \oraresz{Juh{\' a}sz A.}{a}{an3}{A.3}{2}%
\\ \oraresz{Sz{\' e}kely}{a}{an1}{F.11}{2}%%
  \end{cellatartalom}}
\ora{\mitrakjonle}{12}{2}{2}
  {\begin{cellatartalom}%
     \oraresz{Balanyi}{r}{r}{R.}{4}%
\\ \oraresz{Darab{\' a}nt}{mt}{n{\' e}}{F.18}{4}%%
  \end{cellatartalom}}
\ora{\mitrakjonle}{12}{2}{3}
  {\begin{cellatartalom}%
     \oraresz{S{\' o}skuthyn{\' e}}{bi}{-}{B.ea}{8}%%
  \end{cellatartalom}}
\ora{\mitrakjonle}{12}{2}{4}
  {\begin{cellatartalom}%
     \oraresz{Lad{\' a}nyi}{te}{L}{Tt.}{2.66666666666667}%
\\ \oraresz{De{\' a}k}{te}{F}{Tt.}{2.66666666666667}%
\\ \oraresz{Nagy Sz.}{te}{L}{Tt.}{2.66666666666667}%%
  \end{cellatartalom}}
\ora{\mitrakjonle}{12}{2}{5}
  {\begin{cellatartalom}%
     \oraresz{Lad{\' a}nyi}{te}{L}{Tt.}{2.66666666666667}%
\\ \oraresz{De{\' a}k}{te}{F}{Tt.}{2.66666666666667}%
\\ \oraresz{Nagy Sz.}{te}{L}{Tt.}{2.66666666666667}%%
  \end{cellatartalom}}
\ora{\mitrakjonle}{12}{2}{6}
  {\begin{cellatartalom}%
     \oraresz{V. Bartha}{i}{-}{F.18}{8}%%
  \end{cellatartalom}}
\ora{\mitrakjonle}{12}{2}{7}
  {\begin{cellatartalom}%
     \oraresz{Adamis}{f{\" o}}{-}{F.18}{8}%%
  \end{cellatartalom}}
\ora{\mitrakjonle}{12}{2}{8}
  {\begin{cellatartalom}%
     \oraresz{Krajny{\' a}k}{inf}{szak}{Szt.1}{8}%%
  \end{cellatartalom}}
\ora{\mitrakjonle}{12}{2}{9}
  {\begin{cellatartalom}%
    %
  \end{cellatartalom}}
\ora{\mitrakjonle}{12}{3}{0}
  {\begin{cellatartalom}%
     \oraresz{V. Bartha}{ny}{-}{F.18}{8}%%
  \end{cellatartalom}}
\ora{\mitrakjonle}{12}{3}{1}
  {\begin{cellatartalom}%
     \oraresz{Tandory}{a}{an3}{A.3}{1}%
\\ \oraresz{Kudella}{fr}{fr}{Fr.}{1}%
\\ \oraresz{Pap}{sp}{sp2}{Sp.}{1}%
\\ \oraresz{Szabados}{a}{an2}{I.31}{1}%
\\ \oraresz{T{\' o}th}{a}{an1}{II.43}{1}%
\\ \oraresz{J.-Varga}{n{\' e}}{n{\' e}4}{N.}{1}%
\\ \oraresz{Gianelli}{ol}{ol5}{I.34}{1}%
\\ \oraresz{Szab{\' o} A.}{or}{or}{F.16}{1}%%
  \end{cellatartalom}}
\ora{\mitrakjonle}{12}{3}{2}
  {\begin{cellatartalom}%
     \oraresz{Tandory}{a}{an3}{A.3}{1}%
\\ \oraresz{Kudella}{fr}{fr}{Fr.}{1}%
\\ \oraresz{Pap}{sp}{sp2}{Sp.}{1}%
\\ \oraresz{Szabados}{a}{an2}{I.31}{1}%
\\ \oraresz{T{\' o}th}{a}{an1}{II.43}{1}%
\\ \oraresz{J.-Varga}{n{\' e}}{n{\' e}4}{N.}{1}%
\\ \oraresz{Gianelli}{ol}{ol5}{I.34}{1}%
\\ \oraresz{Szab{\' o} A.}{or}{or}{F.16}{1}%%
  \end{cellatartalom}}
\ora{\mitrakjonle}{12}{3}{3}
  {\begin{cellatartalom}%
     \oraresz{Hajba}{a}{an2}{Szt.3}{2}%
\\ \oraresz{Juh{\' a}sz A.}{a}{an3}{A.3}{2}%
\\ \oraresz{Sz{\' e}kely}{a}{an1}{II.44}{2}%
\\ \oraresz{J.-Varga}{n{\' e}}{n{\' e}1}{Fr.}{2}%%
  \end{cellatartalom}}
\ora{\mitrakjonle}{12}{3}{4}
  {\begin{cellatartalom}%
     \oraresz{S{\' o}skuthyn{\' e}}{k{\' e}}{-}{K.ea}{4}%
\\ \oraresz{Sum}{mid}{mid}{FS.t}{4}%%
  \end{cellatartalom}}
\ora{\mitrakjonle}{12}{3}{5}
  {\begin{cellatartalom}%
     \oraresz{Balanyi}{r}{r}{R.}{4}%
\\ \oraresz{Darab{\' a}nt}{mt}{n{\' e}}{F.18}{4}%%
  \end{cellatartalom}}
\ora{\mitrakjonle}{12}{3}{6}
  {\begin{cellatartalom}%
     \oraresz{Pat{\' o}}{t{\" o}}{n{\' e}}{II.42}{4}%
\\ \oraresz{Schlachter}{mt}{r}{M.}{4}%%
  \end{cellatartalom}}
\ora{\mitrakjonle}{12}{3}{7}
  {\begin{cellatartalom}%
     \oraresz{Schlachter}{mt}{r}{M.}{8}%%
  \end{cellatartalom}}
\ora{\mitrakjonle}{12}{3}{8}
  {\begin{cellatartalom}%
    %
  \end{cellatartalom}}
\ora{\mitrakjonle}{12}{3}{9}
  {\begin{cellatartalom}%
    %
  \end{cellatartalom}}
\ora{\mitrakjonle}{12}{4}{0}
  {\begin{cellatartalom}%
    %
  \end{cellatartalom}}
\ora{\mitrakjonle}{12}{4}{1}
  {\begin{cellatartalom}%
     \oraresz{J.-Varga}{of}{-}{F.18}{8}%%
  \end{cellatartalom}}
\ora{\mitrakjonle}{12}{4}{2}
  {\begin{cellatartalom}%
     \oraresz{Darab{\' a}nt}{mt}{n{\' e}}{M.}{4}%
\\ \oraresz{Schlachter}{mt}{r}{T.2}{4}%%
  \end{cellatartalom}}
\ora{\mitrakjonle}{12}{4}{3}
  {\begin{cellatartalom}%
     \oraresz{Lad{\' a}nyi}{te}{L}{Tt.}{2.66666666666667}%
\\ \oraresz{De{\' a}k}{te}{F}{Tt.}{2.66666666666667}%
\\ \oraresz{Nagy Sz.}{te}{L}{Tt.}{2.66666666666667}%%
  \end{cellatartalom}}
\ora{\mitrakjonle}{12}{4}{4}
  {\begin{cellatartalom}%
     \oraresz{Kudella}{fr}{fr}{Fr.}{1}%
\\ \oraresz{Gianelli}{ol}{ol5}{Ol.2}{1}%
\\ \oraresz{Pap}{sp}{sp2}{Sp.}{1}%
\\ \oraresz{Tandory}{a}{an3}{I.34}{1}%
\\ \oraresz{T{\' o}th}{a}{an1}{A.2}{1}%
\\ \oraresz{Szabados}{a}{an2}{{\' E}.}{1}%
\\ \oraresz{J.-Varga}{n{\' e}}{n{\' e}4}{Szt.3}{1}%
\\ \oraresz{Szab{\' o} A.}{or}{or}{I.30}{1}%%
  \end{cellatartalom}}
\ora{\mitrakjonle}{12}{4}{5}
  {\begin{cellatartalom}%
     \oraresz{Barth{\' a}n{\' e}}{{\' e}n}{-}{{\' E}.}{8}%%
  \end{cellatartalom}}
\ora{\mitrakjonle}{12}{4}{6}
  {\begin{cellatartalom}%
     \oraresz{Dickerson}{a}{an1}{A.2}{2}%
\\ \oraresz{Hajba}{a}{an2}{F.10}{2}%
\\ \oraresz{Juh{\' a}sz A.}{a}{an3}{I.32}{2}%
\\ \oraresz{J.-Varga}{n{\' e}}{n{\' e}1}{R.}{2}%%
  \end{cellatartalom}}
\ora{\mitrakjonle}{12}{4}{7}
  {\begin{cellatartalom}%
     \oraresz{Dickerson}{a}{an1}{A.2}{8}%%
  \end{cellatartalom}}
\ora{\mitrakjonle}{12}{4}{8}
  {\begin{cellatartalom}%
    %
  \end{cellatartalom}}
\ora{\mitrakjonle}{12}{4}{9}
  {\begin{cellatartalom}%
    %
  \end{cellatartalom}}
\ora{\mitrakjonle}{12}{5}{0}
  {\begin{cellatartalom}%
     \oraresz{Adamis}{t{\" o}}{r}{F.18}{8}%%
  \end{cellatartalom}}
\ora{\mitrakjonle}{12}{5}{1}
  {\begin{cellatartalom}%
     \oraresz{Adamis}{f{\" o}}{-}{F.18}{8}%%
  \end{cellatartalom}}
\ora{\mitrakjonle}{12}{5}{2}
  {\begin{cellatartalom}%
     \oraresz{P{\' e}teri}{fi}{-}{Fiz.}{8}%%
  \end{cellatartalom}}
\ora{\mitrakjonle}{12}{5}{3}
  {\begin{cellatartalom}%
     \oraresz{S{\' o}skuthyn{\' e}}{bi}{-}{F.18}{8}%%
  \end{cellatartalom}}
\ora{\mitrakjonle}{12}{5}{4}
  {\begin{cellatartalom}%
     \oraresz{Pat{\' o}}{t{\" o}}{n{\' e}}{II.50}{4}%
\\ \oraresz{P{\' o}k}{r}{r}{R.}{4}%%
  \end{cellatartalom}}
\ora{\mitrakjonle}{12}{5}{5}
  {\begin{cellatartalom}%
     \oraresz{V. Bartha}{i}{-}{F.18}{8}%%
  \end{cellatartalom}}
\ora{\mitrakjonle}{12}{5}{6}
  {\begin{cellatartalom}%
     \oraresz{V. Bartha}{i}{-}{F.18}{8}%%
  \end{cellatartalom}}
\ora{\mitrakjonle}{12}{5}{7}
  {\begin{cellatartalom}%
    %
  \end{cellatartalom}}
\ora{\mitrakjonle}{12}{5}{8}
  {\begin{cellatartalom}%
    %
  \end{cellatartalom}}
\ora{\mitrakjonle}{12}{5}{9}
  {\begin{cellatartalom}%
    %
  \end{cellatartalom}}
\ora{\mitrakjonle}{13}{1}{0}
  {\begin{cellatartalom}%
    %
  \end{cellatartalom}}
\ora{\mitrakjonle}{13}{1}{1}
  {\begin{cellatartalom}%
     \oraresz{Barth{\' a}n{\' e}}{{\' e}n}{-}{{\' E}.}{5}%%
  \end{cellatartalom}}
\ora{\mitrakjonle}{13}{1}{2}
  {\begin{cellatartalom}%
     \oraresz{Kerling}{t{\" o}}{2}{II.50}{2.5}%
\\ \oraresz{Malatinszky}{bi}{1}{B.ea}{2.5}%%
  \end{cellatartalom}}
\ora{\mitrakjonle}{13}{1}{3}
  {\begin{cellatartalom}%
     \oraresz{Bellesi}{ol}{ol4}{Ol.1}{1}%
\\ \oraresz{M{\' e}sz{\' a}ros}{n{\' e}}{n{\' e}3}{Kt.}{1}%
\\ \oraresz{Kudella}{fr}{fr}{K.ea}{1}%
\\ \oraresz{Pap}{sp}{sp1}{I.31}{1}%
\\ \oraresz{G{\' a}l}{n{\' e}}{n{\' e}2}{F.10}{1}%%
  \end{cellatartalom}}
\ora{\mitrakjonle}{13}{1}{4}
  {\begin{cellatartalom}%
     \oraresz{Hajba}{a}{an3}{Szt.3}{1.25}%
\\ \oraresz{Dickerson}{a}{an1}{I.32}{1.25}%
\\ \oraresz{Sz{\' e}kely}{a}{an2}{B.ea}{1.25}%
\\ \oraresz{Csan{\' a}di}{a}{an4}{{\' E}.}{1.25}%%
  \end{cellatartalom}}
\ora{\mitrakjonle}{13}{1}{5}
  {\begin{cellatartalom}%
     \oraresz{Horv{\' a}th G.}{k{\' e}}{1}{K.gy}{2.5}%
\\ \oraresz{Malatinszky}{bi}{2}{B.gy}{2.5}%%
  \end{cellatartalom}}
\ora{\mitrakjonle}{13}{1}{6}
  {\begin{cellatartalom}%
     \oraresz{Horv{\' a}th G.}{k{\' e}}{1}{K.gy}{2.5}%
\\ \oraresz{Malatinszky}{bi}{2}{B.gy}{2.5}%%
  \end{cellatartalom}}
\ora{\mitrakjonle}{13}{1}{7}
  {\begin{cellatartalom}%
     \oraresz{Oravecz}{ny}{-}{II.48}{5}%%
  \end{cellatartalom}}
\ora{\mitrakjonle}{13}{1}{8}
  {\begin{cellatartalom}%
    %
  \end{cellatartalom}}
\ora{\mitrakjonle}{13}{1}{9}
  {\begin{cellatartalom}%
    %
  \end{cellatartalom}}
\ora{\mitrakjonle}{13}{2}{0}
  {\begin{cellatartalom}%
     \oraresz{Heged{\H u}s}{te}{F}{Tt.}{5}%%
  \end{cellatartalom}}
\ora{\mitrakjonle}{13}{2}{1}
  {\begin{cellatartalom}%
     \oraresz{Kerling}{t{\" o}}{2}{II.48}{2.5}%
\\ \oraresz{Horv{\' a}th G.}{k{\' e}}{1}{K.gy}{2.5}%%
  \end{cellatartalom}}
\ora{\mitrakjonle}{13}{2}{2}
  {\begin{cellatartalom}%
     \oraresz{Jantner}{mt}{2}{M.}{2.5}%
\\ \oraresz{Sebesty{\' e}n}{mt}{1}{II.48}{2.5}%%
  \end{cellatartalom}}
\ora{\mitrakjonle}{13}{2}{3}
  {\begin{cellatartalom}%
     \oraresz{Tandory}{f{\" o}}{-}{II.48}{5}%%
  \end{cellatartalom}}
\ora{\mitrakjonle}{13}{2}{4}
  {\begin{cellatartalom}%
     \oraresz{Hajba}{a}{an3}{F.18}{1.25}%
\\ \oraresz{Juh{\' a}sz A.}{a}{an1}{A.2}{1.25}%
\\ \oraresz{Sz{\' e}kely}{a}{an2}{A.1}{1.25}%
\\ \oraresz{Csan{\' a}di}{a}{an4}{A.3}{1.25}%%
  \end{cellatartalom}}
\ora{\mitrakjonle}{13}{2}{5}
  {\begin{cellatartalom}%
     \oraresz{Hajba}{a}{an3}{F.18}{1.25}%
\\ \oraresz{Juh{\' a}sz A.}{a}{an1}{A.2}{1.25}%
\\ \oraresz{Sz{\' e}kely}{a}{an2}{A.1}{1.25}%
\\ \oraresz{Csan{\' a}di}{a}{an4}{A.3}{1.25}%%
  \end{cellatartalom}}
\ora{\mitrakjonle}{13}{2}{6}
  {\begin{cellatartalom}%
     \oraresz{Horv{\' a}th}{fi}{-}{II.48}{5}%%
  \end{cellatartalom}}
\ora{\mitrakjonle}{13}{2}{7}
  {\begin{cellatartalom}%
     \oraresz{Lad{\' a}nyi}{te}{L}{Tt.}{5}%%
  \end{cellatartalom}}
\ora{\mitrakjonle}{13}{2}{8}
  {\begin{cellatartalom}%
    %
  \end{cellatartalom}}
\ora{\mitrakjonle}{13}{2}{9}
  {\begin{cellatartalom}%
    %
  \end{cellatartalom}}
\ora{\mitrakjonle}{13}{3}{0}
  {\begin{cellatartalom}%
    %
  \end{cellatartalom}}
\ora{\mitrakjonle}{13}{3}{1}
  {\begin{cellatartalom}%
     \oraresz{Jantner}{mt}{2}{II.48}{2.5}%
\\ \oraresz{Sebesty{\' e}n}{mt}{1}{M.}{2.5}%%
  \end{cellatartalom}}
\ora{\mitrakjonle}{13}{3}{2}
  {\begin{cellatartalom}%
     \oraresz{Oravecz}{i}{-}{II.48}{5}%%
  \end{cellatartalom}}
\ora{\mitrakjonle}{13}{3}{3}
  {\begin{cellatartalom}%
     \oraresz{Oravecz}{i}{-}{II.48}{5}%%
  \end{cellatartalom}}
\ora{\mitrakjonle}{13}{3}{4}
  {\begin{cellatartalom}%
     \oraresz{Hajba}{a}{an3}{II.44}{1.25}%
\\ \oraresz{Juh{\' a}sz A.}{a}{an1}{II.48}{1.25}%
\\ \oraresz{Sz{\' e}kely}{a}{an2}{K.gy}{1.25}%
\\ \oraresz{Csan{\' a}di}{a}{an4}{A.3}{1.25}%%
  \end{cellatartalom}}
\ora{\mitrakjonle}{13}{3}{5}
  {\begin{cellatartalom}%
     \oraresz{Kerling}{t{\" o}}{1}{II.48}{2.5}%
\\ \oraresz{Malatinszky}{bi}{2}{B.gy}{2.5}%%
  \end{cellatartalom}}
\ora{\mitrakjonle}{13}{3}{6}
  {\begin{cellatartalom}%
     \oraresz{Balanyi}{r}{-}{R.}{5}%%
  \end{cellatartalom}}
\ora{\mitrakjonle}{13}{3}{7}
  {\begin{cellatartalom}%
    %
  \end{cellatartalom}}
\ora{\mitrakjonle}{13}{3}{8}
  {\begin{cellatartalom}%
    %
  \end{cellatartalom}}
\ora{\mitrakjonle}{13}{3}{9}
  {\begin{cellatartalom}%
    %
  \end{cellatartalom}}
\ora{\mitrakjonle}{13}{4}{0}
  {\begin{cellatartalom}%
    %
  \end{cellatartalom}}
\ora{\mitrakjonle}{13}{4}{1}
  {\begin{cellatartalom}%
     \oraresz{Bellesi}{ol}{ol4}{Ol.1}{1}%
\\ \oraresz{Kudella}{fr}{fr}{Fr.}{1}%
\\ \oraresz{G{\' a}l}{n{\' e}}{n{\' e}2}{N.}{1}%
\\ \oraresz{Pap}{sp}{sp1}{Sp.}{1}%
\\ \oraresz{M{\' e}sz{\' a}ros}{n{\' e}}{n{\' e}3}{{\' E}.}{1}%%
  \end{cellatartalom}}
\ora{\mitrakjonle}{13}{4}{2}
  {\begin{cellatartalom}%
     \oraresz{Heged{\" u}s}{ol}{ol4}{Ol.1}{1}%
\\ \oraresz{Kudella}{fr}{fr}{Fr.}{1}%
\\ \oraresz{G{\' a}l}{n{\' e}}{n{\' e}2}{N.}{1}%
\\ \oraresz{Pap}{sp}{sp1}{Sp.}{1}%
\\ \oraresz{M{\' e}sz{\' a}ros}{n{\' e}}{n{\' e}3}{{\' E}.}{1}%%
  \end{cellatartalom}}
\ora{\mitrakjonle}{13}{4}{3}
  {\begin{cellatartalom}%
     \oraresz{Horv{\' a}th G.}{k{\' e}}{2}{K.gy}{2.5}%
\\ \oraresz{Malatinszky}{bi}{1}{B.gy}{2.5}%%
  \end{cellatartalom}}
\ora{\mitrakjonle}{13}{4}{4}
  {\begin{cellatartalom}%
     \oraresz{Horv{\' a}th G.}{k{\' e}}{2}{K.gy}{2.5}%
\\ \oraresz{Malatinszky}{bi}{1}{B.gy}{2.5}%%
  \end{cellatartalom}}
\ora{\mitrakjonle}{13}{4}{5}
  {\begin{cellatartalom}%
     \oraresz{Hajba}{a}{an3}{A.2}{1.25}%
\\ \oraresz{Juh{\' a}sz A.}{a}{an1}{II.48}{1.25}%
\\ \oraresz{Sz{\' e}kely}{a}{an2}{R.}{1.25}%
\\ \oraresz{Csan{\' a}di}{a}{an4}{A.3}{1.25}%%
  \end{cellatartalom}}
\ora{\mitrakjonle}{13}{4}{6}
  {\begin{cellatartalom}%
     \oraresz{Tandory}{f{\" o}}{-}{II.48}{5}%%
  \end{cellatartalom}}
\ora{\mitrakjonle}{13}{4}{7}
  {\begin{cellatartalom}%
     \oraresz{Horv{\' a}th}{fi}{-}{Fiz.}{5}%%
  \end{cellatartalom}}
\ora{\mitrakjonle}{13}{4}{8}
  {\begin{cellatartalom}%
    %
  \end{cellatartalom}}
\ora{\mitrakjonle}{13}{4}{9}
  {\begin{cellatartalom}%
    %
  \end{cellatartalom}}
\ora{\mitrakjonle}{13}{5}{0}
  {\begin{cellatartalom}%
     \oraresz{Heged{\H u}s}{te}{F}{Tt.}{2.5}%
\\ \oraresz{Lad{\' a}nyi}{te}{L}{Tt.}{2.5}%%
  \end{cellatartalom}}
\ora{\mitrakjonle}{13}{5}{1}
  {\begin{cellatartalom}%
     \oraresz{Heged{\H u}s}{te}{F}{Tt.}{2.5}%
\\ \oraresz{Lad{\' a}nyi}{te}{L}{Tt.}{2.5}%%
  \end{cellatartalom}}
\ora{\mitrakjonle}{13}{5}{2}
  {\begin{cellatartalom}%
     \oraresz{Oravecz}{i}{-}{II.48}{5}%%
  \end{cellatartalom}}
\ora{\mitrakjonle}{13}{5}{3}
  {\begin{cellatartalom}%
     \oraresz{Bellesi}{ol}{ol4}{L.}{1}%
\\ \oraresz{M{\' e}sz{\' a}ros}{n{\' e}}{n{\' e}3}{N.}{1}%
\\ \oraresz{Kudella}{fr}{fr}{Fr.}{1}%
\\ \oraresz{Pap}{sp}{sp1}{Sp.}{1}%
\\ \oraresz{G{\' a}l}{n{\' e}}{n{\' e}2}{{\' E}.}{1}%%
  \end{cellatartalom}}
\ora{\mitrakjonle}{13}{5}{4}
  {\begin{cellatartalom}%
     \oraresz{Jantner}{mt}{2}{II.43}{2.5}%
\\ \oraresz{Sebesty{\' e}n}{mt}{1}{II.48}{2.5}%%
  \end{cellatartalom}}
\ora{\mitrakjonle}{13}{5}{5}
  {\begin{cellatartalom}%
     \oraresz{Kerling}{t{\" o}}{1}{II.50}{2.5}%
\\ \oraresz{Horv{\' a}th G.}{k{\' e}}{2}{K.ea}{2.5}%%
  \end{cellatartalom}}
\ora{\mitrakjonle}{13}{5}{6}
  {\begin{cellatartalom}%
     \oraresz{Malatinszky}{of}{-}{II.48}{5}%%
  \end{cellatartalom}}
\ora{\mitrakjonle}{13}{5}{7}
  {\begin{cellatartalom}%
    %
  \end{cellatartalom}}
\ora{\mitrakjonle}{13}{5}{8}
  {\begin{cellatartalom}%
    %
  \end{cellatartalom}}
\ora{\mitrakjonle}{13}{5}{9}
  {\begin{cellatartalom}%
    %
  \end{cellatartalom}}
\ora{\mitrakjonle}{14}{1}{0}
  {\begin{cellatartalom}%
    %
  \end{cellatartalom}}
\ora{\mitrakjonle}{14}{1}{1}
  {\begin{cellatartalom}%
     \oraresz{Horv{\' a}th}{fi}{-}{II.51}{5}%%
  \end{cellatartalom}}
\ora{\mitrakjonle}{14}{1}{2}
  {\begin{cellatartalom}%
     \oraresz{Varga Be.}{mt}{tk}{II.51}{2.5}%
\\ \oraresz{Darab{\' a}nt}{mt}{szt}{II.48}{2.5}%%
  \end{cellatartalom}}
\ora{\mitrakjonle}{14}{1}{3}
  {\begin{cellatartalom}%
     \oraresz{Bellesi}{ol}{ol4}{Ol.1}{1}%
\\ \oraresz{M{\' e}sz{\' a}ros}{n{\' e}}{n{\' e}3}{Kt.}{1}%
\\ \oraresz{Kudella}{fr}{fr}{K.ea}{1}%
\\ \oraresz{Pap}{sp}{sp1}{I.31}{1}%
\\ \oraresz{G{\' a}l}{n{\' e}}{n{\' e}2}{F.10}{1}%%
  \end{cellatartalom}}
\ora{\mitrakjonle}{14}{1}{4}
  {\begin{cellatartalom}%
     \oraresz{Hajba}{a}{an3}{Szt.3}{1.25}%
\\ \oraresz{Dickerson}{a}{an1}{I.32}{1.25}%
\\ \oraresz{Sz{\' e}kely}{a}{an2}{B.ea}{1.25}%
\\ \oraresz{Csan{\' a}di}{a}{an4}{{\' E}.}{1.25}%%
  \end{cellatartalom}}
\ora{\mitrakjonle}{14}{1}{5}
  {\begin{cellatartalom}%
     \oraresz{Moln{\' a}r}{inf}{szt}{Szt.2}{2.5}%
\\ \oraresz{Pat{\' o}}{t{\" o}}{tk}{II.51}{2.5}%%
  \end{cellatartalom}}
\ora{\mitrakjonle}{14}{1}{6}
  {\begin{cellatartalom}%
     \oraresz{Pat{\' o}}{i}{-}{II.51}{5}%%
  \end{cellatartalom}}
\ora{\mitrakjonle}{14}{1}{7}
  {\begin{cellatartalom}%
     \oraresz{Tandory}{f{\" o}}{-}{II.44}{5}%%
  \end{cellatartalom}}
\ora{\mitrakjonle}{14}{1}{8}
  {\begin{cellatartalom}%
    %
  \end{cellatartalom}}
\ora{\mitrakjonle}{14}{1}{9}
  {\begin{cellatartalom}%
    %
  \end{cellatartalom}}
\ora{\mitrakjonle}{14}{2}{0}
  {\begin{cellatartalom}%
     \oraresz{Moln{\' a}r}{inf}{szt}{Szt.2}{5}%%
  \end{cellatartalom}}
\ora{\mitrakjonle}{14}{2}{1}
  {\begin{cellatartalom}%
     \oraresz{Fekete R.}{bi}{-}{B.ea}{5}%%
  \end{cellatartalom}}
\ora{\mitrakjonle}{14}{2}{2}
  {\begin{cellatartalom}%
     \oraresz{Fekete R.}{k{\' e}}{-}{K.ea}{5}%%
  \end{cellatartalom}}
\ora{\mitrakjonle}{14}{2}{3}
  {\begin{cellatartalom}%
     \oraresz{Varga Be.}{mt}{tk}{II.51}{2.5}%
\\ \oraresz{Pat{\' o}}{t{\" o}}{szt}{II.50}{2.5}%%
  \end{cellatartalom}}
\ora{\mitrakjonle}{14}{2}{4}
  {\begin{cellatartalom}%
     \oraresz{Hajba}{a}{an3}{F.18}{1.25}%
\\ \oraresz{Juh{\' a}sz A.}{a}{an1}{A.2}{1.25}%
\\ \oraresz{Sz{\' e}kely}{a}{an2}{A.1}{1.25}%
\\ \oraresz{Csan{\' a}di}{a}{an4}{A.3}{1.25}%%
  \end{cellatartalom}}
\ora{\mitrakjonle}{14}{2}{5}
  {\begin{cellatartalom}%
     \oraresz{Hajba}{a}{an3}{F.18}{1.25}%
\\ \oraresz{Juh{\' a}sz A.}{a}{an1}{A.2}{1.25}%
\\ \oraresz{Sz{\' e}kely}{a}{an2}{A.1}{1.25}%
\\ \oraresz{Csan{\' a}di}{a}{an4}{A.3}{1.25}%%
  \end{cellatartalom}}
\ora{\mitrakjonle}{14}{2}{6}
  {\begin{cellatartalom}%
     \oraresz{Heged{\H u}s}{te}{F}{Tt.}{2.5}%
\\ \oraresz{Bard{\' o}czky}{te}{L}{Tt.}{2.5}%%
  \end{cellatartalom}}
\ora{\mitrakjonle}{14}{2}{7}
  {\begin{cellatartalom}%
     \oraresz{Heged{\H u}s}{te}{F}{Tt.}{2.5}%
\\ \oraresz{Bard{\' o}czky}{te}{L}{Tt.}{2.5}%%
  \end{cellatartalom}}
\ora{\mitrakjonle}{14}{2}{8}
  {\begin{cellatartalom}%
     \oraresz{Pat{\' o}}{ny}{-}{II.51}{5}%%
  \end{cellatartalom}}
\ora{\mitrakjonle}{14}{2}{9}
  {\begin{cellatartalom}%
    %
  \end{cellatartalom}}
\ora{\mitrakjonle}{14}{3}{0}
  {\begin{cellatartalom}%
     \oraresz{Heged{\H u}s}{te}{F}{Tt.}{2.5}%
\\ \oraresz{Bard{\' o}czky}{te}{L}{Tt.}{2.5}%%
  \end{cellatartalom}}
\ora{\mitrakjonle}{14}{3}{1}
  {\begin{cellatartalom}%
     \oraresz{Pat{\' o}}{ny}{-}{II.51}{5}%%
  \end{cellatartalom}}
\ora{\mitrakjonle}{14}{3}{2}
  {\begin{cellatartalom}%
     \oraresz{Varga Be.}{mt}{tk}{II.42}{2.5}%
\\ \oraresz{Pat{\' o}}{t{\" o}}{szt}{II.51}{2.5}%%
  \end{cellatartalom}}
\ora{\mitrakjonle}{14}{3}{3}
  {\begin{cellatartalom}%
     \oraresz{Horv{\' a}th}{fi}{-}{Fiz.}{5}%%
  \end{cellatartalom}}
\ora{\mitrakjonle}{14}{3}{4}
  {\begin{cellatartalom}%
     \oraresz{Hajba}{a}{an3}{II.44}{1.25}%
\\ \oraresz{Juh{\' a}sz A.}{a}{an1}{II.48}{1.25}%
\\ \oraresz{Sz{\' e}kely}{a}{an2}{K.gy}{1.25}%
\\ \oraresz{Csan{\' a}di}{a}{an4}{A.3}{1.25}%%
  \end{cellatartalom}}
\ora{\mitrakjonle}{14}{3}{5}
  {\begin{cellatartalom}%
     \oraresz{Tandory}{f{\" o}}{-}{I.34}{5}%%
  \end{cellatartalom}}
\ora{\mitrakjonle}{14}{3}{6}
  {\begin{cellatartalom}%
     \oraresz{Papp L.}{m{\' e}}{tk}{Tk.}{2.5}%
\\ \oraresz{Darab{\' a}nt}{mt}{szt}{I.30}{2.5}%%
  \end{cellatartalom}}
\ora{\mitrakjonle}{14}{3}{7}
  {\begin{cellatartalom}%
     \oraresz{Papp L.}{m{\' e}}{tk}{Tk.}{5}%%
  \end{cellatartalom}}
\ora{\mitrakjonle}{14}{3}{8}
  {\begin{cellatartalom}%
    %
  \end{cellatartalom}}
\ora{\mitrakjonle}{14}{3}{9}
  {\begin{cellatartalom}%
    %
  \end{cellatartalom}}
\ora{\mitrakjonle}{14}{4}{0}
  {\begin{cellatartalom}%
    %
  \end{cellatartalom}}
\ora{\mitrakjonle}{14}{4}{1}
  {\begin{cellatartalom}%
     \oraresz{Bellesi}{ol}{ol4}{Ol.1}{1}%
\\ \oraresz{Kudella}{fr}{fr}{Fr.}{1}%
\\ \oraresz{G{\' a}l}{n{\' e}}{n{\' e}2}{N.}{1}%
\\ \oraresz{Pap}{sp}{sp1}{Sp.}{1}%
\\ \oraresz{M{\' e}sz{\' a}ros}{n{\' e}}{n{\' e}3}{{\' E}.}{1}%%
  \end{cellatartalom}}
\ora{\mitrakjonle}{14}{4}{2}
  {\begin{cellatartalom}%
     \oraresz{Heged{\" u}s}{ol}{ol4}{Ol.1}{1}%
\\ \oraresz{Kudella}{fr}{fr}{Fr.}{1}%
\\ \oraresz{G{\' a}l}{n{\' e}}{n{\' e}2}{N.}{1}%
\\ \oraresz{Pap}{sp}{sp1}{Sp.}{1}%
\\ \oraresz{M{\' e}sz{\' a}ros}{n{\' e}}{n{\' e}3}{{\' E}.}{1}%%
  \end{cellatartalom}}
\ora{\mitrakjonle}{14}{4}{3}
  {\begin{cellatartalom}%
     \oraresz{Papp L.}{m{\' e}}{tk}{Tk.}{2.5}%
\\ \oraresz{Moln{\' a}r}{inf}{szt}{Szt.2}{2.5}%%
  \end{cellatartalom}}
\ora{\mitrakjonle}{14}{4}{4}
  {\begin{cellatartalom}%
     \oraresz{Papp L.}{m{\' e}}{tk}{Tk.}{2.5}%
\\ \oraresz{Moln{\' a}r}{inf}{szt}{Szt.2}{2.5}%%
  \end{cellatartalom}}
\ora{\mitrakjonle}{14}{4}{5}
  {\begin{cellatartalom}%
     \oraresz{Hajba}{a}{an3}{A.2}{1.25}%
\\ \oraresz{Juh{\' a}sz A.}{a}{an1}{II.48}{1.25}%
\\ \oraresz{Sz{\' e}kely}{a}{an2}{R.}{1.25}%
\\ \oraresz{Csan{\' a}di}{a}{an4}{A.3}{1.25}%%
  \end{cellatartalom}}
\ora{\mitrakjonle}{14}{4}{6}
  {\begin{cellatartalom}%
     \oraresz{Pat{\' o}}{t{\" o}}{tk}{II.51}{2.5}%
\\ \oraresz{Darab{\' a}nt}{mt}{szt}{II.44}{2.5}%%
  \end{cellatartalom}}
\ora{\mitrakjonle}{14}{4}{7}
  {\begin{cellatartalom}%
     \oraresz{Pat{\' o}}{of}{-}{II.51}{5}%%
  \end{cellatartalom}}
\ora{\mitrakjonle}{14}{4}{8}
  {\begin{cellatartalom}%
    %
  \end{cellatartalom}}
\ora{\mitrakjonle}{14}{4}{9}
  {\begin{cellatartalom}%
    %
  \end{cellatartalom}}
\ora{\mitrakjonle}{14}{5}{0}
  {\begin{cellatartalom}%
    %
  \end{cellatartalom}}
\ora{\mitrakjonle}{14}{5}{1}
  {\begin{cellatartalom}%
     \oraresz{Fekete R.}{k{\' e}}{-}{B.ea}{5}%%
  \end{cellatartalom}}
\ora{\mitrakjonle}{14}{5}{2}
  {\begin{cellatartalom}%
     \oraresz{Fekete R.}{bi}{-}{II.51}{5}%%
  \end{cellatartalom}}
\ora{\mitrakjonle}{14}{5}{3}
  {\begin{cellatartalom}%
     \oraresz{Bellesi}{ol}{ol4}{L.}{1}%
\\ \oraresz{M{\' e}sz{\' a}ros}{n{\' e}}{n{\' e}3}{N.}{1}%
\\ \oraresz{Kudella}{fr}{fr}{Fr.}{1}%
\\ \oraresz{Pap}{sp}{sp1}{Sp.}{1}%
\\ \oraresz{G{\' a}l}{n{\' e}}{n{\' e}2}{{\' E}.}{1}%%
  \end{cellatartalom}}
\ora{\mitrakjonle}{14}{5}{4}
  {\begin{cellatartalom}%
     \oraresz{Varga Be.}{mt}{tk}{F.16}{2.5}%
\\ \oraresz{Darab{\' a}nt}{mt}{szt}{I.32}{2.5}%%
  \end{cellatartalom}}
\ora{\mitrakjonle}{14}{5}{5}
  {\begin{cellatartalom}%
     \oraresz{Pat{\' o}}{i}{-}{II.51}{5}%%
  \end{cellatartalom}}
\ora{\mitrakjonle}{14}{5}{6}
  {\begin{cellatartalom}%
     \oraresz{Pat{\' o}}{i}{-}{II.51}{5}%%
  \end{cellatartalom}}
\ora{\mitrakjonle}{14}{5}{7}
  {\begin{cellatartalom}%
    %
  \end{cellatartalom}}
\ora{\mitrakjonle}{14}{5}{8}
  {\begin{cellatartalom}%
    %
  \end{cellatartalom}}
\ora{\mitrakjonle}{14}{5}{9}
  {\begin{cellatartalom}%
    %
  \end{cellatartalom}}
\ora{\mitrakjonle}{15}{1}{0}
  {\begin{cellatartalom}%
     \oraresz{Nagy Sz.}{te}{L}{Tt.}{1.33333333333333}%
\\ \oraresz{Furi{\' a}k}{te}{F}{Tt.}{1.33333333333333}%
\\ \oraresz{De{\' a}k}{te}{L}{Tt.}{1.33333333333333}%%
  \end{cellatartalom}}
\ora{\mitrakjonle}{15}{1}{1}
  {\begin{cellatartalom}%
     \oraresz{Nagy Sz.}{te}{L}{Tt.}{1.33333333333333}%
\\ \oraresz{Furi{\' a}k}{te}{F}{Tt.}{1.33333333333333}%
\\ \oraresz{De{\' a}k}{te}{L}{Tt.}{1.33333333333333}%%
  \end{cellatartalom}}
\ora{\mitrakjonle}{15}{1}{2}
  {\begin{cellatartalom}%
     \oraresz{Sch{\H o}n}{fi}{-}{Fiz.}{4}%%
  \end{cellatartalom}}
\ora{\mitrakjonle}{15}{1}{3}
  {\begin{cellatartalom}%
     \oraresz{Krasznay}{of}{-}{I.39}{2}%
\\ \oraresz{S{\" u}le}{i}{F}{F.8}{2}%%
  \end{cellatartalom}}
\ora{\mitrakjonle}{15}{1}{4}
  {\begin{cellatartalom}%
     \oraresz{Eletto}{t{\" o}}{A}{F.10}{1.33333333333333}%
\\ \oraresz{Lovas}{t{\" o}}{F}{II.40}{1.33333333333333}%
\\ \oraresz{S{\" u}le}{ny}{F}{F.8}{1.33333333333333}%%
  \end{cellatartalom}}
\ora{\mitrakjonle}{15}{1}{5}
  {\begin{cellatartalom}%
     \oraresz{Sch{\H o}n}{mt}{A2}{M.}{1.33333333333333}%
\\ \oraresz{Fekete G.}{mt}{A1}{F.9}{1.33333333333333}%
\\ \oraresz{Dezs{\H o}}{mt}{F}{T.2}{1.33333333333333}%%
  \end{cellatartalom}}
\ora{\mitrakjonle}{15}{1}{6}
  {\begin{cellatartalom}%
     \oraresz{Krasznay}{ol}{A1}{F.9}{2}%
\\ \oraresz{Eletto}{civ}{2}{R.}{2}%%
  \end{cellatartalom}}
\ora{\mitrakjonle}{15}{1}{7}
  {\begin{cellatartalom}%
     \oraresz{Dezs{\H o}}{mt}{A2}{T.2}{1.33333333333333}%
\\ \oraresz{Krasznay}{ol}{A1}{F.9}{1.33333333333333}%
\\ \oraresz{Juh{\' a}sz A.}{a}{oktv}{A.1}{1.33333333333333}%%
  \end{cellatartalom}}
\ora{\mitrakjonle}{15}{1}{8}
  {\begin{cellatartalom}%
     \oraresz{Juh{\' a}sz A.}{a}{oktv}{A.1}{4}%%
  \end{cellatartalom}}
\ora{\mitrakjonle}{15}{1}{9}
  {\begin{cellatartalom}%
    %
  \end{cellatartalom}}
\ora{\mitrakjonle}{15}{2}{0}
  {\begin{cellatartalom}%
     \oraresz{S{\" u}le}{i}{F}{F.12}{1.33333333333333}%
\\ \oraresz{Lovas}{t{\" o}}{F}{II.40}{1.33333333333333}%
\\ \oraresz{Menyh{\' a}rt}{bi}{F}{B.gy}{1.33333333333333}%%
  \end{cellatartalom}}
\ora{\mitrakjonle}{15}{2}{1}
  {\begin{cellatartalom}%
     \oraresz{S{\" u}le}{i}{F}{F.12}{1.33333333333333}%
\\ \oraresz{Lovas}{t{\" o}}{F}{II.40}{1.33333333333333}%
\\ \oraresz{Menyh{\' a}rt}{bi}{F}{B.gy}{1.33333333333333}%%
  \end{cellatartalom}}
\ora{\mitrakjonle}{15}{2}{2}
  {\begin{cellatartalom}%
     \oraresz{Krasny{\' a}nszki}{bi}{-}{B.ea}{4}%%
  \end{cellatartalom}}
\ora{\mitrakjonle}{15}{2}{3}
  {\begin{cellatartalom}%
     \oraresz{Gismondi}{mt{\" o}}{2}{I.39}{2}%
\\ \oraresz{Krasznay}{ol}{A1}{Ol.1}{2}%%
  \end{cellatartalom}}
\ora{\mitrakjonle}{15}{2}{4}
  {\begin{cellatartalom}%
     \oraresz{S{\" u}le}{i}{A}{F.10}{4}%%
  \end{cellatartalom}}
\ora{\mitrakjonle}{15}{2}{5}
  {\begin{cellatartalom}%
     \oraresz{S{\" u}le}{i}{A}{F.10}{4}%%
  \end{cellatartalom}}
\ora{\mitrakjonle}{15}{2}{6}
  {\begin{cellatartalom}%
     \oraresz{Gismondi}{mt{\" o}}{1}{I.39}{2}%
\\ \oraresz{Krasznay}{ol}{A2}{II.40}{2}%%
  \end{cellatartalom}}
\ora{\mitrakjonle}{15}{2}{7}
  {\begin{cellatartalom}%
     \oraresz{Eletto}{civ}{1}{F.10}{2}%
\\ \oraresz{Krasznay}{ol}{A2}{II.40}{2}%%
  \end{cellatartalom}}
\ora{\mitrakjonle}{15}{2}{8}
  {\begin{cellatartalom}%
     \oraresz{Krajny{\' a}k}{inf}{szak}{Szt.1}{4}%%
  \end{cellatartalom}}
\ora{\mitrakjonle}{15}{2}{9}
  {\begin{cellatartalom}%
    %
  \end{cellatartalom}}
\ora{\mitrakjonle}{15}{3}{0}
  {\begin{cellatartalom}%
     \oraresz{S{\" u}le}{i}{F}{I.34}{1.33333333333333}%
\\ \oraresz{Dezs{\H o}}{mt}{A1}{T.2}{1.33333333333333}%
\\ \oraresz{Krasznay}{ol}{A2}{Ol.2}{1.33333333333333}%%
  \end{cellatartalom}}
\ora{\mitrakjonle}{15}{3}{1}
  {\begin{cellatartalom}%
     \oraresz{Gismondi}{mt{\" o}}{1}{I.39}{2}%
\\ \oraresz{Krasznay}{ol}{A2}{Ol.2}{2}%%
  \end{cellatartalom}}
\ora{\mitrakjonle}{15}{3}{2}
  {\begin{cellatartalom}%
     \oraresz{Lovas}{t{\" o}}{A}{II.40}{2}%
\\ \oraresz{Eletto}{t{\" o}}{F}{F.9}{2}%%
  \end{cellatartalom}}
\ora{\mitrakjonle}{15}{3}{3}
  {\begin{cellatartalom}%
     \oraresz{Gismondi}{mt{\" o}}{2}{I.39}{2}%
\\ \oraresz{Krasznay}{ol}{A1}{Ol.2}{2}%%
  \end{cellatartalom}}
\ora{\mitrakjonle}{15}{3}{4}
  {\begin{cellatartalom}%
     \oraresz{Sch{\H o}n}{mt}{A2}{F.9}{1.33333333333333}%
\\ \oraresz{Fekete G.}{mt}{A1}{II.41}{1.33333333333333}%
\\ \oraresz{Dezs{\H o}}{mt}{F}{T.2}{1.33333333333333}%%
  \end{cellatartalom}}
\ora{\mitrakjonle}{15}{3}{5}
  {\begin{cellatartalom}%
     \oraresz{Dickerson}{a}{an3}{A.2}{1}%
\\ \oraresz{Liss{\' a}k}{a}{an2}{A.1}{1}%
\\ \oraresz{Csan{\' a}di}{a}{an1}{II.44}{1}%
\\ \oraresz{J.-Varga}{n{\' e}}{n{\' e}2}{II.41}{1}%%
  \end{cellatartalom}}
\ora{\mitrakjonle}{15}{3}{6}
  {\begin{cellatartalom}%
     \oraresz{Dickerson}{a}{an3}{A.2}{1}%
\\ \oraresz{Liss{\' a}k}{a}{an2}{A.1}{1}%
\\ \oraresz{Csan{\' a}di}{a}{an1}{II.44}{1}%
\\ \oraresz{J.-Varga}{n{\' e}}{n{\' e}2}{II.41}{1}%%
  \end{cellatartalom}}
\ora{\mitrakjonle}{15}{3}{7}
  {\begin{cellatartalom}%
     \oraresz{Tandory}{f{\" o}}{F}{II.48}{1.33333333333333}%
\\ \oraresz{Menyh{\' a}rt}{bi}{F}{B.ea}{1.33333333333333}%
\\ \oraresz{F{\" o}ldesi}{inf}{F}{Szt.2}{1.33333333333333}%%
  \end{cellatartalom}}
\ora{\mitrakjonle}{15}{3}{8}
  {\begin{cellatartalom}%
     \oraresz{Tandory}{f{\" o}}{F}{II.48}{1.33333333333333}%
\\ \oraresz{Menyh{\' a}rt}{bi}{F}{B.ea}{1.33333333333333}%
\\ \oraresz{F{\" o}ldesi}{inf}{F}{Szt.2}{1.33333333333333}%%
  \end{cellatartalom}}
\ora{\mitrakjonle}{15}{3}{9}
  {\begin{cellatartalom}%
    %
  \end{cellatartalom}}
\ora{\mitrakjonle}{15}{4}{0}
  {\begin{cellatartalom}%
     \oraresz{S{\' o}skuthyn{\' e}}{k{\' e}}{F}{K.ea}{2}%
\\ \oraresz{Dezs{\H o}}{mt}{F}{T.2}{2}%%
  \end{cellatartalom}}
\ora{\mitrakjonle}{15}{4}{1}
  {\begin{cellatartalom}%
     \oraresz{S{\' o}skuthyn{\' e}}{k{\' e}}{F}{K.ea}{2}%
\\ \oraresz{Dezs{\H o}}{mt}{F}{T.2}{2}%%
  \end{cellatartalom}}
\ora{\mitrakjonle}{15}{4}{2}
  {\begin{cellatartalom}%
     \oraresz{S{\" u}le}{i}{F}{I.30}{2}%
\\ \oraresz{Sch{\H o}n}{fi}{-}{Fiz.}{2}%%
  \end{cellatartalom}}
\ora{\mitrakjonle}{15}{4}{3}
  {\begin{cellatartalom}%
     \oraresz{Husz{\' a}r}{a}{an3}{A.1}{1}%
\\ \oraresz{Liss{\' a}k}{a}{an2}{II.51}{1}%
\\ \oraresz{Csan{\' a}di}{a}{an1}{II.40}{1}%
\\ \oraresz{J.-Varga}{n{\' e}}{n{\' e}2}{F.18}{1}%%
  \end{cellatartalom}}
\ora{\mitrakjonle}{15}{4}{4}
  {\begin{cellatartalom}%
     \oraresz{Eletto}{t{\" o}}{A}{F.11}{2}%
\\ \oraresz{Lovas}{t{\" o}}{F}{II.40}{2}%%
  \end{cellatartalom}}
\ora{\mitrakjonle}{15}{4}{5}
  {\begin{cellatartalom}%
     \oraresz{Krasznay}{ol}{A1}{Ol.1}{2}%
\\ \oraresz{Eletto}{civ}{2}{I.30}{2}%%
  \end{cellatartalom}}
\ora{\mitrakjonle}{15}{4}{6}
  {\begin{cellatartalom}%
     \oraresz{Nagy Sz.}{te}{L}{Tt.}{1.33333333333333}%
\\ \oraresz{Furi{\' a}k}{te}{F}{Tt.}{1.33333333333333}%
\\ \oraresz{De{\' a}k}{te}{L}{Tt.}{1.33333333333333}%%
  \end{cellatartalom}}
\ora{\mitrakjonle}{15}{4}{7}
  {\begin{cellatartalom}%
     \oraresz{S{\" u}le}{ny}{A}{II.43}{4}%%
  \end{cellatartalom}}
\ora{\mitrakjonle}{15}{4}{8}
  {\begin{cellatartalom}%
    %
  \end{cellatartalom}}
\ora{\mitrakjonle}{15}{4}{9}
  {\begin{cellatartalom}%
    %
  \end{cellatartalom}}
\ora{\mitrakjonle}{15}{5}{0}
  {\begin{cellatartalom}%
     \oraresz{Dezs{\H o}}{mt}{F}{T.2}{4}%%
  \end{cellatartalom}}
\ora{\mitrakjonle}{15}{5}{1}
  {\begin{cellatartalom}%
     \oraresz{Sch{\H o}n}{mt}{A2}{F.16}{1.33333333333333}%
\\ \oraresz{Fekete G.}{mt}{A1}{F.9}{1.33333333333333}%
\\ \oraresz{Dezs{\H o}}{mt}{F}{T.2}{1.33333333333333}%%
  \end{cellatartalom}}
\ora{\mitrakjonle}{15}{5}{2}
  {\begin{cellatartalom}%
     \oraresz{Dickerson}{a}{an3}{A.2}{1}%
\\ \oraresz{Liss{\' a}k}{a}{an2}{A.3}{1}%
\\ \oraresz{Csan{\' a}di}{a}{an1}{A.1}{1}%
\\ \oraresz{J.-Varga}{n{\' e}}{n{\' e}2}{I.39}{1}%%
  \end{cellatartalom}}
\ora{\mitrakjonle}{15}{5}{3}
  {\begin{cellatartalom}%
     \oraresz{Krasny{\' a}nszki}{bi}{-}{F.9}{4}%%
  \end{cellatartalom}}
\ora{\mitrakjonle}{15}{5}{4}
  {\begin{cellatartalom}%
     \oraresz{Lovas}{t{\" o}}{A}{II.40}{2}%
\\ \oraresz{Eletto}{t{\" o}}{F}{F.9}{2}%%
  \end{cellatartalom}}
\ora{\mitrakjonle}{15}{5}{5}
  {\begin{cellatartalom}%
     \oraresz{Eletto}{civ}{1}{F.9}{2}%
\\ \oraresz{Krasznay}{ol}{A2}{Ol.1}{2}%%
  \end{cellatartalom}}
\ora{\mitrakjonle}{15}{5}{6}
  {\begin{cellatartalom}%
     \oraresz{S{\" u}le}{i}{A}{F.9}{4}%%
  \end{cellatartalom}}
\ora{\mitrakjonle}{15}{5}{7}
  {\begin{cellatartalom}%
     \oraresz{Cserv{\' a}k}{inf}{F}{Szt.1}{4}%%
  \end{cellatartalom}}
\ora{\mitrakjonle}{15}{5}{8}
  {\begin{cellatartalom}%
     \oraresz{Cserv{\' a}k}{inf}{F}{Szt.1}{4}%%
  \end{cellatartalom}}
\ora{\mitrakjonle}{15}{5}{9}
  {\begin{cellatartalom}%
    %
  \end{cellatartalom}}
\ora{\mitrakjonle}{16}{1}{0}
  {\begin{cellatartalom}%
     \oraresz{Sebesty{\' e}n}{mt}{mt}{II.43}{5}%%
  \end{cellatartalom}}
\ora{\mitrakjonle}{16}{1}{1}
  {\begin{cellatartalom}%
     \oraresz{Sebesty{\' e}n}{mt}{mt}{II.43}{2.5}%
\\ \oraresz{V. Bartha}{ny}{an}{II.42}{2.5}%%
  \end{cellatartalom}}
\ora{\mitrakjonle}{16}{1}{2}
  {\begin{cellatartalom}%
     \oraresz{P{\' e}teri}{mt}{F}{M.}{1.66666666666667}%
\\ \oraresz{Sebesty{\' e}n}{mt}{an}{II.43}{1.66666666666667}%
\\ \oraresz{V. Bartha}{ny}{mt}{II.42}{1.66666666666667}%%
  \end{cellatartalom}}
\ora{\mitrakjonle}{16}{1}{3}
  {\begin{cellatartalom}%
     \oraresz{Kerling}{t{\" o}}{A}{II.43}{1.66666666666667}%
\\ \oraresz{Szendrei}{t{\" o}}{F}{II.48}{1.66666666666667}%
\\ \oraresz{S{\" u}le}{i}{F}{F.8}{1.66666666666667}%%
  \end{cellatartalom}}
\ora{\mitrakjonle}{16}{1}{4}
  {\begin{cellatartalom}%
     \oraresz{Sebesty{\' e}n}{of}{-}{II.43}{2.5}%
\\ \oraresz{S{\" u}le}{ny}{F}{F.8}{2.5}%%
  \end{cellatartalom}}
\ora{\mitrakjonle}{16}{1}{5}
  {\begin{cellatartalom}%
     \oraresz{Husz{\' a}r}{a}{an1}{A.1}{1.25}%
\\ \oraresz{Nagy Kata}{a}{an3}{F.11}{1.25}%
\\ \oraresz{Liss{\' a}k}{a}{an2}{II.44}{1.25}%
\\ \oraresz{Dickerson}{a}{an1}{II.40}{1.25}%%
  \end{cellatartalom}}
\ora{\mitrakjonle}{16}{1}{6}
  {\begin{cellatartalom}%
     \oraresz{Husz{\' a}r}{a}{an1}{A.1}{1.25}%
\\ \oraresz{Nagy Kata}{a}{an3}{F.11}{1.25}%
\\ \oraresz{Liss{\' a}k}{a}{an2}{II.44}{1.25}%
\\ \oraresz{Kassai}{a}{an1}{Kt.}{1.25}%%
  \end{cellatartalom}}
\ora{\mitrakjonle}{16}{1}{7}
  {\begin{cellatartalom}%
     \oraresz{Barth{\' a}n{\' e}}{zt}{zt}{{\' E}.}{1.66666666666667}%
\\ \oraresz{P{\' o}k}{zt}{{\' a}g}{R.}{1.66666666666667}%
\\ \oraresz{Juh{\' a}sz A.}{a}{oktv}{A.1}{1.66666666666667}%%
  \end{cellatartalom}}
\ora{\mitrakjonle}{16}{1}{8}
  {\begin{cellatartalom}%
     \oraresz{Barth{\' a}n{\' e}}{zt}{zt}{{\' E}.}{1.66666666666667}%
\\ \oraresz{P{\' o}k}{zt}{{\' a}g}{R.}{1.66666666666667}%
\\ \oraresz{Juh{\' a}sz A.}{a}{oktv}{A.1}{1.66666666666667}%%
  \end{cellatartalom}}
\ora{\mitrakjonle}{16}{1}{9}
  {\begin{cellatartalom}%
    %
  \end{cellatartalom}}
\ora{\mitrakjonle}{16}{2}{0}
  {\begin{cellatartalom}%
     \oraresz{S{\" u}le}{i}{F}{F.12}{1.25}%
\\ \oraresz{Szendrei}{t{\" o}}{F}{II.50}{1.25}%
\\ \oraresz{P{\' e}teri}{mt}{F}{II.43}{1.25}%
\\ \oraresz{Menyh{\' a}rt}{bi}{F}{B.gy}{1.25}%%
  \end{cellatartalom}}
\ora{\mitrakjonle}{16}{2}{1}
  {\begin{cellatartalom}%
     \oraresz{S{\" u}le}{i}{F}{F.12}{1.25}%
\\ \oraresz{Szendrei}{t{\" o}}{F}{II.50}{1.25}%
\\ \oraresz{P{\' e}teri}{mt}{F}{II.43}{1.25}%
\\ \oraresz{Menyh{\' a}rt}{bi}{F}{B.gy}{1.25}%%
  \end{cellatartalom}}
\ora{\mitrakjonle}{16}{2}{2}
  {\begin{cellatartalom}%
     \oraresz{Husz{\' a}r}{a}{an1}{Szt.3}{1.25}%
\\ \oraresz{Liss{\' a}k}{a}{an2}{F.8}{1.25}%
\\ \oraresz{Nagy Kata}{a}{an3}{A.3}{1.25}%
\\ \oraresz{Kassai}{a}{an1}{A.2}{1.25}%%
  \end{cellatartalom}}
\ora{\mitrakjonle}{16}{2}{3}
  {\begin{cellatartalom}%
     \oraresz{Szabolcsi}{fr}{fr}{Fr.}{1}%
\\ \oraresz{Pap}{sp}{sp2}{Sp.}{1}%
\\ \oraresz{Csuka}{a}{an2}{A.2}{1}%
\\ \oraresz{Tomor}{ol}{ol4}{F.9}{1}%
\\ \oraresz{J.-Varga}{n{\' e}}{n{\' e}5}{I.33}{1}%%
  \end{cellatartalom}}
\ora{\mitrakjonle}{16}{2}{4}
  {\begin{cellatartalom}%
     \oraresz{P{\' e}teri}{fi}{-}{Fiz.}{5}%%
  \end{cellatartalom}}
\ora{\mitrakjonle}{16}{2}{5}
  {\begin{cellatartalom}%
     \oraresz{V. Bartha}{i}{-}{II.43}{5}%%
  \end{cellatartalom}}
\ora{\mitrakjonle}{16}{2}{6}
  {\begin{cellatartalom}%
     \oraresz{Krasny{\' a}nszki}{bi}{an}{B.gy}{2.5}%
\\ \oraresz{Sebesty{\' e}n}{mt}{mt}{M.}{2.5}%%
  \end{cellatartalom}}
\ora{\mitrakjonle}{16}{2}{7}
  {\begin{cellatartalom}%
     \oraresz{Sebesty{\' e}n}{mt}{mt}{M.}{2.5}%
\\ \oraresz{Barth{\' a}n{\' e}}{zt}{zt}{{\' E}.}{2.5}%%
  \end{cellatartalom}}
\ora{\mitrakjonle}{16}{2}{8}
  {\begin{cellatartalom}%
     \oraresz{Barth{\' a}n{\' e}}{zt}{zt}{{\' E}.}{5}%%
  \end{cellatartalom}}
\ora{\mitrakjonle}{16}{2}{9}
  {\begin{cellatartalom}%
    %
  \end{cellatartalom}}
\ora{\mitrakjonle}{16}{3}{0}
  {\begin{cellatartalom}%
     \oraresz{S{\" u}le}{i}{F}{I.34}{2.5}%
\\ \oraresz{Sebesty{\' e}n}{mt}{mt}{M.}{2.5}%%
  \end{cellatartalom}}
\ora{\mitrakjonle}{16}{3}{1}
  {\begin{cellatartalom}%
     \oraresz{Heged{\H u}s}{te}{F}{Tt.}{2.5}%
\\ \oraresz{Bard{\' o}czky}{te}{L}{Tt.}{2.5}%%
  \end{cellatartalom}}
\ora{\mitrakjonle}{16}{3}{2}
  {\begin{cellatartalom}%
     \oraresz{Heged{\H u}s}{te}{F}{Tt.}{2.5}%
\\ \oraresz{Bard{\' o}czky}{te}{L}{Tt.}{2.5}%%
  \end{cellatartalom}}
\ora{\mitrakjonle}{16}{3}{3}
  {\begin{cellatartalom}%
     \oraresz{V. Bartha}{i}{-}{II.43}{5}%%
  \end{cellatartalom}}
\ora{\mitrakjonle}{16}{3}{4}
  {\begin{cellatartalom}%
     \oraresz{V. Bartha}{i}{-}{II.43}{5}%%
  \end{cellatartalom}}
\ora{\mitrakjonle}{16}{3}{5}
  {\begin{cellatartalom}%
     \oraresz{Krasny{\' a}nszki}{bi}{mt}{II.42}{1.66666666666667}%
\\ \oraresz{P{\' e}teri}{mt}{F}{M.}{1.66666666666667}%
\\ \oraresz{Sebesty{\' e}n}{mt}{an}{II.43}{1.66666666666667}%%
  \end{cellatartalom}}
\ora{\mitrakjonle}{16}{3}{6}
  {\begin{cellatartalom}%
     \oraresz{Kerling}{t{\" o}}{A}{II.43}{2.5}%
\\ \oraresz{Szendrei}{t{\" o}}{F}{II.48}{2.5}%%
  \end{cellatartalom}}
\ora{\mitrakjonle}{16}{3}{7}
  {\begin{cellatartalom}%
     \oraresz{Tandory}{f{\" o}}{F}{II.48}{1.25}%
\\ \oraresz{Menyh{\' a}rt}{bi}{F}{B.ea}{1.25}%
\\ \oraresz{P{\' o}k}{zt}{{\' a}g}{II.51}{1.25}%
\\ \oraresz{F{\" o}ldesi}{inf}{F}{Szt.2}{1.25}%%
  \end{cellatartalom}}
\ora{\mitrakjonle}{16}{3}{8}
  {\begin{cellatartalom}%
     \oraresz{Tandory}{f{\" o}}{F}{II.48}{1.25}%
\\ \oraresz{Menyh{\' a}rt}{bi}{F}{B.ea}{1.25}%
\\ \oraresz{P{\' o}k}{zt}{{\' a}g}{II.51}{1.25}%
\\ \oraresz{F{\" o}ldesi}{inf}{F}{Szt.2}{1.25}%%
  \end{cellatartalom}}
\ora{\mitrakjonle}{16}{3}{9}
  {\begin{cellatartalom}%
    %
  \end{cellatartalom}}
\ora{\mitrakjonle}{16}{4}{0}
  {\begin{cellatartalom}%
     \oraresz{Horv{\' a}th}{fi}{F}{Fiz.}{2.5}%
\\ \oraresz{S{\' o}skuthyn{\' e}}{k{\' e}}{F}{K.ea}{2.5}%%
  \end{cellatartalom}}
\ora{\mitrakjonle}{16}{4}{1}
  {\begin{cellatartalom}%
     \oraresz{Horv{\' a}th}{fi}{F}{Fiz.}{2.5}%
\\ \oraresz{S{\' o}skuthyn{\' e}}{k{\' e}}{F}{K.ea}{2.5}%%
  \end{cellatartalom}}
\ora{\mitrakjonle}{16}{4}{2}
  {\begin{cellatartalom}%
     \oraresz{S{\" u}le}{i}{F}{I.30}{2.5}%
\\ \oraresz{V. Bartha}{i}{-}{II.43}{2.5}%%
  \end{cellatartalom}}
\ora{\mitrakjonle}{16}{4}{3}
  {\begin{cellatartalom}%
     \oraresz{Krasny{\' a}nszki}{bi}{mt}{B.ea}{1.66666666666667}%
\\ \oraresz{P{\' e}teri}{mt}{F}{T.2}{1.66666666666667}%
\\ \oraresz{Sebesty{\' e}n}{mt}{an}{II.43}{1.66666666666667}%%
  \end{cellatartalom}}
\ora{\mitrakjonle}{16}{4}{4}
  {\begin{cellatartalom}%
     \oraresz{Krasny{\' a}nszki}{bi}{an}{B.ea}{2.5}%
\\ \oraresz{Sebesty{\' e}n}{mt}{mt}{II.43}{2.5}%%
  \end{cellatartalom}}
\ora{\mitrakjonle}{16}{4}{5}
  {\begin{cellatartalom}%
     \oraresz{Szabolcsi}{fr}{fr}{Fr.}{1}%
\\ \oraresz{Pap}{sp}{sp2}{Sp.}{1}%
\\ \oraresz{Csuka}{a}{an2}{II.51}{1}%
\\ \oraresz{Tomor}{ol}{ol4}{Ol.2}{1}%
\\ \oraresz{J.-Varga}{n{\' e}}{n{\' e}5}{F.12}{1}%%
  \end{cellatartalom}}
\ora{\mitrakjonle}{16}{4}{6}
  {\begin{cellatartalom}%
     \oraresz{Husz{\' a}r}{a}{an1}{A.1}{1.25}%
\\ \oraresz{Liss{\' a}k}{a}{an2}{I.31}{1.25}%
\\ \oraresz{Nagy Kata}{a}{an3}{I.30}{1.25}%
\\ \oraresz{Kassai}{a}{an1}{I.39}{1.25}%%
  \end{cellatartalom}}
\ora{\mitrakjonle}{16}{4}{7}
  {\begin{cellatartalom}%
     \oraresz{Barth{\' a}n{\' e}}{{\' e}k}{{\' e}k}{{\' E}.}{1.66666666666667}%
\\ \oraresz{P{\' o}k}{zt}{m{\H u}t}{II.48}{1.66666666666667}%
\\ \oraresz{Balanyi}{zt}{kgy}{R.}{1.66666666666667}%%
  \end{cellatartalom}}
\ora{\mitrakjonle}{16}{4}{8}
  {\begin{cellatartalom}%
     \oraresz{Barth{\' a}n{\' e}}{{\' e}k}{{\' e}k}{{\' E}.}{1.66666666666667}%
\\ \oraresz{P{\' o}k}{zt}{m{\H u}t}{II.48}{1.66666666666667}%
\\ \oraresz{Balanyi}{zt}{kgy}{R.}{1.66666666666667}%%
  \end{cellatartalom}}
\ora{\mitrakjonle}{16}{4}{9}
  {\begin{cellatartalom}%
    %
  \end{cellatartalom}}
\ora{\mitrakjonle}{16}{5}{0}
  {\begin{cellatartalom}%
     \oraresz{Dickerson}{a}{an1}{A.2}{2.5}%
\\ \oraresz{Kassai}{a}{an1}{A.3}{2.5}%%
  \end{cellatartalom}}
\ora{\mitrakjonle}{16}{5}{1}
  {\begin{cellatartalom}%
     \oraresz{Husz{\' a}r}{a}{an1}{A.2}{1.25}%
\\ \oraresz{Kassai}{a}{an1}{A.3}{1.25}%
\\ \oraresz{Liss{\' a}k}{a}{an2}{A.1}{1.25}%
\\ \oraresz{Nagy Kata}{a}{an3}{I.30}{1.25}%%
  \end{cellatartalom}}
\ora{\mitrakjonle}{16}{5}{2}
  {\begin{cellatartalom}%
     \oraresz{Heged{\H u}s}{te}{F}{Tt.}{2.5}%
\\ \oraresz{Bard{\' o}czky}{te}{L}{Tt.}{2.5}%%
  \end{cellatartalom}}
\ora{\mitrakjonle}{16}{5}{3}
  {\begin{cellatartalom}%
     \oraresz{Kerling}{t{\" o}}{A}{II.43}{2.5}%
\\ \oraresz{Szendrei}{t{\" o}}{F}{II.48}{2.5}%%
  \end{cellatartalom}}
\ora{\mitrakjonle}{16}{5}{4}
  {\begin{cellatartalom}%
     \oraresz{Szabolcsi}{fr}{fr}{Fr.}{1}%
\\ \oraresz{Pap}{sp}{sp2}{Sp.}{1}%
\\ \oraresz{Tomor}{ol}{ol4}{L.}{1}%
\\ \oraresz{Csuka}{a}{an2}{F.17}{1}%
\\ \oraresz{J.-Varga}{n{\' e}}{n{\' e}5}{F.10}{1}%%
  \end{cellatartalom}}
\ora{\mitrakjonle}{16}{5}{5}
  {\begin{cellatartalom}%
     \oraresz{Szabolcsi}{fr}{fr}{Fr.}{1}%
\\ \oraresz{Pap}{sp}{sp2}{Sp.}{1}%
\\ \oraresz{Tomor}{ol}{ol4}{L.}{1}%
\\ \oraresz{Csuka}{a}{an2}{F.17}{1}%
\\ \oraresz{J.-Varga}{n{\' e}}{n{\' e}5}{F.10}{1}%%
  \end{cellatartalom}}
\ora{\mitrakjonle}{16}{5}{6}
  {\begin{cellatartalom}%
     \oraresz{P{\' e}teri}{fi}{-}{Fiz.}{5}%%
  \end{cellatartalom}}
\ora{\mitrakjonle}{16}{5}{7}
  {\begin{cellatartalom}%
     \oraresz{Cserv{\' a}k}{inf}{F}{Szt.1}{5}%%
  \end{cellatartalom}}
\ora{\mitrakjonle}{16}{5}{8}
  {\begin{cellatartalom}%
     \oraresz{Cserv{\' a}k}{inf}{F}{Szt.1}{5}%%
  \end{cellatartalom}}
\ora{\mitrakjonle}{16}{5}{9}
  {\begin{cellatartalom}%
    %
  \end{cellatartalom}}
\ora{\mitrakjonle}{17}{1}{0}
  {\begin{cellatartalom}%
    %
  \end{cellatartalom}}
\ora{\mitrakjonle}{17}{1}{1}
  {\begin{cellatartalom}%
     \oraresz{Teremy}{t{\" o}}{F}{I.34}{3.5}%
\\ \oraresz{P.-P{\' a}sztor}{t{\" o}}{A}{II.50}{3.5}%%
  \end{cellatartalom}}
\ora{\mitrakjonle}{17}{1}{2}
  {\begin{cellatartalom}%
     \oraresz{P{\' e}teri}{mt}{F}{M.}{2.33333333333333}%
\\ \oraresz{Fekete G.}{mt}{r}{F.9}{2.33333333333333}%
\\ \oraresz{P{\' a}lmay}{mt}{n{\' e}}{I.39}{2.33333333333333}%%
  \end{cellatartalom}}
\ora{\mitrakjonle}{17}{1}{3}
  {\begin{cellatartalom}%
     \oraresz{S{\" u}le}{i}{F}{F.8}{3.5}%
\\ \oraresz{P.-P{\' a}sztor}{i}{A}{F.9}{3.5}%%
  \end{cellatartalom}}
\ora{\mitrakjonle}{17}{1}{4}
  {\begin{cellatartalom}%
     \oraresz{S{\" u}le}{ny}{F}{F.8}{3.5}%
\\ \oraresz{P.-P{\' a}sztor}{i}{A}{F.9}{3.5}%%
  \end{cellatartalom}}
\ora{\mitrakjonle}{17}{1}{5}
  {\begin{cellatartalom}%
     \oraresz{G{\' a}l}{n{\' e}}{n{\' e}1}{N.}{1.75}%
\\ \oraresz{Nagy Kata}{a}{an3}{F.11}{1.75}%
\\ \oraresz{Liss{\' a}k}{a}{an2}{II.44}{1.75}%
\\ \oraresz{Dickerson}{a}{an1}{II.40}{1.75}%%
  \end{cellatartalom}}
\ora{\mitrakjonle}{17}{1}{6}
  {\begin{cellatartalom}%
     \oraresz{M{\' e}sz{\' a}ros}{n{\' e}}{n{\' e}1}{N.}{1.75}%
\\ \oraresz{Nagy Kata}{a}{an3}{F.11}{1.75}%
\\ \oraresz{Liss{\' a}k}{a}{an2}{II.44}{1.75}%
\\ \oraresz{Kassai}{a}{an1}{Kt.}{1.75}%%
  \end{cellatartalom}}
\ora{\mitrakjonle}{17}{1}{7}
  {\begin{cellatartalom}%
     \oraresz{Balanyi}{r}{r}{II.51}{1.75}%
\\ \oraresz{Barth{\' a}n{\' e}}{zt}{zt}{{\' E}.}{1.75}%
\\ \oraresz{P{\' o}k}{zt}{{\' a}g}{R.}{1.75}%
\\ \oraresz{Juh{\' a}sz A.}{a}{oktv}{A.1}{1.75}%%
  \end{cellatartalom}}
\ora{\mitrakjonle}{17}{1}{8}
  {\begin{cellatartalom}%
     \oraresz{Barth{\' a}n{\' e}}{zt}{zt}{{\' E}.}{2.33333333333333}%
\\ \oraresz{P{\' o}k}{zt}{{\' a}g}{R.}{2.33333333333333}%
\\ \oraresz{Juh{\' a}sz A.}{a}{oktv}{A.1}{2.33333333333333}%%
  \end{cellatartalom}}
\ora{\mitrakjonle}{17}{1}{9}
  {\begin{cellatartalom}%
    %
  \end{cellatartalom}}
\ora{\mitrakjonle}{17}{2}{0}
  {\begin{cellatartalom}%
     \oraresz{S{\" u}le}{i}{F}{F.12}{2.33333333333333}%
\\ \oraresz{P{\' e}teri}{mt}{F}{II.43}{2.33333333333333}%
\\ \oraresz{Menyh{\' a}rt}{bi}{F}{B.gy}{2.33333333333333}%%
  \end{cellatartalom}}
\ora{\mitrakjonle}{17}{2}{1}
  {\begin{cellatartalom}%
     \oraresz{S{\" u}le}{i}{F}{F.12}{2.33333333333333}%
\\ \oraresz{P{\' e}teri}{mt}{F}{II.43}{2.33333333333333}%
\\ \oraresz{Menyh{\' a}rt}{bi}{F}{B.gy}{2.33333333333333}%%
  \end{cellatartalom}}
\ora{\mitrakjonle}{17}{2}{2}
  {\begin{cellatartalom}%
     \oraresz{M{\' e}sz{\' a}ros}{n{\' e}}{n{\' e}1}{F.9}{1.75}%
\\ \oraresz{Liss{\' a}k}{a}{an2}{F.8}{1.75}%
\\ \oraresz{Nagy Kata}{a}{an3}{A.3}{1.75}%
\\ \oraresz{Kassai}{a}{an1}{A.2}{1.75}%%
  \end{cellatartalom}}
\ora{\mitrakjonle}{17}{2}{3}
  {\begin{cellatartalom}%
     \oraresz{Szabolcsi}{fr}{fr}{Fr.}{1}%
\\ \oraresz{Pap}{sp}{sp2}{Sp.}{1}%
\\ \oraresz{Csuka}{a}{an2}{A.2}{1}%
\\ \oraresz{Nagy Kata}{a}{an3}{F.18}{1}%
\\ \oraresz{Varga Bo.}{a}{an1}{A.3}{1}%
\\ \oraresz{Tomor}{ol}{ol4}{F.9}{1}%
\\ \oraresz{J.-Varga}{n{\' e}}{n{\' e}5}{I.33}{1}%%
  \end{cellatartalom}}
\ora{\mitrakjonle}{17}{2}{4}
  {\begin{cellatartalom}%
     \oraresz{P.-P{\' a}sztor}{of}{-}{F.9}{7}%%
  \end{cellatartalom}}
\ora{\mitrakjonle}{17}{2}{5}
  {\begin{cellatartalom}%
     \oraresz{Menyh{\' a}rt}{bi}{-}{B.ea}{7}%%
  \end{cellatartalom}}
\ora{\mitrakjonle}{17}{2}{6}
  {\begin{cellatartalom}%
     \oraresz{Teremy}{t{\" o}}{F}{II.50}{7}%%
  \end{cellatartalom}}
\ora{\mitrakjonle}{17}{2}{7}
  {\begin{cellatartalom}%
     \oraresz{Balanyi}{r}{r}{R.}{3.5}%
\\ \oraresz{Barth{\' a}n{\' e}}{zt}{zt}{{\' E}.}{3.5}%%
  \end{cellatartalom}}
\ora{\mitrakjonle}{17}{2}{8}
  {\begin{cellatartalom}%
     \oraresz{Balanyi}{r}{r}{R.}{3.5}%
\\ \oraresz{Barth{\' a}n{\' e}}{zt}{zt}{{\' E}.}{3.5}%%
  \end{cellatartalom}}
\ora{\mitrakjonle}{17}{2}{9}
  {\begin{cellatartalom}%
    %
  \end{cellatartalom}}
\ora{\mitrakjonle}{17}{3}{0}
  {\begin{cellatartalom}%
     \oraresz{S{\" u}le}{i}{F}{I.34}{3.5}%
\\ \oraresz{P.-P{\' a}sztor}{i}{A}{F.9}{3.5}%%
  \end{cellatartalom}}
\ora{\mitrakjonle}{17}{3}{1}
  {\begin{cellatartalom}%
     \oraresz{Sch{\H o}n}{fi}{-}{F.9}{7}%%
  \end{cellatartalom}}
\ora{\mitrakjonle}{17}{3}{2}
  {\begin{cellatartalom}%
     \oraresz{Menyh{\' a}rt}{bi}{-}{K.ea}{7}%%
  \end{cellatartalom}}
\ora{\mitrakjonle}{17}{3}{3}
  {\begin{cellatartalom}%
     \oraresz{Vas E.}{te}{L}{Tt.}{2.33333333333333}%
\\ \oraresz{Lad{\' a}nyi}{te}{L}{Tt.}{2.33333333333333}%
\\ \oraresz{De{\' a}k}{te}{F}{Tt.}{2.33333333333333}%%
  \end{cellatartalom}}
\ora{\mitrakjonle}{17}{3}{4}
  {\begin{cellatartalom}%
     \oraresz{Vas E.}{te}{L}{Tt.}{2.33333333333333}%
\\ \oraresz{Lad{\' a}nyi}{te}{L}{Tt.}{2.33333333333333}%
\\ \oraresz{De{\' a}k}{te}{F}{Tt.}{2.33333333333333}%%
  \end{cellatartalom}}
\ora{\mitrakjonle}{17}{3}{5}
  {\begin{cellatartalom}%
     \oraresz{P{\' e}teri}{mt}{F}{M.}{2.33333333333333}%
\\ \oraresz{Fekete G.}{mt}{r}{II.51}{2.33333333333333}%
\\ \oraresz{P{\' a}lmay}{mt}{n{\' e}}{F.9}{2.33333333333333}%%
  \end{cellatartalom}}
\ora{\mitrakjonle}{17}{3}{6}
  {\begin{cellatartalom}%
     \oraresz{P.-P{\' a}sztor}{t{\" o}}{A}{F.9}{7}%%
  \end{cellatartalom}}
\ora{\mitrakjonle}{17}{3}{7}
  {\begin{cellatartalom}%
     \oraresz{Tandory}{f{\" o}}{F}{II.48}{1.75}%
\\ \oraresz{Menyh{\' a}rt}{bi}{F}{B.ea}{1.75}%
\\ \oraresz{P{\' o}k}{zt}{{\' a}g}{II.51}{1.75}%
\\ \oraresz{F{\" o}ldesi}{inf}{F}{Szt.2}{1.75}%%
  \end{cellatartalom}}
\ora{\mitrakjonle}{17}{3}{8}
  {\begin{cellatartalom}%
     \oraresz{Tandory}{f{\" o}}{F}{II.48}{1.75}%
\\ \oraresz{Menyh{\' a}rt}{bi}{F}{B.ea}{1.75}%
\\ \oraresz{P{\' o}k}{zt}{{\' a}g}{II.51}{1.75}%
\\ \oraresz{F{\" o}ldesi}{inf}{F}{Szt.2}{1.75}%%
  \end{cellatartalom}}
\ora{\mitrakjonle}{17}{3}{9}
  {\begin{cellatartalom}%
    %
  \end{cellatartalom}}
\ora{\mitrakjonle}{17}{4}{0}
  {\begin{cellatartalom}%
     \oraresz{S{\' o}skuthyn{\' e}}{k{\' e}}{F}{K.ea}{3.5}%
\\ \oraresz{Teremy}{t{\" o}}{F}{F.9}{3.5}%%
  \end{cellatartalom}}
\ora{\mitrakjonle}{17}{4}{1}
  {\begin{cellatartalom}%
     \oraresz{S{\' o}skuthyn{\' e}}{k{\' e}}{F}{K.ea}{3.5}%
\\ \oraresz{Teremy}{t{\" o}}{F}{F.9}{3.5}%%
  \end{cellatartalom}}
\ora{\mitrakjonle}{17}{4}{2}
  {\begin{cellatartalom}%
     \oraresz{S{\" u}le}{i}{F}{I.30}{3.5}%
\\ \oraresz{P.-P{\' a}sztor}{ny}{A}{II.50}{3.5}%%
  \end{cellatartalom}}
\ora{\mitrakjonle}{17}{4}{3}
  {\begin{cellatartalom}%
     \oraresz{P{\' e}teri}{mt}{F}{T.2}{2.33333333333333}%
\\ \oraresz{Fekete G.}{mt}{r}{F.12}{2.33333333333333}%
\\ \oraresz{P{\' a}lmay}{mt}{n{\' e}}{F.9}{2.33333333333333}%%
  \end{cellatartalom}}
\ora{\mitrakjonle}{17}{4}{4}
  {\begin{cellatartalom}%
     \oraresz{Sch{\H o}n}{fi}{-}{F.9}{7}%%
  \end{cellatartalom}}
\ora{\mitrakjonle}{17}{4}{5}
  {\begin{cellatartalom}%
     \oraresz{Szabolcsi}{fr}{fr}{Fr.}{1}%
\\ \oraresz{Pap}{sp}{sp2}{Sp.}{1}%
\\ \oraresz{Csuka}{a}{an2}{II.51}{1}%
\\ \oraresz{Nagy Kata}{a}{an3}{II.42}{1}%
\\ \oraresz{Varga Bo.}{a}{an1}{Szt.3}{1}%
\\ \oraresz{Tomor}{ol}{ol4}{Ol.2}{1}%
\\ \oraresz{J.-Varga}{n{\' e}}{n{\' e}5}{F.12}{1}%%
  \end{cellatartalom}}
\ora{\mitrakjonle}{17}{4}{6}
  {\begin{cellatartalom}%
     \oraresz{M{\' e}sz{\' a}ros}{n{\' e}}{n{\' e}1}{N.}{1.75}%
\\ \oraresz{Liss{\' a}k}{a}{an2}{I.31}{1.75}%
\\ \oraresz{Nagy Kata}{a}{an3}{I.30}{1.75}%
\\ \oraresz{Kassai}{a}{an1}{I.39}{1.75}%%
  \end{cellatartalom}}
\ora{\mitrakjonle}{17}{4}{7}
  {\begin{cellatartalom}%
     \oraresz{Barth{\' a}n{\' e}}{{\' e}k}{{\' e}k}{{\' E}.}{2.33333333333333}%
\\ \oraresz{P{\' o}k}{zt}{m{\H u}t}{II.48}{2.33333333333333}%
\\ \oraresz{Balanyi}{zt}{kgy}{R.}{2.33333333333333}%%
  \end{cellatartalom}}
\ora{\mitrakjonle}{17}{4}{8}
  {\begin{cellatartalom}%
     \oraresz{Barth{\' a}n{\' e}}{{\' e}k}{{\' e}k}{{\' E}.}{2.33333333333333}%
\\ \oraresz{P{\' o}k}{zt}{m{\H u}t}{II.48}{2.33333333333333}%
\\ \oraresz{Balanyi}{zt}{kgy}{R.}{2.33333333333333}%%
  \end{cellatartalom}}
\ora{\mitrakjonle}{17}{4}{9}
  {\begin{cellatartalom}%
    %
  \end{cellatartalom}}
\ora{\mitrakjonle}{17}{5}{0}
  {\begin{cellatartalom}%
     \oraresz{Kassai}{a}{an1}{A.3}{3.5}%
\\ \oraresz{G{\' a}l}{n{\' e}}{n{\' e}1}{N.}{3.5}%%
  \end{cellatartalom}}
\ora{\mitrakjonle}{17}{5}{1}
  {\begin{cellatartalom}%
     \oraresz{M{\' e}sz{\' a}ros}{n{\' e}}{n{\' e}1}{N.}{1.75}%
\\ \oraresz{Kassai}{a}{an1}{A.3}{1.75}%
\\ \oraresz{Liss{\' a}k}{a}{an2}{A.1}{1.75}%
\\ \oraresz{Nagy Kata}{a}{an3}{I.30}{1.75}%%
  \end{cellatartalom}}
\ora{\mitrakjonle}{17}{5}{2}
  {\begin{cellatartalom}%
     \oraresz{Teremy}{t{\" o}}{F}{II.43}{3.5}%
\\ \oraresz{P.-P{\' a}sztor}{t{\" o}}{A}{F.9}{3.5}%%
  \end{cellatartalom}}
\ora{\mitrakjonle}{17}{5}{3}
  {\begin{cellatartalom}%
     \oraresz{Vas E.}{te}{L}{Tt.}{2.33333333333333}%
\\ \oraresz{Lad{\' a}nyi}{te}{L}{Tt.}{2.33333333333333}%
\\ \oraresz{De{\' a}k}{te}{F}{Tt.}{2.33333333333333}%%
  \end{cellatartalom}}
\ora{\mitrakjonle}{17}{5}{4}
  {\begin{cellatartalom}%
     \oraresz{Nagy Kata}{a}{an3}{A.2}{1}%
\\ \oraresz{Szabolcsi}{fr}{fr}{Fr.}{1}%
\\ \oraresz{Pap}{sp}{sp2}{Sp.}{1}%
\\ \oraresz{Tomor}{ol}{ol4}{L.}{1}%
\\ \oraresz{Csuka}{a}{an2}{F.17}{1}%
\\ \oraresz{Varga Bo.}{a}{an1}{{\' E}.}{1}%
\\ \oraresz{J.-Varga}{n{\' e}}{n{\' e}5}{F.10}{1}%%
  \end{cellatartalom}}
\ora{\mitrakjonle}{17}{5}{5}
  {\begin{cellatartalom}%
     \oraresz{Nagy Kata}{a}{an3}{A.2}{1}%
\\ \oraresz{Szabolcsi}{fr}{fr}{Fr.}{1}%
\\ \oraresz{Pap}{sp}{sp2}{Sp.}{1}%
\\ \oraresz{Tomor}{ol}{ol4}{L.}{1}%
\\ \oraresz{Csuka}{a}{an2}{F.17}{1}%
\\ \oraresz{Varga Bo.}{a}{an1}{{\' E}.}{1}%
\\ \oraresz{J.-Varga}{n{\' e}}{n{\' e}5}{F.10}{1}%%
  \end{cellatartalom}}
\ora{\mitrakjonle}{17}{5}{6}
  {\begin{cellatartalom}%
     \oraresz{P{\' o}k}{r}{r}{R.}{7}%%
  \end{cellatartalom}}
\ora{\mitrakjonle}{17}{5}{7}
  {\begin{cellatartalom}%
     \oraresz{Cserv{\' a}k}{inf}{F}{Szt.1}{7}%%
  \end{cellatartalom}}
\ora{\mitrakjonle}{17}{5}{8}
  {\begin{cellatartalom}%
     \oraresz{Cserv{\' a}k}{inf}{F}{Szt.1}{7}%%
  \end{cellatartalom}}
\ora{\mitrakjonle}{17}{5}{9}
  {\begin{cellatartalom}%
    %
  \end{cellatartalom}}
\ora{\mitrakjonle}{18}{1}{0}
  {\begin{cellatartalom}%
     \oraresz{Hal{\' a}sz}{n{\' e}}{n{\' e}}{F.18}{7}%%
  \end{cellatartalom}}
\ora{\mitrakjonle}{18}{1}{1}
  {\begin{cellatartalom}%
     \oraresz{Teremy}{t{\" o}}{F}{I.34}{3.5}%
\\ \oraresz{Kerling}{t{\" o}}{A}{II.44}{3.5}%%
  \end{cellatartalom}}
\ora{\mitrakjonle}{18}{1}{2}
  {\begin{cellatartalom}%
     \oraresz{P{\' e}teri}{mt}{F}{M.}{1}%
\\ \oraresz{Kudella}{fr}{fr}{Fr.}{1}%
\\ \oraresz{Sum}{sp}{sp1}{Sp.}{1}%
\\ \oraresz{Csuka}{ol}{ol3}{II.40}{1}%
\\ \oraresz{G{\' a}l}{n{\' e}}{n{\' e}2}{N.}{1}%
\\ \oraresz{Hal{\' a}sz}{n{\' e}}{n{\' e}4}{Kt.}{1}%
\\ \oraresz{S{\' o}tin{\' e}}{n{\' e}}{n{\' e}3}{F.18}{1}%%
  \end{cellatartalom}}
\ora{\mitrakjonle}{18}{1}{3}
  {\begin{cellatartalom}%
     \oraresz{S{\" u}le}{i}{F}{F.8}{3.5}%
\\ \oraresz{Oravecz}{i}{A}{II.44}{3.5}%%
  \end{cellatartalom}}
\ora{\mitrakjonle}{18}{1}{4}
  {\begin{cellatartalom}%
     \oraresz{S{\" u}le}{ny}{F}{F.8}{3.5}%
\\ \oraresz{Oravecz}{ny}{A}{II.44}{3.5}%%
  \end{cellatartalom}}
\ora{\mitrakjonle}{18}{1}{5}
  {\begin{cellatartalom}%
     \oraresz{P{\' e}teri}{fi}{-}{B.ea}{7}%%
  \end{cellatartalom}}
\ora{\mitrakjonle}{18}{1}{6}
  {\begin{cellatartalom}%
     \oraresz{Hajba}{a}{an6}{{\' E}.}{3.5}%
\\ \oraresz{Dickerson}{a}{an5}{F.16}{3.5}%%
  \end{cellatartalom}}
\ora{\mitrakjonle}{18}{1}{7}
  {\begin{cellatartalom}%
     \oraresz{Barth{\' a}n{\' e}}{zt}{zt}{{\' E}.}{2.33333333333333}%
\\ \oraresz{P{\' o}k}{zt}{{\' a}g}{R.}{2.33333333333333}%
\\ \oraresz{Juh{\' a}sz A.}{a}{oktv}{A.1}{2.33333333333333}%%
  \end{cellatartalom}}
\ora{\mitrakjonle}{18}{1}{8}
  {\begin{cellatartalom}%
     \oraresz{Barth{\' a}n{\' e}}{zt}{zt}{{\' E}.}{2.33333333333333}%
\\ \oraresz{P{\' o}k}{zt}{{\' a}g}{R.}{2.33333333333333}%
\\ \oraresz{Juh{\' a}sz A.}{a}{oktv}{A.1}{2.33333333333333}%%
  \end{cellatartalom}}
\ora{\mitrakjonle}{18}{1}{9}
  {\begin{cellatartalom}%
    %
  \end{cellatartalom}}
\ora{\mitrakjonle}{18}{2}{0}
  {\begin{cellatartalom}%
     \oraresz{S{\" u}le}{i}{F}{F.12}{2.33333333333333}%
\\ \oraresz{S{\' o}skuthyn{\' e}}{k{\' e}}{F}{K.ea}{2.33333333333333}%
\\ \oraresz{P{\' e}teri}{mt}{F}{II.43}{2.33333333333333}%%
  \end{cellatartalom}}
\ora{\mitrakjonle}{18}{2}{1}
  {\begin{cellatartalom}%
     \oraresz{S{\" u}le}{i}{F}{F.12}{2.33333333333333}%
\\ \oraresz{S{\' o}skuthyn{\' e}}{k{\' e}}{F}{K.ea}{2.33333333333333}%
\\ \oraresz{P{\' e}teri}{mt}{F}{II.43}{2.33333333333333}%%
  \end{cellatartalom}}
\ora{\mitrakjonle}{18}{2}{2}
  {\begin{cellatartalom}%
     \oraresz{Dickerson}{a}{an6}{B.gy}{3.5}%
\\ \oraresz{Hajba}{a}{an5}{II.43}{3.5}%%
  \end{cellatartalom}}
\ora{\mitrakjonle}{18}{2}{3}
  {\begin{cellatartalom}%
     \oraresz{Horv{\' a}th G.}{bi}{2}{B.gy}{3.5}%
\\ \oraresz{Darab{\' a}nt}{mt}{A1}{II.44}{3.5}%%
  \end{cellatartalom}}
\ora{\mitrakjonle}{18}{2}{4}
  {\begin{cellatartalom}%
     \oraresz{Horv{\' a}th G.}{bi}{1}{B.gy}{3.5}%
\\ \oraresz{Varga Be.}{mt}{A2}{II.42}{3.5}%%
  \end{cellatartalom}}
\ora{\mitrakjonle}{18}{2}{5}
  {\begin{cellatartalom}%
     \oraresz{Horv{\' a}th G.}{bi}{1}{B.gy}{3.5}%
\\ \oraresz{Varga Be.}{mt}{A2}{II.42}{3.5}%%
  \end{cellatartalom}}
\ora{\mitrakjonle}{18}{2}{6}
  {\begin{cellatartalom}%
     \oraresz{Teremy}{t{\" o}}{F}{II.50}{3.5}%
\\ \oraresz{Kerling}{t{\" o}}{A}{II.44}{3.5}%%
  \end{cellatartalom}}
\ora{\mitrakjonle}{18}{2}{7}
  {\begin{cellatartalom}%
     \oraresz{Barth{\' a}n{\' e}}{zt}{zt}{{\' E}.}{7}%%
  \end{cellatartalom}}
\ora{\mitrakjonle}{18}{2}{8}
  {\begin{cellatartalom}%
     \oraresz{Barth{\' a}n{\' e}}{zt}{zt}{{\' E}.}{7}%%
  \end{cellatartalom}}
\ora{\mitrakjonle}{18}{2}{9}
  {\begin{cellatartalom}%
    %
  \end{cellatartalom}}
\ora{\mitrakjonle}{18}{3}{0}
  {\begin{cellatartalom}%
     \oraresz{S{\" u}le}{i}{F}{I.34}{7}%%
  \end{cellatartalom}}
\ora{\mitrakjonle}{18}{3}{1}
  {\begin{cellatartalom}%
     \oraresz{Dickerson}{a}{an5}{Szt.3}{3.5}%
\\ \oraresz{Hajba}{a}{an6}{F.18}{3.5}%%
  \end{cellatartalom}}
\ora{\mitrakjonle}{18}{3}{2}
  {\begin{cellatartalom}%
     \oraresz{Horv{\' a}th G.}{bi}{2}{B.gy}{3.5}%
\\ \oraresz{Darab{\' a}nt}{mt}{A1}{II.41}{3.5}%%
  \end{cellatartalom}}
\ora{\mitrakjonle}{18}{3}{3}
  {\begin{cellatartalom}%
     \oraresz{Horv{\' a}th G.}{bi}{2}{B.gy}{3.5}%
\\ \oraresz{Darab{\' a}nt}{mt}{A1}{II.41}{3.5}%%
  \end{cellatartalom}}
\ora{\mitrakjonle}{18}{3}{4}
  {\begin{cellatartalom}%
     \oraresz{Horv{\' a}th G.}{bi}{1}{B.gy}{3.5}%
\\ \oraresz{Varga Be.}{mt}{A2}{II.42}{3.5}%%
  \end{cellatartalom}}
\ora{\mitrakjonle}{18}{3}{5}
  {\begin{cellatartalom}%
     \oraresz{P{\' e}teri}{mt}{F}{M.}{2.33333333333333}%
\\ \oraresz{Nagy Sz.}{te}{L}{Tt.}{2.33333333333333}%
\\ \oraresz{De{\' a}k}{te}{F}{Tt.}{2.33333333333333}%%
  \end{cellatartalom}}
\ora{\mitrakjonle}{18}{3}{6}
  {\begin{cellatartalom}%
     \oraresz{P{\' e}teri}{fi}{-}{Fiz.}{7}%%
  \end{cellatartalom}}
\ora{\mitrakjonle}{18}{3}{7}
  {\begin{cellatartalom}%
     \oraresz{Barth{\' a}n{\' e}}{zt}{zt}{{\' E}.}{3.5}%
\\ \oraresz{P{\' o}k}{zt}{{\' a}g}{II.51}{3.5}%%
  \end{cellatartalom}}
\ora{\mitrakjonle}{18}{3}{8}
  {\begin{cellatartalom}%
     \oraresz{Barth{\' a}n{\' e}}{zt}{zt}{{\' E}.}{3.5}%
\\ \oraresz{P{\' o}k}{zt}{{\' a}g}{II.51}{3.5}%%
  \end{cellatartalom}}
\ora{\mitrakjonle}{18}{3}{9}
  {\begin{cellatartalom}%
    %
  \end{cellatartalom}}
\ora{\mitrakjonle}{18}{4}{0}
  {\begin{cellatartalom}%
     \oraresz{S{\' o}skuthyn{\' e}}{k{\' e}}{F}{K.ea}{3.5}%
\\ \oraresz{Teremy}{t{\" o}}{F}{F.9}{3.5}%%
  \end{cellatartalom}}
\ora{\mitrakjonle}{18}{4}{1}
  {\begin{cellatartalom}%
     \oraresz{S{\' o}skuthyn{\' e}}{k{\' e}}{F}{K.ea}{3.5}%
\\ \oraresz{Teremy}{t{\" o}}{F}{F.9}{3.5}%%
  \end{cellatartalom}}
\ora{\mitrakjonle}{18}{4}{2}
  {\begin{cellatartalom}%
     \oraresz{S{\" u}le}{i}{F}{I.30}{3.5}%
\\ \oraresz{Oravecz}{i}{A}{II.44}{3.5}%%
  \end{cellatartalom}}
\ora{\mitrakjonle}{18}{4}{3}
  {\begin{cellatartalom}%
     \oraresz{P{\' e}teri}{mt}{F}{T.2}{2.33333333333333}%
\\ \oraresz{Hajba}{a}{an6}{A.3}{2.33333333333333}%
\\ \oraresz{Dickerson}{a}{an5}{II.44}{2.33333333333333}%%
  \end{cellatartalom}}
\ora{\mitrakjonle}{18}{4}{4}
  {\begin{cellatartalom}%
     \oraresz{Hajba}{a}{an6}{A.3}{3.5}%
\\ \oraresz{Dickerson}{a}{an5}{II.44}{3.5}%%
  \end{cellatartalom}}
\ora{\mitrakjonle}{18}{4}{5}
  {\begin{cellatartalom}%
     \oraresz{Horv{\' a}th G.}{bi}{2}{B.ea}{3.5}%
\\ \oraresz{Darab{\' a}nt}{mt}{A1}{M.}{3.5}%%
  \end{cellatartalom}}
\ora{\mitrakjonle}{18}{4}{6}
  {\begin{cellatartalom}%
     \oraresz{Kudella}{fr}{fr}{Fr.}{1.16666666666667}%
\\ \oraresz{Sum}{sp}{sp1}{Sp.}{1.16666666666667}%
\\ \oraresz{Csuka}{ol}{ol3}{L.}{1.16666666666667}%
\\ \oraresz{G{\' a}l}{n{\' e}}{n{\' e}2}{II.40}{1.16666666666667}%
\\ \oraresz{Hal{\' a}sz}{n{\' e}}{n{\' e}4}{II.41}{1.16666666666667}%
\\ \oraresz{S{\' o}tin{\' e}}{n{\' e}}{n{\' e}3}{{\' E}.}{1.16666666666667}%%
  \end{cellatartalom}}
\ora{\mitrakjonle}{18}{4}{7}
  {\begin{cellatartalom}%
     \oraresz{Barth{\' a}n{\' e}}{{\' e}k}{{\' e}k}{{\' E}.}{2.33333333333333}%
\\ \oraresz{P{\' o}k}{zt}{m{\H u}t}{II.48}{2.33333333333333}%
\\ \oraresz{Balanyi}{zt}{kgy}{R.}{2.33333333333333}%%
  \end{cellatartalom}}
\ora{\mitrakjonle}{18}{4}{8}
  {\begin{cellatartalom}%
     \oraresz{Barth{\' a}n{\' e}}{{\' e}k}{{\' e}k}{{\' E}.}{2.33333333333333}%
\\ \oraresz{P{\' o}k}{zt}{m{\H u}t}{II.48}{2.33333333333333}%
\\ \oraresz{Balanyi}{zt}{kgy}{R.}{2.33333333333333}%%
  \end{cellatartalom}}
\ora{\mitrakjonle}{18}{4}{9}
  {\begin{cellatartalom}%
    %
  \end{cellatartalom}}
\ora{\mitrakjonle}{18}{5}{0}
  {\begin{cellatartalom}%
    %
  \end{cellatartalom}}
\ora{\mitrakjonle}{18}{5}{1}
  {\begin{cellatartalom}%
     \oraresz{Kudella}{fr}{fr}{Fr.}{1.16666666666667}%
\\ \oraresz{Sum}{sp}{sp1}{Sp.}{1.16666666666667}%
\\ \oraresz{Csuka}{ol}{ol3}{I.32}{1.16666666666667}%
\\ \oraresz{G{\' a}l}{n{\' e}}{n{\' e}2}{I.33}{1.16666666666667}%
\\ \oraresz{Hal{\' a}sz}{n{\' e}}{n{\' e}4}{II.41}{1.16666666666667}%
\\ \oraresz{S{\' o}tin{\' e}}{n{\' e}}{n{\' e}3}{{\' E}.}{1.16666666666667}%%
  \end{cellatartalom}}
\ora{\mitrakjonle}{18}{5}{2}
  {\begin{cellatartalom}%
     \oraresz{Teremy}{t{\" o}}{F}{II.43}{3.5}%
\\ \oraresz{Kerling}{t{\" o}}{A}{II.44}{3.5}%%
  \end{cellatartalom}}
\ora{\mitrakjonle}{18}{5}{3}
  {\begin{cellatartalom}%
     \oraresz{Horv{\' a}th G.}{bi}{1}{K.gy}{3.5}%
\\ \oraresz{Varga Be.}{mt}{A2}{II.41}{3.5}%%
  \end{cellatartalom}}
\ora{\mitrakjonle}{18}{5}{4}
  {\begin{cellatartalom}%
     \oraresz{Nagy Sz.}{te}{L}{Tt.}{3.5}%
\\ \oraresz{De{\' a}k}{te}{F}{Tt.}{3.5}%%
  \end{cellatartalom}}
\ora{\mitrakjonle}{18}{5}{5}
  {\begin{cellatartalom}%
     \oraresz{Nagy Sz.}{te}{L}{Tt.}{3.5}%
\\ \oraresz{De{\' a}k}{te}{F}{Tt.}{3.5}%%
  \end{cellatartalom}}
\ora{\mitrakjonle}{18}{5}{6}
  {\begin{cellatartalom}%
     \oraresz{S{\' o}skuthyn{\' e}}{of}{-}{II.44}{7}%%
  \end{cellatartalom}}
\ora{\mitrakjonle}{18}{5}{7}
  {\begin{cellatartalom}%
    %
  \end{cellatartalom}}
\ora{\mitrakjonle}{18}{5}{8}
  {\begin{cellatartalom}%
    %
  \end{cellatartalom}}
\ora{\mitrakjonle}{18}{5}{9}
  {\begin{cellatartalom}%
    %
  \end{cellatartalom}}
\ora{\mitrakjonle}{19}{1}{0}
  {\begin{cellatartalom}%
     \oraresz{Kerling}{t{\" o}}{F}{II.50}{6}%%
  \end{cellatartalom}}
\ora{\mitrakjonle}{19}{1}{1}
  {\begin{cellatartalom}%
     \oraresz{Sch{\H o}n}{fi}{-}{Fiz.}{6}%%
  \end{cellatartalom}}
\ora{\mitrakjonle}{19}{1}{2}
  {\begin{cellatartalom}%
     \oraresz{Kudella}{fr}{fr}{Fr.}{1}%
\\ \oraresz{Sum}{sp}{sp1}{Sp.}{1}%
\\ \oraresz{Csuka}{ol}{ol3}{II.40}{1}%
\\ \oraresz{G{\' a}l}{n{\' e}}{n{\' e}2}{N.}{1}%
\\ \oraresz{Hal{\' a}sz}{n{\' e}}{n{\' e}4}{Kt.}{1}%
\\ \oraresz{S{\' o}tin{\' e}}{n{\' e}}{n{\' e}3}{F.18}{1}%%
  \end{cellatartalom}}
\ora{\mitrakjonle}{19}{1}{3}
  {\begin{cellatartalom}%
     \oraresz{Both}{i}{A}{I.33}{3}%
\\ \oraresz{S{\" u}le}{i}{F}{F.8}{3}%%
  \end{cellatartalom}}
\ora{\mitrakjonle}{19}{1}{4}
  {\begin{cellatartalom}%
     \oraresz{Both}{ny}{A}{I.33}{3}%
\\ \oraresz{S{\" u}le}{ny}{F}{F.8}{3}%%
  \end{cellatartalom}}
\ora{\mitrakjonle}{19}{1}{5}
  {\begin{cellatartalom}%
     \oraresz{Tandory}{f{\" o}}{-}{I.34}{6}%%
  \end{cellatartalom}}
\ora{\mitrakjonle}{19}{1}{6}
  {\begin{cellatartalom}%
     \oraresz{Furi{\' a}k}{te}{F}{Tt.}{3}%
\\ \oraresz{Vas E.}{te}{L}{Tt.}{3}%%
  \end{cellatartalom}}
\ora{\mitrakjonle}{19}{1}{7}
  {\begin{cellatartalom}%
     \oraresz{Barth{\' a}n{\' e}}{zt}{zt}{{\' E}.}{1.2}%
\\ \oraresz{Furi{\' a}k}{te}{F}{Tt.}{1.2}%
\\ \oraresz{Vas E.}{te}{L}{Tt.}{1.2}%
\\ \oraresz{P{\' o}k}{zt}{{\' a}g}{R.}{1.2}%
\\ \oraresz{Juh{\' a}sz A.}{a}{oktv}{A.1}{1.2}%%
  \end{cellatartalom}}
\ora{\mitrakjonle}{19}{1}{8}
  {\begin{cellatartalom}%
     \oraresz{Moln{\' a}r}{inf}{inf}{Szt.2}{1.5}%
\\ \oraresz{Barth{\' a}n{\' e}}{zt}{zt}{{\' E}.}{1.5}%
\\ \oraresz{P{\' o}k}{zt}{{\' a}g}{R.}{1.5}%
\\ \oraresz{Juh{\' a}sz A.}{a}{oktv}{A.1}{1.5}%%
  \end{cellatartalom}}
\ora{\mitrakjonle}{19}{1}{9}
  {\begin{cellatartalom}%
    %
  \end{cellatartalom}}
\ora{\mitrakjonle}{19}{2}{0}
  {\begin{cellatartalom}%
     \oraresz{S{\" u}le}{i}{F}{F.12}{3}%
\\ \oraresz{Mester}{mt}{F}{I.39}{3}%%
  \end{cellatartalom}}
\ora{\mitrakjonle}{19}{2}{1}
  {\begin{cellatartalom}%
     \oraresz{S{\" u}le}{i}{F}{F.12}{3}%
\\ \oraresz{Mester}{mt}{F}{I.39}{3}%%
  \end{cellatartalom}}
\ora{\mitrakjonle}{19}{2}{2}
  {\begin{cellatartalom}%
     \oraresz{Kerling}{t{\" o}}{F}{I.34}{3}%
\\ \oraresz{Adamis}{t{\" o}}{A}{II.51}{3}%%
  \end{cellatartalom}}
\ora{\mitrakjonle}{19}{2}{3}
  {\begin{cellatartalom}%
     \oraresz{Moln{\' a}r}{inf}{inf}{Szt.2}{3}%
\\ \oraresz{P.-P{\' a}sztor}{m{\' e}}{tk}{Tk.}{3}%%
  \end{cellatartalom}}
\ora{\mitrakjonle}{19}{2}{4}
  {\begin{cellatartalom}%
     \oraresz{Husz{\' a}r}{a}{an7}{I.34}{3}%
\\ \oraresz{T{\' o}th}{a}{an8}{II.41}{3}%%
  \end{cellatartalom}}
\ora{\mitrakjonle}{19}{2}{5}
  {\begin{cellatartalom}%
     \oraresz{Husz{\' a}r}{a}{an7}{I.34}{3}%
\\ \oraresz{T{\' o}th}{a}{an8}{II.41}{3}%%
  \end{cellatartalom}}
\ora{\mitrakjonle}{19}{2}{6}
  {\begin{cellatartalom}%
     \oraresz{Furi{\' a}k}{te}{F}{Tt.}{3}%
\\ \oraresz{Vas E.}{te}{L}{Tt.}{3}%%
  \end{cellatartalom}}
\ora{\mitrakjonle}{19}{2}{7}
  {\begin{cellatartalom}%
     \oraresz{Barth{\' a}n{\' e}}{zt}{zt}{{\' E}.}{6}%%
  \end{cellatartalom}}
\ora{\mitrakjonle}{19}{2}{8}
  {\begin{cellatartalom}%
     \oraresz{Barth{\' a}n{\' e}}{zt}{zt}{{\' E}.}{6}%%
  \end{cellatartalom}}
\ora{\mitrakjonle}{19}{2}{9}
  {\begin{cellatartalom}%
    %
  \end{cellatartalom}}
\ora{\mitrakjonle}{19}{3}{0}
  {\begin{cellatartalom}%
     \oraresz{S{\" u}le}{i}{F}{I.34}{6}%%
  \end{cellatartalom}}
\ora{\mitrakjonle}{19}{3}{1}
  {\begin{cellatartalom}%
     \oraresz{Kerling}{t{\" o}}{F}{F.8}{3}%
\\ \oraresz{Adamis}{t{\" o}}{A}{II.50}{3}%%
  \end{cellatartalom}}
\ora{\mitrakjonle}{19}{3}{2}
  {\begin{cellatartalom}%
     \oraresz{Krasny{\' a}nszki}{bi}{-}{B.ea}{6}%%
  \end{cellatartalom}}
\ora{\mitrakjonle}{19}{3}{3}
  {\begin{cellatartalom}%
     \oraresz{Moln{\' a}r}{inf}{inf}{Szt.2}{3}%
\\ \oraresz{P.-P{\' a}sztor}{m{\' e}}{tk}{Tk.}{3}%%
  \end{cellatartalom}}
\ora{\mitrakjonle}{19}{3}{4}
  {\begin{cellatartalom}%
     \oraresz{Moln{\' a}r}{inf}{inf}{Szt.2}{3}%
\\ \oraresz{P{\' o}k}{m{\' e}}{tk}{Tk.}{3}%%
  \end{cellatartalom}}
\ora{\mitrakjonle}{19}{3}{5}
  {\begin{cellatartalom}%
     \oraresz{Husz{\' a}r}{a}{an7}{{\' E}.}{3}%
\\ \oraresz{T{\' o}th}{a}{an8}{A.3}{3}%%
  \end{cellatartalom}}
\ora{\mitrakjonle}{19}{3}{6}
  {\begin{cellatartalom}%
     \oraresz{T{\' o}th}{of}{-}{I.33}{6}%%
  \end{cellatartalom}}
\ora{\mitrakjonle}{19}{3}{7}
  {\begin{cellatartalom}%
     \oraresz{P{\' o}k}{zt}{{\' a}g}{II.51}{3}%
\\ \oraresz{F{\" o}ldesi}{inf}{F}{Szt.2}{3}%%
  \end{cellatartalom}}
\ora{\mitrakjonle}{19}{3}{8}
  {\begin{cellatartalom}%
     \oraresz{P{\' o}k}{zt}{{\' a}g}{II.51}{3}%
\\ \oraresz{F{\" o}ldesi}{inf}{F}{Szt.2}{3}%%
  \end{cellatartalom}}
\ora{\mitrakjonle}{19}{3}{9}
  {\begin{cellatartalom}%
    %
  \end{cellatartalom}}
\ora{\mitrakjonle}{19}{4}{0}
  {\begin{cellatartalom}%
     \oraresz{Horv{\' a}th}{fi}{F}{Fiz.}{3}%
\\ \oraresz{Kerling}{t{\" o}}{F}{II.50}{3}%%
  \end{cellatartalom}}
\ora{\mitrakjonle}{19}{4}{1}
  {\begin{cellatartalom}%
     \oraresz{Horv{\' a}th}{fi}{F}{Fiz.}{3}%
\\ \oraresz{Kerling}{t{\" o}}{F}{II.50}{3}%%
  \end{cellatartalom}}
\ora{\mitrakjonle}{19}{4}{2}
  {\begin{cellatartalom}%
     \oraresz{Both}{i}{A}{I.39}{3}%
\\ \oraresz{S{\" u}le}{i}{F}{I.30}{3}%%
  \end{cellatartalom}}
\ora{\mitrakjonle}{19}{4}{3}
  {\begin{cellatartalom}%
     \oraresz{Varga Be.}{mt}{A}{F.16}{3}%
\\ \oraresz{Mester}{mt}{F}{II.48}{3}%%
  \end{cellatartalom}}
\ora{\mitrakjonle}{19}{4}{4}
  {\begin{cellatartalom}%
     \oraresz{Varga Be.}{mt}{A}{F.16}{3}%
\\ \oraresz{Mester}{mt}{F}{II.48}{3}%%
  \end{cellatartalom}}
\ora{\mitrakjonle}{19}{4}{5}
  {\begin{cellatartalom}%
     \oraresz{Husz{\' a}r}{a}{an7}{F.18}{3}%
\\ \oraresz{T{\' o}th}{a}{an8}{I.34}{3}%%
  \end{cellatartalom}}
\ora{\mitrakjonle}{19}{4}{6}
  {\begin{cellatartalom}%
     \oraresz{Kudella}{fr}{fr}{Fr.}{1}%
\\ \oraresz{Sum}{sp}{sp1}{Sp.}{1}%
\\ \oraresz{Csuka}{ol}{ol3}{L.}{1}%
\\ \oraresz{G{\' a}l}{n{\' e}}{n{\' e}2}{II.40}{1}%
\\ \oraresz{Hal{\' a}sz}{n{\' e}}{n{\' e}4}{II.41}{1}%
\\ \oraresz{S{\' o}tin{\' e}}{n{\' e}}{n{\' e}3}{{\' E}.}{1}%%
  \end{cellatartalom}}
\ora{\mitrakjonle}{19}{4}{7}
  {\begin{cellatartalom}%
     \oraresz{Papp L.}{m{\' e}}{tk}{Tk.}{1.5}%
\\ \oraresz{Barth{\' a}n{\' e}}{{\' e}k}{{\' e}k}{{\' E}.}{1.5}%
\\ \oraresz{P{\' o}k}{zt}{m{\H u}t}{II.48}{1.5}%
\\ \oraresz{Balanyi}{zt}{kgy}{R.}{1.5}%%
  \end{cellatartalom}}
\ora{\mitrakjonle}{19}{4}{8}
  {\begin{cellatartalom}%
     \oraresz{Barth{\' a}n{\' e}}{{\' e}k}{{\' e}k}{{\' E}.}{2}%
\\ \oraresz{P{\' o}k}{zt}{m{\H u}t}{II.48}{2}%
\\ \oraresz{Balanyi}{zt}{kgy}{R.}{2}%%
  \end{cellatartalom}}
\ora{\mitrakjonle}{19}{4}{9}
  {\begin{cellatartalom}%
    %
  \end{cellatartalom}}
\ora{\mitrakjonle}{19}{5}{0}
  {\begin{cellatartalom}%
     \oraresz{Moln{\' a}r}{inf}{inf}{Szt.2}{6}%%
  \end{cellatartalom}}
\ora{\mitrakjonle}{19}{5}{1}
  {\begin{cellatartalom}%
     \oraresz{Kudella}{fr}{fr}{Fr.}{1}%
\\ \oraresz{Sum}{sp}{sp1}{Sp.}{1}%
\\ \oraresz{Csuka}{ol}{ol3}{I.32}{1}%
\\ \oraresz{G{\' a}l}{n{\' e}}{n{\' e}2}{I.33}{1}%
\\ \oraresz{Hal{\' a}sz}{n{\' e}}{n{\' e}4}{II.41}{1}%
\\ \oraresz{S{\' o}tin{\' e}}{n{\' e}}{n{\' e}3}{{\' E}.}{1}%%
  \end{cellatartalom}}
\ora{\mitrakjonle}{19}{5}{2}
  {\begin{cellatartalom}%
     \oraresz{Krasny{\' a}nszki}{bi}{-}{K.ea}{6}%%
  \end{cellatartalom}}
\ora{\mitrakjonle}{19}{5}{3}
  {\begin{cellatartalom}%
     \oraresz{Sch{\H o}n}{fi}{-}{Fiz.}{6}%%
  \end{cellatartalom}}
\ora{\mitrakjonle}{19}{5}{4}
  {\begin{cellatartalom}%
     \oraresz{Tandory}{f{\" o}}{-}{I.39}{6}%%
  \end{cellatartalom}}
\ora{\mitrakjonle}{19}{5}{5}
  {\begin{cellatartalom}%
     \oraresz{Varga Be.}{mt}{A}{I.39}{3}%
\\ \oraresz{Mester}{mt}{F}{I.33}{3}%%
  \end{cellatartalom}}
\ora{\mitrakjonle}{19}{5}{6}
  {\begin{cellatartalom}%
     \oraresz{Husz{\' a}r}{a}{an7}{A.2}{3}%
\\ \oraresz{T{\' o}th}{a}{an8}{A.3}{3}%%
  \end{cellatartalom}}
\ora{\mitrakjonle}{19}{5}{7}
  {\begin{cellatartalom}%
     \oraresz{Cserv{\' a}k}{inf}{F}{Szt.1}{6}%%
  \end{cellatartalom}}
\ora{\mitrakjonle}{19}{5}{8}
  {\begin{cellatartalom}%
     \oraresz{Cserv{\' a}k}{inf}{F}{Szt.1}{6}%%
  \end{cellatartalom}}
\ora{\mitrakjonle}{19}{5}{9}
  {\begin{cellatartalom}%
    %
  \end{cellatartalom}}
\ora{\mitrakjonle}{20}{1}{0}
  {\begin{cellatartalom}%
    %
  \end{cellatartalom}}
\ora{\mitrakjonle}{20}{1}{1}
  {\begin{cellatartalom}%
    %
  \end{cellatartalom}}
\ora{\mitrakjonle}{20}{1}{2}
  {\begin{cellatartalom}%
    %
  \end{cellatartalom}}
\ora{\mitrakjonle}{20}{1}{3}
  {\begin{cellatartalom}%
    %
  \end{cellatartalom}}
\ora{\mitrakjonle}{20}{1}{4}
  {\begin{cellatartalom}%
    %
  \end{cellatartalom}}
\ora{\mitrakjonle}{20}{1}{5}
  {\begin{cellatartalom}%
    %
  \end{cellatartalom}}
\ora{\mitrakjonle}{20}{1}{6}
  {\begin{cellatartalom}%
    %
  \end{cellatartalom}}
\ora{\mitrakjonle}{20}{1}{7}
  {\begin{cellatartalom}%
     \oraresz{Juh{\' a}sz A.}{a}{oktv}{A.1}{1}%%
  \end{cellatartalom}}
\ora{\mitrakjonle}{20}{1}{8}
  {\begin{cellatartalom}%
     \oraresz{Juh{\' a}sz A.}{a}{oktv}{A.1}{1}%%
  \end{cellatartalom}}
\ora{\mitrakjonle}{20}{1}{9}
  {\begin{cellatartalom}%
    %
  \end{cellatartalom}}
\ora{\mitrakjonle}{20}{2}{0}
  {\begin{cellatartalom}%
    %
  \end{cellatartalom}}
\ora{\mitrakjonle}{20}{2}{1}
  {\begin{cellatartalom}%
    %
  \end{cellatartalom}}
\ora{\mitrakjonle}{20}{2}{2}
  {\begin{cellatartalom}%
    %
  \end{cellatartalom}}
\ora{\mitrakjonle}{20}{2}{3}
  {\begin{cellatartalom}%
    %
  \end{cellatartalom}}
\ora{\mitrakjonle}{20}{2}{4}
  {\begin{cellatartalom}%
    %
  \end{cellatartalom}}
\ora{\mitrakjonle}{20}{2}{5}
  {\begin{cellatartalom}%
    %
  \end{cellatartalom}}
\ora{\mitrakjonle}{20}{2}{6}
  {\begin{cellatartalom}%
    %
  \end{cellatartalom}}
\ora{\mitrakjonle}{20}{2}{7}
  {\begin{cellatartalom}%
    %
  \end{cellatartalom}}
\ora{\mitrakjonle}{20}{2}{8}
  {\begin{cellatartalom}%
     \oraresz{Krajny{\' a}k}{inf}{szak}{Szt.1}{1}%%
  \end{cellatartalom}}
\ora{\mitrakjonle}{20}{2}{9}
  {\begin{cellatartalom}%
    %
  \end{cellatartalom}}
\ora{\mitrakjonle}{20}{3}{0}
  {\begin{cellatartalom}%
    %
  \end{cellatartalom}}
\ora{\mitrakjonle}{20}{3}{1}
  {\begin{cellatartalom}%
    %
  \end{cellatartalom}}
\ora{\mitrakjonle}{20}{3}{2}
  {\begin{cellatartalom}%
    %
  \end{cellatartalom}}
\ora{\mitrakjonle}{20}{3}{3}
  {\begin{cellatartalom}%
    %
  \end{cellatartalom}}
\ora{\mitrakjonle}{20}{3}{4}
  {\begin{cellatartalom}%
    %
  \end{cellatartalom}}
\ora{\mitrakjonle}{20}{3}{5}
  {\begin{cellatartalom}%
    %
  \end{cellatartalom}}
\ora{\mitrakjonle}{20}{3}{6}
  {\begin{cellatartalom}%
    %
  \end{cellatartalom}}
\ora{\mitrakjonle}{20}{3}{7}
  {\begin{cellatartalom}%
    %
  \end{cellatartalom}}
\ora{\mitrakjonle}{20}{3}{8}
  {\begin{cellatartalom}%
    %
  \end{cellatartalom}}
\ora{\mitrakjonle}{20}{3}{9}
  {\begin{cellatartalom}%
    %
  \end{cellatartalom}}
\ora{\mitrakjonle}{20}{4}{0}
  {\begin{cellatartalom}%
    %
  \end{cellatartalom}}
\ora{\mitrakjonle}{20}{4}{1}
  {\begin{cellatartalom}%
    %
  \end{cellatartalom}}
\ora{\mitrakjonle}{20}{4}{2}
  {\begin{cellatartalom}%
    %
  \end{cellatartalom}}
\ora{\mitrakjonle}{20}{4}{3}
  {\begin{cellatartalom}%
    %
  \end{cellatartalom}}
\ora{\mitrakjonle}{20}{4}{4}
  {\begin{cellatartalom}%
    %
  \end{cellatartalom}}
\ora{\mitrakjonle}{20}{4}{5}
  {\begin{cellatartalom}%
    %
  \end{cellatartalom}}
\ora{\mitrakjonle}{20}{4}{6}
  {\begin{cellatartalom}%
    %
  \end{cellatartalom}}
\ora{\mitrakjonle}{20}{4}{7}
  {\begin{cellatartalom}%
    %
  \end{cellatartalom}}
\ora{\mitrakjonle}{20}{4}{8}
  {\begin{cellatartalom}%
    %
  \end{cellatartalom}}
\ora{\mitrakjonle}{20}{4}{9}
  {\begin{cellatartalom}%
    %
  \end{cellatartalom}}
\ora{\mitrakjonle}{20}{5}{0}
  {\begin{cellatartalom}%
    %
  \end{cellatartalom}}
\ora{\mitrakjonle}{20}{5}{1}
  {\begin{cellatartalom}%
    %
  \end{cellatartalom}}
\ora{\mitrakjonle}{20}{5}{2}
  {\begin{cellatartalom}%
    %
  \end{cellatartalom}}
\ora{\mitrakjonle}{20}{5}{3}
  {\begin{cellatartalom}%
    %
  \end{cellatartalom}}
\ora{\mitrakjonle}{20}{5}{4}
  {\begin{cellatartalom}%
    %
  \end{cellatartalom}}
\ora{\mitrakjonle}{20}{5}{5}
  {\begin{cellatartalom}%
    %
  \end{cellatartalom}}
\ora{\mitrakjonle}{20}{5}{6}
  {\begin{cellatartalom}%
    %
  \end{cellatartalom}}
\ora{\mitrakjonle}{20}{5}{7}
  {\begin{cellatartalom}%
    %
  \end{cellatartalom}}
\ora{\mitrakjonle}{20}{5}{8}
  {\begin{cellatartalom}%
    %
  \end{cellatartalom}}
\ora{\mitrakjonle}{20}{5}{9}
  {\begin{cellatartalom}%
    %
  \end{cellatartalom}}
\ora{\mitrakjonle}{21}{1}{0}
  {\begin{cellatartalom}%
    %
  \end{cellatartalom}}
\ora{\mitrakjonle}{21}{1}{1}
  {\begin{cellatartalom}%
    %
  \end{cellatartalom}}
\ora{\mitrakjonle}{21}{1}{2}
  {\begin{cellatartalom}%
    %
  \end{cellatartalom}}
\ora{\mitrakjonle}{21}{1}{3}
  {\begin{cellatartalom}%
    %
  \end{cellatartalom}}
\ora{\mitrakjonle}{21}{1}{4}
  {\begin{cellatartalom}%
    %
  \end{cellatartalom}}
\ora{\mitrakjonle}{21}{1}{5}
  {\begin{cellatartalom}%
    %
  \end{cellatartalom}}
\ora{\mitrakjonle}{21}{1}{6}
  {\begin{cellatartalom}%
    %
  \end{cellatartalom}}
\ora{\mitrakjonle}{21}{1}{7}
  {\begin{cellatartalom}%
     \oraresz{Juh{\' a}sz A.}{a}{oktv}{A.1}{1}%%
  \end{cellatartalom}}
\ora{\mitrakjonle}{21}{1}{8}
  {\begin{cellatartalom}%
     \oraresz{Juh{\' a}sz A.}{a}{oktv}{A.1}{1}%%
  \end{cellatartalom}}
\ora{\mitrakjonle}{21}{1}{9}
  {\begin{cellatartalom}%
    %
  \end{cellatartalom}}
\ora{\mitrakjonle}{21}{2}{0}
  {\begin{cellatartalom}%
    %
  \end{cellatartalom}}
\ora{\mitrakjonle}{21}{2}{1}
  {\begin{cellatartalom}%
    %
  \end{cellatartalom}}
\ora{\mitrakjonle}{21}{2}{2}
  {\begin{cellatartalom}%
    %
  \end{cellatartalom}}
\ora{\mitrakjonle}{21}{2}{3}
  {\begin{cellatartalom}%
    %
  \end{cellatartalom}}
\ora{\mitrakjonle}{21}{2}{4}
  {\begin{cellatartalom}%
    %
  \end{cellatartalom}}
\ora{\mitrakjonle}{21}{2}{5}
  {\begin{cellatartalom}%
    %
  \end{cellatartalom}}
\ora{\mitrakjonle}{21}{2}{6}
  {\begin{cellatartalom}%
    %
  \end{cellatartalom}}
\ora{\mitrakjonle}{21}{2}{7}
  {\begin{cellatartalom}%
    %
  \end{cellatartalom}}
\ora{\mitrakjonle}{21}{2}{8}
  {\begin{cellatartalom}%
    %
  \end{cellatartalom}}
\ora{\mitrakjonle}{21}{2}{9}
  {\begin{cellatartalom}%
    %
  \end{cellatartalom}}
\ora{\mitrakjonle}{21}{3}{0}
  {\begin{cellatartalom}%
    %
  \end{cellatartalom}}
\ora{\mitrakjonle}{21}{3}{1}
  {\begin{cellatartalom}%
    %
  \end{cellatartalom}}
\ora{\mitrakjonle}{21}{3}{2}
  {\begin{cellatartalom}%
    %
  \end{cellatartalom}}
\ora{\mitrakjonle}{21}{3}{3}
  {\begin{cellatartalom}%
    %
  \end{cellatartalom}}
\ora{\mitrakjonle}{21}{3}{4}
  {\begin{cellatartalom}%
    %
  \end{cellatartalom}}
\ora{\mitrakjonle}{21}{3}{5}
  {\begin{cellatartalom}%
    %
  \end{cellatartalom}}
\ora{\mitrakjonle}{21}{3}{6}
  {\begin{cellatartalom}%
    %
  \end{cellatartalom}}
\ora{\mitrakjonle}{21}{3}{7}
  {\begin{cellatartalom}%
    %
  \end{cellatartalom}}
\ora{\mitrakjonle}{21}{3}{8}
  {\begin{cellatartalom}%
    %
  \end{cellatartalom}}
\ora{\mitrakjonle}{21}{3}{9}
  {\begin{cellatartalom}%
    %
  \end{cellatartalom}}
\ora{\mitrakjonle}{21}{4}{0}
  {\begin{cellatartalom}%
    %
  \end{cellatartalom}}
\ora{\mitrakjonle}{21}{4}{1}
  {\begin{cellatartalom}%
    %
  \end{cellatartalom}}
\ora{\mitrakjonle}{21}{4}{2}
  {\begin{cellatartalom}%
    %
  \end{cellatartalom}}
\ora{\mitrakjonle}{21}{4}{3}
  {\begin{cellatartalom}%
    %
  \end{cellatartalom}}
\ora{\mitrakjonle}{21}{4}{4}
  {\begin{cellatartalom}%
    %
  \end{cellatartalom}}
\ora{\mitrakjonle}{21}{4}{5}
  {\begin{cellatartalom}%
    %
  \end{cellatartalom}}
\ora{\mitrakjonle}{21}{4}{6}
  {\begin{cellatartalom}%
    %
  \end{cellatartalom}}
\ora{\mitrakjonle}{21}{4}{7}
  {\begin{cellatartalom}%
    %
  \end{cellatartalom}}
\ora{\mitrakjonle}{21}{4}{8}
  {\begin{cellatartalom}%
    %
  \end{cellatartalom}}
\ora{\mitrakjonle}{21}{4}{9}
  {\begin{cellatartalom}%
    %
  \end{cellatartalom}}
\ora{\mitrakjonle}{21}{5}{0}
  {\begin{cellatartalom}%
    %
  \end{cellatartalom}}
\ora{\mitrakjonle}{21}{5}{1}
  {\begin{cellatartalom}%
    %
  \end{cellatartalom}}
\ora{\mitrakjonle}{21}{5}{2}
  {\begin{cellatartalom}%
    %
  \end{cellatartalom}}
\ora{\mitrakjonle}{21}{5}{3}
  {\begin{cellatartalom}%
    %
  \end{cellatartalom}}
\ora{\mitrakjonle}{21}{5}{4}
  {\begin{cellatartalom}%
    %
  \end{cellatartalom}}
\ora{\mitrakjonle}{21}{5}{5}
  {\begin{cellatartalom}%
    %
  \end{cellatartalom}}
\ora{\mitrakjonle}{21}{5}{6}
  {\begin{cellatartalom}%
    %
  \end{cellatartalom}}
\ora{\mitrakjonle}{21}{5}{7}
  {\begin{cellatartalom}%
    %
  \end{cellatartalom}}
\ora{\mitrakjonle}{21}{5}{8}
  {\begin{cellatartalom}%
    %
  \end{cellatartalom}}
\ora{\mitrakjonle}{21}{5}{9}
  {\begin{cellatartalom}%
    %
  \end{cellatartalom}}
\ora{\mitrakjonle}{22}{1}{0}
  {\begin{cellatartalom}%
    %
  \end{cellatartalom}}
\ora{\mitrakjonle}{22}{1}{1}
  {\begin{cellatartalom}%
    %
  \end{cellatartalom}}
\ora{\mitrakjonle}{22}{1}{2}
  {\begin{cellatartalom}%
    %
  \end{cellatartalom}}
\ora{\mitrakjonle}{22}{1}{3}
  {\begin{cellatartalom}%
    %
  \end{cellatartalom}}
\ora{\mitrakjonle}{22}{1}{4}
  {\begin{cellatartalom}%
    %
  \end{cellatartalom}}
\ora{\mitrakjonle}{22}{1}{5}
  {\begin{cellatartalom}%
    %
  \end{cellatartalom}}
\ora{\mitrakjonle}{22}{1}{6}
  {\begin{cellatartalom}%
    %
  \end{cellatartalom}}
\ora{\mitrakjonle}{22}{1}{7}
  {\begin{cellatartalom}%
     \oraresz{Juh{\' a}sz A.}{a}{oktv}{A.1}{1}%%
  \end{cellatartalom}}
\ora{\mitrakjonle}{22}{1}{8}
  {\begin{cellatartalom}%
     \oraresz{Juh{\' a}sz A.}{a}{oktv}{A.1}{1}%%
  \end{cellatartalom}}
\ora{\mitrakjonle}{22}{1}{9}
  {\begin{cellatartalom}%
    %
  \end{cellatartalom}}
\ora{\mitrakjonle}{22}{2}{0}
  {\begin{cellatartalom}%
    %
  \end{cellatartalom}}
\ora{\mitrakjonle}{22}{2}{1}
  {\begin{cellatartalom}%
    %
  \end{cellatartalom}}
\ora{\mitrakjonle}{22}{2}{2}
  {\begin{cellatartalom}%
    %
  \end{cellatartalom}}
\ora{\mitrakjonle}{22}{2}{3}
  {\begin{cellatartalom}%
    %
  \end{cellatartalom}}
\ora{\mitrakjonle}{22}{2}{4}
  {\begin{cellatartalom}%
    %
  \end{cellatartalom}}
\ora{\mitrakjonle}{22}{2}{5}
  {\begin{cellatartalom}%
    %
  \end{cellatartalom}}
\ora{\mitrakjonle}{22}{2}{6}
  {\begin{cellatartalom}%
    %
  \end{cellatartalom}}
\ora{\mitrakjonle}{22}{2}{7}
  {\begin{cellatartalom}%
    %
  \end{cellatartalom}}
\ora{\mitrakjonle}{22}{2}{8}
  {\begin{cellatartalom}%
    %
  \end{cellatartalom}}
\ora{\mitrakjonle}{22}{2}{9}
  {\begin{cellatartalom}%
    %
  \end{cellatartalom}}
\ora{\mitrakjonle}{22}{3}{0}
  {\begin{cellatartalom}%
    %
  \end{cellatartalom}}
\ora{\mitrakjonle}{22}{3}{1}
  {\begin{cellatartalom}%
    %
  \end{cellatartalom}}
\ora{\mitrakjonle}{22}{3}{2}
  {\begin{cellatartalom}%
    %
  \end{cellatartalom}}
\ora{\mitrakjonle}{22}{3}{3}
  {\begin{cellatartalom}%
    %
  \end{cellatartalom}}
\ora{\mitrakjonle}{22}{3}{4}
  {\begin{cellatartalom}%
    %
  \end{cellatartalom}}
\ora{\mitrakjonle}{22}{3}{5}
  {\begin{cellatartalom}%
     \oraresz{J.-Varga}{n{\' e}}{n{\' e}2}{II.41}{1}%%
  \end{cellatartalom}}
\ora{\mitrakjonle}{22}{3}{6}
  {\begin{cellatartalom}%
     \oraresz{J.-Varga}{n{\' e}}{n{\' e}2}{II.41}{1}%%
  \end{cellatartalom}}
\ora{\mitrakjonle}{22}{3}{7}
  {\begin{cellatartalom}%
    %
  \end{cellatartalom}}
\ora{\mitrakjonle}{22}{3}{8}
  {\begin{cellatartalom}%
    %
  \end{cellatartalom}}
\ora{\mitrakjonle}{22}{3}{9}
  {\begin{cellatartalom}%
    %
  \end{cellatartalom}}
\ora{\mitrakjonle}{22}{4}{0}
  {\begin{cellatartalom}%
    %
  \end{cellatartalom}}
\ora{\mitrakjonle}{22}{4}{1}
  {\begin{cellatartalom}%
    %
  \end{cellatartalom}}
\ora{\mitrakjonle}{22}{4}{2}
  {\begin{cellatartalom}%
    %
  \end{cellatartalom}}
\ora{\mitrakjonle}{22}{4}{3}
  {\begin{cellatartalom}%
     \oraresz{J.-Varga}{n{\' e}}{n{\' e}2}{F.18}{1}%%
  \end{cellatartalom}}
\ora{\mitrakjonle}{22}{4}{4}
  {\begin{cellatartalom}%
    %
  \end{cellatartalom}}
\ora{\mitrakjonle}{22}{4}{5}
  {\begin{cellatartalom}%
    %
  \end{cellatartalom}}
\ora{\mitrakjonle}{22}{4}{6}
  {\begin{cellatartalom}%
    %
  \end{cellatartalom}}
\ora{\mitrakjonle}{22}{4}{7}
  {\begin{cellatartalom}%
    %
  \end{cellatartalom}}
\ora{\mitrakjonle}{22}{4}{8}
  {\begin{cellatartalom}%
    %
  \end{cellatartalom}}
\ora{\mitrakjonle}{22}{4}{9}
  {\begin{cellatartalom}%
    %
  \end{cellatartalom}}
\ora{\mitrakjonle}{22}{5}{0}
  {\begin{cellatartalom}%
    %
  \end{cellatartalom}}
\ora{\mitrakjonle}{22}{5}{1}
  {\begin{cellatartalom}%
    %
  \end{cellatartalom}}
\ora{\mitrakjonle}{22}{5}{2}
  {\begin{cellatartalom}%
     \oraresz{J.-Varga}{n{\' e}}{n{\' e}2}{I.39}{1}%%
  \end{cellatartalom}}
\ora{\mitrakjonle}{22}{5}{3}
  {\begin{cellatartalom}%
    %
  \end{cellatartalom}}
\ora{\mitrakjonle}{22}{5}{4}
  {\begin{cellatartalom}%
    %
  \end{cellatartalom}}
\ora{\mitrakjonle}{22}{5}{5}
  {\begin{cellatartalom}%
    %
  \end{cellatartalom}}
\ora{\mitrakjonle}{22}{5}{6}
  {\begin{cellatartalom}%
    %
  \end{cellatartalom}}
\ora{\mitrakjonle}{22}{5}{7}
  {\begin{cellatartalom}%
    %
  \end{cellatartalom}}
\ora{\mitrakjonle}{22}{5}{8}
  {\begin{cellatartalom}%
    %
  \end{cellatartalom}}
\ora{\mitrakjonle}{22}{5}{9}
  {\begin{cellatartalom}%
    %
  \end{cellatartalom}}
\ora{\mitrakjonle}{23}{1}{0}
  {\begin{cellatartalom}%
     \oraresz{Nagy Sz.}{te}{L}{Tt.}{2}%
\\ \oraresz{Furi{\' a}k}{te}{F}{Tt.}{2}%%
  \end{cellatartalom}}
\ora{\mitrakjonle}{23}{1}{1}
  {\begin{cellatartalom}%
     \oraresz{Nagy Sz.}{te}{L}{Tt.}{2}%
\\ \oraresz{Furi{\' a}k}{te}{F}{Tt.}{2}%%
  \end{cellatartalom}}
\ora{\mitrakjonle}{23}{1}{2}
  {\begin{cellatartalom}%
    %
  \end{cellatartalom}}
\ora{\mitrakjonle}{23}{1}{3}
  {\begin{cellatartalom}%
    %
  \end{cellatartalom}}
\ora{\mitrakjonle}{23}{1}{4}
  {\begin{cellatartalom}%
    %
  \end{cellatartalom}}
\ora{\mitrakjonle}{23}{1}{5}
  {\begin{cellatartalom}%
    %
  \end{cellatartalom}}
\ora{\mitrakjonle}{23}{1}{6}
  {\begin{cellatartalom}%
    %
  \end{cellatartalom}}
\ora{\mitrakjonle}{23}{1}{7}
  {\begin{cellatartalom}%
     \oraresz{Juh{\' a}sz A.}{a}{oktv}{A.1}{4}%%
  \end{cellatartalom}}
\ora{\mitrakjonle}{23}{1}{8}
  {\begin{cellatartalom}%
     \oraresz{Juh{\' a}sz A.}{a}{oktv}{A.1}{4}%%
  \end{cellatartalom}}
\ora{\mitrakjonle}{23}{1}{9}
  {\begin{cellatartalom}%
    %
  \end{cellatartalom}}
\ora{\mitrakjonle}{23}{2}{0}
  {\begin{cellatartalom}%
    %
  \end{cellatartalom}}
\ora{\mitrakjonle}{23}{2}{1}
  {\begin{cellatartalom}%
    %
  \end{cellatartalom}}
\ora{\mitrakjonle}{23}{2}{2}
  {\begin{cellatartalom}%
    %
  \end{cellatartalom}}
\ora{\mitrakjonle}{23}{2}{3}
  {\begin{cellatartalom}%
    %
  \end{cellatartalom}}
\ora{\mitrakjonle}{23}{2}{4}
  {\begin{cellatartalom}%
    %
  \end{cellatartalom}}
\ora{\mitrakjonle}{23}{2}{5}
  {\begin{cellatartalom}%
    %
  \end{cellatartalom}}
\ora{\mitrakjonle}{23}{2}{6}
  {\begin{cellatartalom}%
    %
  \end{cellatartalom}}
\ora{\mitrakjonle}{23}{2}{7}
  {\begin{cellatartalom}%
    %
  \end{cellatartalom}}
\ora{\mitrakjonle}{23}{2}{8}
  {\begin{cellatartalom}%
    %
  \end{cellatartalom}}
\ora{\mitrakjonle}{23}{2}{9}
  {\begin{cellatartalom}%
    %
  \end{cellatartalom}}
\ora{\mitrakjonle}{23}{3}{0}
  {\begin{cellatartalom}%
    %
  \end{cellatartalom}}
\ora{\mitrakjonle}{23}{3}{1}
  {\begin{cellatartalom}%
    %
  \end{cellatartalom}}
\ora{\mitrakjonle}{23}{3}{2}
  {\begin{cellatartalom}%
    %
  \end{cellatartalom}}
\ora{\mitrakjonle}{23}{3}{3}
  {\begin{cellatartalom}%
    %
  \end{cellatartalom}}
\ora{\mitrakjonle}{23}{3}{4}
  {\begin{cellatartalom}%
    %
  \end{cellatartalom}}
\ora{\mitrakjonle}{23}{3}{5}
  {\begin{cellatartalom}%
     \oraresz{Dickerson}{a}{an3}{A.2}{1}%
\\ \oraresz{Liss{\' a}k}{a}{an2}{A.1}{1}%
\\ \oraresz{Csan{\' a}di}{a}{an1}{II.44}{1}%
\\ \oraresz{J.-Varga}{n{\' e}}{n{\' e}2}{II.41}{1}%%
  \end{cellatartalom}}
\ora{\mitrakjonle}{23}{3}{6}
  {\begin{cellatartalom}%
     \oraresz{Dickerson}{a}{an3}{A.2}{1}%
\\ \oraresz{Liss{\' a}k}{a}{an2}{A.1}{1}%
\\ \oraresz{Csan{\' a}di}{a}{an1}{II.44}{1}%
\\ \oraresz{J.-Varga}{n{\' e}}{n{\' e}2}{II.41}{1}%%
  \end{cellatartalom}}
\ora{\mitrakjonle}{23}{3}{7}
  {\begin{cellatartalom}%
    %
  \end{cellatartalom}}
\ora{\mitrakjonle}{23}{3}{8}
  {\begin{cellatartalom}%
    %
  \end{cellatartalom}}
\ora{\mitrakjonle}{23}{3}{9}
  {\begin{cellatartalom}%
    %
  \end{cellatartalom}}
\ora{\mitrakjonle}{23}{4}{0}
  {\begin{cellatartalom}%
    %
  \end{cellatartalom}}
\ora{\mitrakjonle}{23}{4}{1}
  {\begin{cellatartalom}%
    %
  \end{cellatartalom}}
\ora{\mitrakjonle}{23}{4}{2}
  {\begin{cellatartalom}%
    %
  \end{cellatartalom}}
\ora{\mitrakjonle}{23}{4}{3}
  {\begin{cellatartalom}%
     \oraresz{Husz{\' a}r}{a}{an3}{A.1}{1}%
\\ \oraresz{Liss{\' a}k}{a}{an2}{II.51}{1}%
\\ \oraresz{Csan{\' a}di}{a}{an1}{II.40}{1}%
\\ \oraresz{J.-Varga}{n{\' e}}{n{\' e}2}{F.18}{1}%%
  \end{cellatartalom}}
\ora{\mitrakjonle}{23}{4}{4}
  {\begin{cellatartalom}%
    %
  \end{cellatartalom}}
\ora{\mitrakjonle}{23}{4}{5}
  {\begin{cellatartalom}%
    %
  \end{cellatartalom}}
\ora{\mitrakjonle}{23}{4}{6}
  {\begin{cellatartalom}%
     \oraresz{Nagy Sz.}{te}{L}{Tt.}{2}%
\\ \oraresz{Furi{\' a}k}{te}{F}{Tt.}{2}%%
  \end{cellatartalom}}
\ora{\mitrakjonle}{23}{4}{7}
  {\begin{cellatartalom}%
    %
  \end{cellatartalom}}
\ora{\mitrakjonle}{23}{4}{8}
  {\begin{cellatartalom}%
    %
  \end{cellatartalom}}
\ora{\mitrakjonle}{23}{4}{9}
  {\begin{cellatartalom}%
    %
  \end{cellatartalom}}
\ora{\mitrakjonle}{23}{5}{0}
  {\begin{cellatartalom}%
    %
  \end{cellatartalom}}
\ora{\mitrakjonle}{23}{5}{1}
  {\begin{cellatartalom}%
    %
  \end{cellatartalom}}
\ora{\mitrakjonle}{23}{5}{2}
  {\begin{cellatartalom}%
     \oraresz{Dickerson}{a}{an3}{A.2}{1}%
\\ \oraresz{Liss{\' a}k}{a}{an2}{A.3}{1}%
\\ \oraresz{Csan{\' a}di}{a}{an1}{A.1}{1}%
\\ \oraresz{J.-Varga}{n{\' e}}{n{\' e}2}{I.39}{1}%%
  \end{cellatartalom}}
\ora{\mitrakjonle}{23}{5}{3}
  {\begin{cellatartalom}%
    %
  \end{cellatartalom}}
\ora{\mitrakjonle}{23}{5}{4}
  {\begin{cellatartalom}%
    %
  \end{cellatartalom}}
\ora{\mitrakjonle}{23}{5}{5}
  {\begin{cellatartalom}%
    %
  \end{cellatartalom}}
\ora{\mitrakjonle}{23}{5}{6}
  {\begin{cellatartalom}%
    %
  \end{cellatartalom}}
\ora{\mitrakjonle}{23}{5}{7}
  {\begin{cellatartalom}%
    %
  \end{cellatartalom}}
\ora{\mitrakjonle}{23}{5}{8}
  {\begin{cellatartalom}%
    %
  \end{cellatartalom}}
\ora{\mitrakjonle}{23}{5}{9}
  {\begin{cellatartalom}%
    %
  \end{cellatartalom}}
\ora{\mitrakjonle}{24}{1}{0}
  {\begin{cellatartalom}%
    %
  \end{cellatartalom}}
\ora{\mitrakjonle}{24}{1}{1}
  {\begin{cellatartalom}%
    %
  \end{cellatartalom}}
\ora{\mitrakjonle}{24}{1}{2}
  {\begin{cellatartalom}%
    %
  \end{cellatartalom}}
\ora{\mitrakjonle}{24}{1}{3}
  {\begin{cellatartalom}%
    %
  \end{cellatartalom}}
\ora{\mitrakjonle}{24}{1}{4}
  {\begin{cellatartalom}%
    %
  \end{cellatartalom}}
\ora{\mitrakjonle}{24}{1}{5}
  {\begin{cellatartalom}%
    %
  \end{cellatartalom}}
\ora{\mitrakjonle}{24}{1}{6}
  {\begin{cellatartalom}%
    %
  \end{cellatartalom}}
\ora{\mitrakjonle}{24}{1}{7}
  {\begin{cellatartalom}%
     \oraresz{Juh{\' a}sz A.}{a}{oktv}{A.1}{1}%%
  \end{cellatartalom}}
\ora{\mitrakjonle}{24}{1}{8}
  {\begin{cellatartalom}%
     \oraresz{Juh{\' a}sz A.}{a}{oktv}{A.1}{1}%%
  \end{cellatartalom}}
\ora{\mitrakjonle}{24}{1}{9}
  {\begin{cellatartalom}%
    %
  \end{cellatartalom}}
\ora{\mitrakjonle}{24}{2}{0}
  {\begin{cellatartalom}%
    %
  \end{cellatartalom}}
\ora{\mitrakjonle}{24}{2}{1}
  {\begin{cellatartalom}%
    %
  \end{cellatartalom}}
\ora{\mitrakjonle}{24}{2}{2}
  {\begin{cellatartalom}%
    %
  \end{cellatartalom}}
\ora{\mitrakjonle}{24}{2}{3}
  {\begin{cellatartalom}%
    %
  \end{cellatartalom}}
\ora{\mitrakjonle}{24}{2}{4}
  {\begin{cellatartalom}%
    %
  \end{cellatartalom}}
\ora{\mitrakjonle}{24}{2}{5}
  {\begin{cellatartalom}%
    %
  \end{cellatartalom}}
\ora{\mitrakjonle}{24}{2}{6}
  {\begin{cellatartalom}%
    %
  \end{cellatartalom}}
\ora{\mitrakjonle}{24}{2}{7}
  {\begin{cellatartalom}%
    %
  \end{cellatartalom}}
\ora{\mitrakjonle}{24}{2}{8}
  {\begin{cellatartalom}%
    %
  \end{cellatartalom}}
\ora{\mitrakjonle}{24}{2}{9}
  {\begin{cellatartalom}%
    %
  \end{cellatartalom}}
\ora{\mitrakjonle}{24}{3}{0}
  {\begin{cellatartalom}%
    %
  \end{cellatartalom}}
\ora{\mitrakjonle}{24}{3}{1}
  {\begin{cellatartalom}%
    %
  \end{cellatartalom}}
\ora{\mitrakjonle}{24}{3}{2}
  {\begin{cellatartalom}%
    %
  \end{cellatartalom}}
\ora{\mitrakjonle}{24}{3}{3}
  {\begin{cellatartalom}%
    %
  \end{cellatartalom}}
\ora{\mitrakjonle}{24}{3}{4}
  {\begin{cellatartalom}%
    %
  \end{cellatartalom}}
\ora{\mitrakjonle}{24}{3}{5}
  {\begin{cellatartalom}%
    %
  \end{cellatartalom}}
\ora{\mitrakjonle}{24}{3}{6}
  {\begin{cellatartalom}%
    %
  \end{cellatartalom}}
\ora{\mitrakjonle}{24}{3}{7}
  {\begin{cellatartalom}%
    %
  \end{cellatartalom}}
\ora{\mitrakjonle}{24}{3}{8}
  {\begin{cellatartalom}%
    %
  \end{cellatartalom}}
\ora{\mitrakjonle}{24}{3}{9}
  {\begin{cellatartalom}%
    %
  \end{cellatartalom}}
\ora{\mitrakjonle}{24}{4}{0}
  {\begin{cellatartalom}%
    %
  \end{cellatartalom}}
\ora{\mitrakjonle}{24}{4}{1}
  {\begin{cellatartalom}%
    %
  \end{cellatartalom}}
\ora{\mitrakjonle}{24}{4}{2}
  {\begin{cellatartalom}%
    %
  \end{cellatartalom}}
\ora{\mitrakjonle}{24}{4}{3}
  {\begin{cellatartalom}%
    %
  \end{cellatartalom}}
\ora{\mitrakjonle}{24}{4}{4}
  {\begin{cellatartalom}%
    %
  \end{cellatartalom}}
\ora{\mitrakjonle}{24}{4}{5}
  {\begin{cellatartalom}%
    %
  \end{cellatartalom}}
\ora{\mitrakjonle}{24}{4}{6}
  {\begin{cellatartalom}%
    %
  \end{cellatartalom}}
\ora{\mitrakjonle}{24}{4}{7}
  {\begin{cellatartalom}%
    %
  \end{cellatartalom}}
\ora{\mitrakjonle}{24}{4}{8}
  {\begin{cellatartalom}%
    %
  \end{cellatartalom}}
\ora{\mitrakjonle}{24}{4}{9}
  {\begin{cellatartalom}%
    %
  \end{cellatartalom}}
\ora{\mitrakjonle}{24}{5}{0}
  {\begin{cellatartalom}%
    %
  \end{cellatartalom}}
\ora{\mitrakjonle}{24}{5}{1}
  {\begin{cellatartalom}%
    %
  \end{cellatartalom}}
\ora{\mitrakjonle}{24}{5}{2}
  {\begin{cellatartalom}%
    %
  \end{cellatartalom}}
\ora{\mitrakjonle}{24}{5}{3}
  {\begin{cellatartalom}%
    %
  \end{cellatartalom}}
\ora{\mitrakjonle}{24}{5}{4}
  {\begin{cellatartalom}%
    %
  \end{cellatartalom}}
\ora{\mitrakjonle}{24}{5}{5}
  {\begin{cellatartalom}%
    %
  \end{cellatartalom}}
\ora{\mitrakjonle}{24}{5}{6}
  {\begin{cellatartalom}%
    %
  \end{cellatartalom}}
\ora{\mitrakjonle}{24}{5}{7}
  {\begin{cellatartalom}%
    %
  \end{cellatartalom}}
\ora{\mitrakjonle}{24}{5}{8}
  {\begin{cellatartalom}%
    %
  \end{cellatartalom}}
\ora{\mitrakjonle}{24}{5}{9}
  {\begin{cellatartalom}%
    %
  \end{cellatartalom}}
\ora{\mitrakjonle}{25}{1}{0}
  {\begin{cellatartalom}%
    %
  \end{cellatartalom}}
\ora{\mitrakjonle}{25}{1}{1}
  {\begin{cellatartalom}%
    %
  \end{cellatartalom}}
\ora{\mitrakjonle}{25}{1}{2}
  {\begin{cellatartalom}%
    %
  \end{cellatartalom}}
\ora{\mitrakjonle}{25}{1}{3}
  {\begin{cellatartalom}%
    %
  \end{cellatartalom}}
\ora{\mitrakjonle}{25}{1}{4}
  {\begin{cellatartalom}%
    %
  \end{cellatartalom}}
\ora{\mitrakjonle}{25}{1}{5}
  {\begin{cellatartalom}%
    %
  \end{cellatartalom}}
\ora{\mitrakjonle}{25}{1}{6}
  {\begin{cellatartalom}%
    %
  \end{cellatartalom}}
\ora{\mitrakjonle}{25}{1}{7}
  {\begin{cellatartalom}%
     \oraresz{Juh{\' a}sz A.}{a}{oktv}{A.1}{1}%%
  \end{cellatartalom}}
\ora{\mitrakjonle}{25}{1}{8}
  {\begin{cellatartalom}%
     \oraresz{Juh{\' a}sz A.}{a}{oktv}{A.1}{1}%%
  \end{cellatartalom}}
\ora{\mitrakjonle}{25}{1}{9}
  {\begin{cellatartalom}%
    %
  \end{cellatartalom}}
\ora{\mitrakjonle}{25}{2}{0}
  {\begin{cellatartalom}%
    %
  \end{cellatartalom}}
\ora{\mitrakjonle}{25}{2}{1}
  {\begin{cellatartalom}%
    %
  \end{cellatartalom}}
\ora{\mitrakjonle}{25}{2}{2}
  {\begin{cellatartalom}%
    %
  \end{cellatartalom}}
\ora{\mitrakjonle}{25}{2}{3}
  {\begin{cellatartalom}%
    %
  \end{cellatartalom}}
\ora{\mitrakjonle}{25}{2}{4}
  {\begin{cellatartalom}%
    %
  \end{cellatartalom}}
\ora{\mitrakjonle}{25}{2}{5}
  {\begin{cellatartalom}%
    %
  \end{cellatartalom}}
\ora{\mitrakjonle}{25}{2}{6}
  {\begin{cellatartalom}%
    %
  \end{cellatartalom}}
\ora{\mitrakjonle}{25}{2}{7}
  {\begin{cellatartalom}%
    %
  \end{cellatartalom}}
\ora{\mitrakjonle}{25}{2}{8}
  {\begin{cellatartalom}%
    %
  \end{cellatartalom}}
\ora{\mitrakjonle}{25}{2}{9}
  {\begin{cellatartalom}%
    %
  \end{cellatartalom}}
\ora{\mitrakjonle}{25}{3}{0}
  {\begin{cellatartalom}%
    %
  \end{cellatartalom}}
\ora{\mitrakjonle}{25}{3}{1}
  {\begin{cellatartalom}%
    %
  \end{cellatartalom}}
\ora{\mitrakjonle}{25}{3}{2}
  {\begin{cellatartalom}%
    %
  \end{cellatartalom}}
\ora{\mitrakjonle}{25}{3}{3}
  {\begin{cellatartalom}%
    %
  \end{cellatartalom}}
\ora{\mitrakjonle}{25}{3}{4}
  {\begin{cellatartalom}%
    %
  \end{cellatartalom}}
\ora{\mitrakjonle}{25}{3}{5}
  {\begin{cellatartalom}%
    %
  \end{cellatartalom}}
\ora{\mitrakjonle}{25}{3}{6}
  {\begin{cellatartalom}%
    %
  \end{cellatartalom}}
\ora{\mitrakjonle}{25}{3}{7}
  {\begin{cellatartalom}%
    %
  \end{cellatartalom}}
\ora{\mitrakjonle}{25}{3}{8}
  {\begin{cellatartalom}%
    %
  \end{cellatartalom}}
\ora{\mitrakjonle}{25}{3}{9}
  {\begin{cellatartalom}%
    %
  \end{cellatartalom}}
\ora{\mitrakjonle}{25}{4}{0}
  {\begin{cellatartalom}%
    %
  \end{cellatartalom}}
\ora{\mitrakjonle}{25}{4}{1}
  {\begin{cellatartalom}%
    %
  \end{cellatartalom}}
\ora{\mitrakjonle}{25}{4}{2}
  {\begin{cellatartalom}%
    %
  \end{cellatartalom}}
\ora{\mitrakjonle}{25}{4}{3}
  {\begin{cellatartalom}%
    %
  \end{cellatartalom}}
\ora{\mitrakjonle}{25}{4}{4}
  {\begin{cellatartalom}%
    %
  \end{cellatartalom}}
\ora{\mitrakjonle}{25}{4}{5}
  {\begin{cellatartalom}%
    %
  \end{cellatartalom}}
\ora{\mitrakjonle}{25}{4}{6}
  {\begin{cellatartalom}%
    %
  \end{cellatartalom}}
\ora{\mitrakjonle}{25}{4}{7}
  {\begin{cellatartalom}%
    %
  \end{cellatartalom}}
\ora{\mitrakjonle}{25}{4}{8}
  {\begin{cellatartalom}%
    %
  \end{cellatartalom}}
\ora{\mitrakjonle}{25}{4}{9}
  {\begin{cellatartalom}%
    %
  \end{cellatartalom}}
\ora{\mitrakjonle}{25}{5}{0}
  {\begin{cellatartalom}%
    %
  \end{cellatartalom}}
\ora{\mitrakjonle}{25}{5}{1}
  {\begin{cellatartalom}%
    %
  \end{cellatartalom}}
\ora{\mitrakjonle}{25}{5}{2}
  {\begin{cellatartalom}%
    %
  \end{cellatartalom}}
\ora{\mitrakjonle}{25}{5}{3}
  {\begin{cellatartalom}%
    %
  \end{cellatartalom}}
\ora{\mitrakjonle}{25}{5}{4}
  {\begin{cellatartalom}%
    %
  \end{cellatartalom}}
\ora{\mitrakjonle}{25}{5}{5}
  {\begin{cellatartalom}%
    %
  \end{cellatartalom}}
\ora{\mitrakjonle}{25}{5}{6}
  {\begin{cellatartalom}%
    %
  \end{cellatartalom}}
\ora{\mitrakjonle}{25}{5}{7}
  {\begin{cellatartalom}%
    %
  \end{cellatartalom}}
\ora{\mitrakjonle}{25}{5}{8}
  {\begin{cellatartalom}%
    %
  \end{cellatartalom}}
\ora{\mitrakjonle}{25}{5}{9}
  {\begin{cellatartalom}%
    %
  \end{cellatartalom}}

		
		% debug: 
%		\draw[->, ultra thick] (evcStart0)--(ELL1-3-4);
	\end{tikzpicture}
	}%endof \foreach \mitrakjonle
}%endof \orarend